%\documentstyle[12pt,german,mathsym,twoside,diplom]{report}
\documentclass[12pt,twoside]{report}

\UseLegacyTextSymbols
\usepackage{latexsym}
\usepackage{german,mathsym,diplom}
\usepackage{epsf}

\newtheorem{theorem}{Theorem}[chapter]
\newtheorem{satz}[theorem]{Satz}
%%\newtheorem{definition}[theorem]{Definition}
\newtheorem{korollar}[theorem]{Korollar}

%\newtheorem{beisp}[theorem]{Beispiel}
%\def\beispiel{\@ifnextchar [{\@beisp}{\begin{@beispiel}\em}}
%\def\endbeispiel{\end{@beispiel}}
%\def\@beisp[#1]{\begin{@beispiel}[#1]\em}
%\newtheorem{@beispiel}[theorem]{Beispiel}
\newtheorem{beispiel}[theorem]{Beispiel}

\newtheorem{bem}[theorem]{Bemerkung}
\newtheorem{kor}[theorem]{Korollar}
\newtheorem{lemma}[theorem]{Lemma}

\newcommand{\ggt}{\mbox{ggT}}
\newcommand{\kgv}{\mbox{kgV}}
\newcommand{\mod}[1]{\mbox{ mod $#1$}}
\newcommand{\euler}[1]{\Phi\left( #1 \right)}
\newcommand{\beweis}{{\bf Beweis: }}
\newcommand{\kong}[1]{\mbox{\shortstack{$\equiv$ \\ $#1$}}}
\newcommand{\leg}[2]{\left(\frac{#1}{#2}\right)}
\newcommand{\glatt}[2]{\Psi\left( #1,#2 \right)}
\newcommand{\unten}[1]{\lfloor #1 \rfloor} 
\newcommand{\oben}[1]{\lceil #1 \rceil} 
\newcommand{\dlog}[2]{\log_{#1} #2}
\newcommand{\koerper}[1]{\F_{#1}}
\newcommand{\ring}[1]{\Z_{#1}}
\newcommand{\EXP}{\mbox{exp}}
\newcommand{\LOG}{\mbox{log}}
\newcommand{\betrag}[1]{\mid #1 \mid}
\newcommand{\grad}{\mbox{grad}}
\newcommand{\Oklasse}[1]{\mbox{{\cal O}}(#1)}
\newcommand{\includeeps}[2]{
   \def\epsfsize##1##2{#2##1}
   \centerline{\epsffile{#1}}
}
\newcommand{\ds}{\displaystyle}
\newcommand{\einheit}[1]{\Z_{#1}^*}
\title{Public Key Kryptographie}

\author{W. Geiselmann}

%\includeonly{authen}

\begin{document}
% erste Seite

\thispagestyle{empty}
\phantom{*}
\vskip 3truecm

\centerline{\Huge\bf Public Key}
\medskip

\centerline{\Huge\bf Kryptographie}
\vskip 3truecm


\centerline{Vorlesungsskript}
\centerline{WS 96/97}
\medskip
\centerline{\bf Willi Geiselmann}
\vfill


\centerline{\Large Vorabversion}
\vfill


\centerline{Ausgearbeitet von}
\medskip
\centerline{\bf Pawel Wocjan}
\vfill


\centerline{Institut f\"ur Algorithmen und Kognitive Systeme}
\centerline{Universit\"at Karlsruhe}
\newpage


%\thispagestyle{empty}
%\phantom{*}

%\cleardoublepage


\thispagestyle{empty}
\phantom{*}
\vskip 5truecm

\centerline{{\Large Vorabversion}\footnote{Korrektur- und 
          Erg\"anzungsvorschl\"age werden mit Interesse entgegengenommen}}

\vfill

%zugefuegt fuer twoside Style, damit nicht Seite nur mit Seitennummer:
%\newpage
% Rueckseite wieder leer machen
%\thispagestyle{empty}
%\phantom{*}

%\cleardoublepage

\setcounter{page}{0}
\tableofcontents

\chapter*{Einf"uhrung, Terminologie}
\addcontentsline{toc}{chapter}{\numberline{}{Einf"uhrung, Terminologie}}

\section*{Terminologie}
\addcontentsline{toc}{section}{\numberline{}{Terminologie}}

Das kleinst m"ogliche Modell in der Kryptographie hat folgende Form:

\begin{itemize}
 \item Sender Alice, Empf"anger Bob;
 \item Nachricht oder Klartext $m$, Chiffrat $c$;
 \item Verschl"usselung $E$ (encryption), Entschl"usselung $D$ (decryption).
\end{itemize}

Alice will Bob eine Nachricht $m$ schicken. Alice will erreichen,
da\3 ihre Nachricht sicher "ubertragen wird, d.h. ein Unbefugter
darf die Nachricht nicht lesen oder ver"andern k"onnen. 
Um dies zu erreichen, mu\3 die Nachricht von
Alice verschl"usselt werden, so da\3 sie nur von Bob
entschl"usselt werden kann.

\begin{center}
  \begin{tabular}{ccc}
              &   \shortstack{verschl"usselt\\ $E$ encryption}  &               \\
    Klartext  &   $\longrightarrow$                             & Chiffrat      \\
    (message) &                                                 & (ciphertext)  \\
    $m$       &   $\longleftarrow$                              & $c$           \\
              &   \shortstack{entschl"usselt\\ $D$ decryption}  & 
  \end{tabular}
\end{center}

Mathematisch kann der Sachverhalt wie folgt beschrieben
werden. Die Funktion $E$ ist eine injektive Funktion von der
Menge der Klartexte in die Menge der Chiffrate, $D$ ist eine
surjektive Funktion in der entgegengesetzten Richtung.

\begin{eqnarray*}
  E(m) & = & c \\
  D(c) & = & m \\
  D(E(m)) & = & m \\
  E(D(c)) & = & c 
\end{eqnarray*}


\section*{Aufgaben der Kryptographie}
\addcontentsline{toc}{section}{\numberline{}{Aufgaben der Kryptographie}}

{\em Kryptographie} oder auch {\em Kryptologie} (aus dem Griechischen
kryp\'os "`versteckt"' und l\'ogos "`Wort"') ist die Wissenschaft, die
sich mit der sicheren (geheimen) Kommunikation befa\3t. Die
Vertraulichkeit ist nur eine der Aufgaben, die die Kryptographie heute
erf"ullen mu\3. "Ublicherweise werden die folgenden vier Eigenschaften
betrachtet:
\begin{itemize}
  \item{\bf Vertraulichkeit:} \\Geheimhaltung der Nachricht.
  \item{\bf Integrit"at:}  \\ Empf"anger kann verifizieren, da\3 die Nachricht nicht ver"andert wurde.
  \item{\bf Verbindlichkeit:} \\ Unabstreitbarkeit, Sender kann die Herkunft nicht ableugnen.
  \item{\bf Authentizit"at:} \\ eindeutige Herkunft, Nachricht ist wirklich vom Sender.
\end{itemize}


%% *************** Algorithmen und Schl"ussel
\section*{Algorithmen und Schl"ussel}
\addcontentsline{toc}{section}{\numberline{}{Algorithmen und Schl"ussel}}

Ein {\em kryptographischer Algorithmus} ist die Realisierung einer
mathematischen Funktion, die zur Ver- und Entschl"usselnung benutzt
wird. (Im allgemeinen gibt es zwei verschiedene Funktionen zur Ver-
und Entschl"usselung.)


\subsection*{Algorithmen ohne Schl"ussel}
\addcontentsline{toc}{subsection}{\numberline{}{Algorithmen ohne Schl"ussel}}

\begin{itemize}
  \item Sicherheit beruht auf der Geheimhaltung des Algorithmus.
  \item Der Algorithmus kann schlecht auf G"ute "uberpr"uft werden, 
        falls kein kompetentes Mitglied in der Gruppe vorhanden ist, 
        das den Algorithmus benutzt.
  \item Personen in mehreren Gruppen ben"otigen mehrere Algorithmen.
  \item Nach einem erfolgreichen Angriff oder nachdem ein Mitglied die Gruppe
        verlassen hat, mu\3 der Algorithmus vollst"andig ge"andert werden.
\end{itemize}


\subsection*{Algoritmen mit Schl"ussel}
\addcontentsline{toc}{subsection}{\numberline{}{Algorithmen mit Schl"ussel}}

Es wird ein Algorithmus verwendet, der von Kryptographen gepr"uft oder
standardisiert werden kann. Die Sicherheit wird durch die Wahl eines
geheimen Schl"ussels erreicht.  Es gibt zwei Klassen von Algorithmen
mit Schl"ussel.
\begin{enumerate}
  \item{\bf Symmetrische Algorithmen}\\
    Die Schl"ussel zur Ver- und Entschl"usselung sind einfach auseinander berechenbar, 
    meistens sind sie identisch. Sender und Empf"anger einigen sich vor dem Senden 
    von Nachrichten auf den (die) Schl"ussel.
  \begin{itemize}
    \item {\bf  Stromchiffren:} Nachrichten werden bit- bzw. byteweise verarbeitet. 
          Das Chiffrat eines Zeichens h"angt von den vorigen Zeichen bzw. von der Stelle im
          Text ab.
    \item {\bf Blockchiffren:} Nachrichten  werden blockweise (64 bit, 128 bit) verarbeitet. 
          Ein Klartextblock wird unabh"angig von der Stelle seines Auftretens auf ein
          Chiffrat abgebildet.
  \end{itemize}
 \item{\bf Asymmetrische Algorithmen} (oder Public-Key-Algorithmen) \\
  Es werden unterschiedliche Schl"ussel verwendet. Der Schl"ussel zur Entschl"usselung 
  l"a\3t sich nicht mit vertretbarem Aufwand aus dem Schl"ussel zum
  Verschl"usseln berechnen. Der Schl"ussel zum Verschl"usseln hei\3t 
  {\em "offentlicher Schl"ussel}, der andere {\em geheimer Schl"ussel}.
\end{enumerate}


%% ************** Kryptoanalyse
\section*{Kryptoanalyse}
\addcontentsline{toc}{section}{\numberline{}{Kryptoanalyse}}

Das Ziel der {\em Kryptoanalyse} ist es, aus einem Chiffrat $c$ die zugeh"orige Nachricht $m$ 
ohne Kenntnis des geheimen Schl"ussels zu gewinnen oder den geheimen
Schl"ussel zu finden. Die Grundannahme ist dabei, 
da"s der Algorithmus bekannt ist. Es gibt vier Typen von Angriffen:
\begin{enumerate}
  \item{\bf Ciphertext-Only} \\
   \begin{tabular}{rl}
     Geg:  & Einige Chiffrate, die mit demselben Schl"ussel und Algorithmus \\
           & verschl"usselt wurden. \\
     Ziel: & Klartexte und/oder Schl"ussel zu finden.
   \end{tabular}
  \item{\bf Known-plaintext} \\
   \begin{tabular}{rl}
     Geg:  & Paare von Klartext und Chiffrat. \\
     Ziel: & Finden des Schl"ussels (um sp"ater unbekannte Klartexte zu finden).\\
   \end{tabular}
  \item{\bf Chosen-plaintext} \\
   \begin{tabular}{rl}
     Geg:  & Klartexte k"onnen frei gew"ahlt werden, deren Chiffrate dann berechnet\\
           & werden. \\
     Ziel: & Schl"ussel zu finden.
   \end{tabular}
  \item{\bf Chosen-ciphertext} \\
   \begin{tabular}{rl}
     Geg:  & Chiffrate k"onnen frei gew"ahlt werden, die dann entschl"usselt werden\\ 
           & (bei Public-Keyverfahren sinnvoll). \\
     Ziel: & Schl"ussel (geheimen Schl"ussel bei PK) zu finden.
   \end{tabular}
\end{enumerate} 
Alle Kryptosysteme k"onnen durch einen Known-plaintext-Angriff gebrochen werden, 
indem man alle m"oglichen Schl"ussel ausprobiert. Diese Art, den Schl"ussel aus einem
Klartext-Chiffrat-Paar zu gewinnen, hei\3t {\em Brute-force-Angriff}. Oft ist dieser 
Angriff wegen der Redundanz in der Nachricht (Sprache, bekannter Anfang, Anrede, usw. ) auch
m"oglich, wenn nur Chiffrate gegeben sind. Die Kryptographie sucht Algorithmen, die mit den zur 
Verf"ugung stehenden Mitteln (heutige oder zuk"unftige) nicht gebrochen werden k"onnen.

\section*{Sicherheit von Algorithemen}
\addcontentsline{toc}{section}{\numberline{}{Sicherheit von Algorithemen}}

Es gibt vier Kategorien, die das Versagen eines kryptographischen Algorithmus beschreiben.
\begin{enumerate}
\item{\bf Vollst"andiger Zusammenbruch} \\ 
  Der (geheime) Schl"ussel wurde gefunden.
\item{\bf Globale Entschl"usselung} (global deduction) \\
 Es wurde ein anderer Algorithmus zum Entschl"usseln gefunden, ohne den Schl"ussel zu kennen.
\item{\bf Lokale Entschl"usselung} \\
  Es wurden von einigen Chiffraten die Klartexte gefunden.
\item{\bf Informationsgewinn} \\
  Teile oder Information "uber den Klartext und/oder den Schl"ussel wurden gefunden.
\end{enumerate}


\section*{Aufwand von Angriffen}
\addcontentsline{toc}{section}{\numberline{}{Aufwand von Angriffen}}

Der Aufwand der Angriffe kann auf drei verschiedene Weisen gemessen werden:
\begin{enumerate}
  \item{\bf Rechenkapazit"at:} \\
    Zeit bzw. Anzahl der Rechenschritte.
  \item{\bf Speicherkomplexit"at:} \\
    ben"otigter Speicher.
  \item{\bf Datenkomplexit"at:} \\
    ben"otigte Anzahl an Chiffraten bzw. Chiffrat-Klartext-Paaren.
\end{enumerate}    
  
  

\chapter{Protokolle}

\section{Einf"uhrung}

Ein Protokoll besteht aus einer Reihe von Schritten, die nacheinander
ausgef"uhrt werden. Es sind mindestens zwei Teilnehmer am Protokoll
beteiligt. Durch das Protokoll soll ein gewisser Zweck erf"ullt
werden. Weiterhin m"ussen noch folgende Bedingungen erf"ullt sein:

\begin{itemize}
  \item Das Protokoll ist jedem Teilnehmer bekannt.
  \item Jeder Teilnehmer ist mit dem Protokoll einverstanden. 
  \item Der Ablauf des Protokolls mu\3 eindeutig sein.
  \item Das Protokoll mu\3 vollst"andig sein (es ist in allen 
        F"allen genau bestimmt, was zu tun ist).
\end{itemize}

Das Ziel eines kryptographischen Protokolls ist es, Abh"oren und
Betr"ugen zu verhindern. Au\3erdem darf kein Teilnehmer mehr erfahren
k"onnen als das, was im Protokoll vorgesehen ist.

\begin{beispiel} Informale Protokolle aus dem t"aglichen Leben:
  \begin{itemize}
    \item Telefonische Bestellung von G"utern;
    \item Bezahlen mit Kreditkarte;
    \item B"urgerentscheid;
    \item Kartenspiele.
  \end{itemize}
\end{beispiel}

Diese Protokolle funktionieren meistens, da man die Teilnehmer kennt,
sieht oder h"ort. Computer brauchen im Gegensatz zu Menschen formale
Protokolle.

%% ************* Protokol mit Notar
\section{Protokoll mit Notar}

Ein {\em Notar} ist eine vertrauensw"urdige Instanz, die keinen der
Teilnehmer bevorzugt. Mit Hilfe einer solchen Instanz lassen sich recht einfach
Protokolle angeben.

\begin{beispiel}[Verkauf eines Autos mit Hilfe eines Notars] 
Der Verkauf kann wie folgt abgewickelt werden:
  \begin{enumerate}
    \item Alice will Bob ein Auto verkaufen.
    \item Bob will mit einem Scheck bezahlen.
    \item Alice gibt den KFZ-Schein dem Notar.
    \item Bob gibt den Scheck an Alice, die ihn dann bei der Bank einl"ost.
    \item Falls der Scheck gedeckt war, gibt der Notar den KFZ-Schein an Bob.
    \item Falls aber der Scheck nicht gedeckt war, mu\3 Alice dies dem Notar 
          beweisen und erh"alt daraufhin den KFZ-Schein vom Notar zur"uck.
  \end{enumerate}
Eine zweite L"osung w"are ein von der Bank zertifizierter Scheck. In
diesem Fall k"onnte man ohne Notar auskommen.
\end{beispiel}

Das obige Beispiel l"a\3t sich auf Rechnernetzwerke "ubertragen. Ein
wesentlicher Nachteil ist es, da\3 bei jeder Transaktion der Notar involviert
ist. Dieser Engpa\3 kann durch den Einsatz mehrerer Notare behoben
werden, was aber die Kosten steigert. Eine andere M"oglichkeit, diesen
Engpa\3 zu beheben, ist es, das Protokoll wie folgt in zwei Teile zu
zerlegen:
\begin{itemize}
  \item{\bf Normalfall}
    \begin{itemize}
      \item Alice und Bob handeln einen Vertrag aus,
      \item Alice unterschreibt,
      \item Bob unterschreibt.
    \end{itemize}
  \item{{\bf Ausnahmefall} (Vertrag wird nicht eingehalten)}
    \begin{itemize}
      \item Alice und Bob erscheinen vor Gericht,
      \item Alice und Bob erl"autern den Fall,
      \item Richter f"allt eine Entscheidung.
    \end{itemize}
\end{itemize}
Dieses Protokoll verhindert nicht einen Betrug, hat aber einen {\em
Abschreckungseffekt}. Das Beste, was man manchmal auch erreichen kann,
ist ein {\em abgesichertes Protokoll}. Dieses Protokoll garantiert,
da\3 keiner der Teilnehmer betr"ugen kann. Sobald sich ein Partner
nicht an das Protokoll h"alt (betr"ugen will), bemerkt das der andere
Teilnehmer und bricht ab.

%% *************** Typen von Angriffen gegen Protokolle
\section{Typen von Angriffen gegen Protokolle}

\begin{itemize}
  \item{\bf passiver Angriff:} \\
    Ein Unbeteiligter h"ort die Nachrichten ab und versucht daraus 
    Information zu gewinnen.
  \item{\bf aktiver Angriff:} \\
    Ein Unbeteiligter "andert die Nachrichten (f"angt beispielsweise 
    Nachrichten ab und schickt andere).
  \item{\bf passiver Betrug:} \\
    Ein Teilnehmer des Protokolls versucht zus"atzliche Informationen 
    zu gewinnen (bei\-spielsweise w"ahlt er bei einer vom Protokoll 
    verlangten Wahl von Zufallszahlen geschickt und nicht zuf"allig).
  \item{\bf aktiver Betrug:} \\
    Ein Teilnehmer h"alt sich nicht an das Protokoll (er gibt 
    beispielsweise einen falschen Namen an).
\end{itemize}

%% ************** Eine erste Public-Key-Idee
\section{Eine erste Public-Key-Idee (Merkles Puzzle)}

Merkles Puzzel basiert auf der Idee, R"atsel zu erzeugen, die f"ur den Sender 
Alice und
den Empf"anger Bob in vertretbarer Zeit zu l"osen sind, nicht aber
f"ur einen Abh"orer Eve. Damit soll das Ziel erreicht werden, da\3 vor
der Kommunikation kein gemeinsamer Schl"ussel zwischen Alice und Bob
ausgetauscht werden mu\3. 

{\bf Vorgehensweise:} 
\begin{enumerate}
  \item Bob erzeugt $2^{20} (\approx 10^6)$ Nachrichten der Form {\bf
    "`R"atselnr.: $x_i$, Schl"ussel: $y_i$"'}, wobei $x_i$ und
    $y_i$ jeweils paarweise verschiedene Zufallszahlen sind. Die $y_i$
    haben eine L"ange von mindestens 40 bit. Bob verschl"usselt diese 
    Nachrichten mit einem symmetrischen Verfahren, wobei er verschiedene 
    Schl"ussel der L"ange 20 bit benutzt, und schickt alles an Alice.
  \item Alice w"ahlt eine der verschl"usselten Nachrichten zuf"allig
    aus und entschl"usselt sie durch einen Brute-force-Angriff
    (Aufwand ungef"ahr $10^6$ Schritte).  Alice erh"alt nach der
    Entschl"usselung die beiden Zufallszahlen $x_0$ und $y_0$.
  \item Alice verschl"usselt die Nachricht $m$ mit dem Schl"ussel
    $y_0$ und schickt das Chiffrat zusammen mit $x_0$ im Klartext an
    Bob.
  \item Bob w"ahlt das zu $x_0$ geh"orende $y_0$ aus seiner Datenbank
    und entschl"usselt das Chiffrat von Alice mit dem Schl"ussel $y_0$.
\end{enumerate}

{\bf Aufwandsbetrachtung: } \\
\begin{tabular}{rcl}
  Bob          & $\approx$ & $2^{20}$ Schritte (Erzeugung der R"atsel) \\
  Senden       & $\approx$ & $2^{20}$ R"atsel von Bob und Chiffrat von Alice\\
  Alice        & $\approx$ & $2^{20}$ Schritte (Brute-force-Angriff) und 
                 Verschl"usselung des Klartextes\\
  Eve          & $\approx$ & $2^{40}$ Schritte (Brute-force-Angriff)
\end{tabular} 

Der Abh"orer Eve mu\3 im schlimmsten Fall $2^{20}$ R"atsel mit
Brute-force entschl"usseln bis er $x_0$ gefunden hat und damit auch
$y_0$. Dazu braucht er aber ungef"ahr $2^{40}$ Schritte, da jedes
R"atsel mit einem 20 bit langen Schl"ussel verschl"usselt wurde.
Es ist f"ur Eve ung"unstiger, das Chiffrat von Alice direkt mit einem
Brute-force-Angriff zu entschl"usseln, da er dazu mehr als
$2^{40}$ Schritte br"auchte, weil die Schl"ussel $y_i$ eine L"ange von mehr als
40 bit haben.

Wenn Alice und Bob 10000 Operationen pro Sekunde ausf"uhren k"onnen, dann
ben"otigen sie jeweils ungef"ahr eine Minute. Eve mu\3 aber bei
vergleichbarer Rechenleistung etwa ein Jahr investieren. Man sieht
also, da\3 der Aufwand f"ur den Abh"orer Eve unvergleichbar schneller
w"achst, wenn die Anzahl der R"atsel und die L"ange der Schl"ussel erh"oht
wird.

%************* Einwegfunktionen
\section{Einwegfunktionen}

Einwegfunktionen bilden einen zentralen Baustein der Public-Key-Protokolle. 

\begin{definition}[Einwegfunktion]
Eine {\em Einwegfunktion} ist eine injektive Funktion $f$, die einfach zu
berechnen ist, deren Umkehrfunktion $f^{-1}$ dagegen sehr schwierig zu
berechnen ist.
\end{definition}

Eine mathematisch exakte Definition w"are: Die Funktion $f$ ist in
polynomialer Zeit zu berechnen, die Umkehrfunktion $f^{-1}$ aber
nicht. Funktionen, die diese Eigenschaft erf"ullen, sind bisher
nicht bekannt.

\begin{definition}[Einwegfunktion mit Fallt"ur] 
Eine {\em Einwegfunktion mit Fallt"ur} ist eine Einwegfunktion $f$,
deren Umkehrfunktion $f^{-1}$ einfach zu berechnen ist, falls man eine
zus"atzliche Information besitzt.
\end{definition}

%************* Hashfunktionen  
\section{Hashfunktionen}

\begin{definition}[Hashfunktion]
Eine {\em Hashfunktion} ist eine Funktion $H$, die eine beliebig lange Eingabe 
auf eine Ausgabe fester L"ange (normalerweise k"urzer) abbildet. 
\end{definition}

\begin{beispiel} 
Das XOR aller Eingabebytes ist eine Hashfunktion.
\end{beispiel}

\begin{definition}[kryptographische Hashfunktion]
  Eine {\em kryptographische Hashfunktion} ist eine Hashfunktion, die
  zus"atzlich noch folgende Bedingungen erf"ullt:
  \begin{itemize}
    \item sie ist einfach zu berechnen,
    \item es ist schwierig zu einer vorgegebenen Ausgabe (Hashwert)
      ein Urbild (eine der h"ochstwahrscheinlich unendlich vielen
      Nachrichten mit diesem Hashwert) zu finden, und 
    \item es ist schwierig zwei Nachrichten mit demselben Hashwert 
      (Kollision) zu finden.
  \end{itemize}
\end{definition}

%% ************************ Hybrid-Kryptosysteme   
\section{Hybrid-Kryptosystem}

Public-Key-Kryptosysteme sind etwa einen Faktor 1000 langsamer als
Blockchiffren bei gleichen Sicherheitsanforderungen. Bei den symmetrischen 
Verfahren mu\3 aber vor der Kommunikation ein geheimer Schl"ussel sicher
ausgetauscht werden. Daher wird in den meisten praktisch eingesetzten
Systemen zun"achst ein geheimer Sitzungsschl"ussel $K$ mit einem
Public-Key-Verfahren "ubermittelt. Danach wird ein schnelleres
symmetrisches Verfahren verwendet, das dann den Sitzungsschl"ussel $K$
benutzt. Diese Systeme hei\3en {\em Hybridsysteme}, da zwei
kryptographische Verfahren eingesetzt werden.
\begin{enumerate}
  \item Bob schickt Alice seinen "offentlichen Schl"ussel $B$.
  \item Alice w"ahlt einen Sitzungsschl"ussel $K$, verschl"usselt ihn
         mit $B$ und schickt $E_B(K)$ an Bob.
  \item Bob entschl"usselt $D_{B'}(E_B(K))=K$ mit seinem geheimen
        Schl"ussel $B'$. Jetzt haben beide den gemeinsamen
        Sitzungsschl"ussel $K$, der nur ihnen bekannt ist.
  \item Ab jetzt l"auft die Kommunikation mit einem schnellen
        symmetrischen Verfahren, das den Sitzungsschl"ussel $K$
        verwendet.
\end{enumerate}

%% ********************* Digitale Signatur
\section{Digitale Signatur}

Eine digitale Signatur soll "ahnlich wie eine handschriftliche Unterschrift 
folgende Eigenschaften erf"ullen:
\begin{itemize}
  \item authentisch, d.h. einer Person eindeutig zugeordnet;
  \item nicht kopierbar, d.h. Teile des Dokuments oder das Dokument sind
    nach der Unterschrift nicht "anderbar;
  \item unabstreitbar;
  \item kurz im Vergleich zum Dokument.
\end{itemize}
Es ist schwierig, elektronisch alle Kritertien gleichzeitig zu realisieren.


\subsection{Signatur mit Hilfe eines Notars und eines symmetrischen Kryptosystems} 

"Ubung: Geben Sie unter Verwendung eines symmetrischen Kryptosystems
und eines Notars ein Signaturverfahren an. Wie werden die folgenden
Aufgaben gel"ost:

\begin{itemize}
   \item Wie "ubermittelt Alice an Bob eine signierte Nachricht?
   \item Wie kann Bob Celia nachweisen, da\3 das Dokument von Alice signiert wurde?
\end{itemize}
Welche Kriterien werden von Ihrem Verfahren erf"ullt?


\subsection{Signatur mit Public-Key-Verfahren}

Public-Key-Verfahren k"onnen auch zur Authentifizierung verwendet
werden, indem man den geheimen Schl"ussel zur Verschl"usselung benutzt.
\begin{enumerate}
  \item Alice verschl"usselt die Nachricht $m$ mit ihrem geheimen
        Schl"ussel $A'$. Sie f"ugt dem Chiffrat $c=D_{A'}(m)$ 
        ihren Namen hinzu und schickt dies an Bob.
  \item Bob entschl"usselt $c$ mit Alices "offentlichen Schl"ussel 
        $A$ und erh"alt die Nachricht $m=E_A(c)$.
\end{enumerate}
Diese Signatur ist kopierbar. Diese Tatsache st"ort nicht bei der 
Unterschrift von Vertr"agen. Werden damit aber Schecks unterschrieben, so 
k"onnen diese mehrfach eingel"ost werden. Es m"ussen Zeitstempel, laufende 
Nummern, etc. eingef"ugt und von der einl"osenden Bank gespeichert werden, um
dieses Problem zu beheben. Diese Unterschrift ist authentisch und
unabstreitbar. Ihre L"ange h"angt von der L"ange der Nachricht $m$ ab.
Sie ist also nicht kurz.

\subsection{Signatur mit Hashfunktionen}

Bei der Signatur mit Public-Key-Verfahren mu\3 die ganze Nachricht $m$
verschl"usselt werden. Dies ist oft mit hohem Rechenaufwand verbunden,
wenn die Nachricht $m$ gro\3 ist. Dieses Problem l"a\3t sich mit Hilfe
einer Hashfunktion beheben.
\begin{enumerate}
  \item Alice berechnet den Hashwert $h=H(m)$ der Nachricht $m$.
  \item Alice verschl"usselt $h$ mit ihrem geheimen Schl"ussel $A'$ 
        und erh"alt ihre Unterschrift $D_{A'}(h)$ der Nachricht $m$.
  \item Alice schickt ihre Nachricht $m$ mit ihrer Unterschrift $D_{A'}(h)$ 
        an Bob.
  \item Bob berechnet den Hashwert der empfangenen Nachricht und
        vergleicht ihn mit $E_{A}(D_{A'}(h))$, wobei $A$ der
        "offentliche Schl"ussel von Alice ist.
\end{enumerate}

{\bf Vorteil:} Deutlich geringerer Rechenaufwand und Trennung von
Nachricht und kurzer Unterschrift. \\
{\bf Nachteil:} Es ist nicht so sicher. Die Sicherheit h"angt von der
Wirksamkeit der Angriffe gegen Hashfunktionen ab.

\subsection{Angriffe gegen Hashfunktionen}

Es soll zu einer Nachricht $m$ eine weitere Nachricht $m'$ gefunden werden, 
so da\3 $H(m)=H(m')$ (Kollision) gilt. Dies kann verwendet werden, um dem
Empf"anger eine falsche Nachricht zuzuspielen, ohne da\3 er dies bemerken kann.
\begin{itemize}
\item{\bf Brute-force-Angriff} \\
   Die Nachricht $m$ und der Hashwert $H(m)$ seien bekannt. Es werden
   solange verschiedene Nachrichten $m'$ ausprobiert, bis $H(m)=H(m')$
   gilt. Durch Ab"andern der Nachricht $m$ an $n$ Stellen (durch
   F"ullworte, Umstellungen) ergeben sich $2^n$ verschiedene Texte
   $m'$. Hat der Hashwert $k$ bit L"ange, so m"ussen etwa $2^k$
   Nachrichten $m'$ ausprobiert werden, bis $H(m)=H(m')$ auftritt.
\item{\bf Meet-in-middle-Angriff} \\
   Man kann Hashfunktionen normalerweise derart modifizieren, da\3 sie
   eine Nachricht von hinten, also vom bekannten Hashwert $H(m)$ der
   Nachricht, r"uckhashen.  Diese Tatsache kann ausgenutzt werden,
   um eine Kollision schneller als mit dem Brute-force-Angriff zu
   finden. Dazu wird die Nachricht $m$ in zwei Teile $m_1$ und $m_2$
   geteilt. Im ersten Teil werden "Anderungen an $k/2$ Stellen
   vorgenommen. Es enstehen dadurch $2^{k/2}$ verschiedene Nachrichten
   $m'_1$ mit Vorw"artshashwerten $H(m'_1)$. Diese Nachrichten werden nach
   dem Hashwert sortiert.

   Mit Hilfe des bekannten Hashwerts $H(m)$ wird solange der  
   R"uckwerts\-hashwert $H_{r"uck}(m'_2,H(m))$ vom entsprechend ver"anderten 
   zweiten Teil $m'_2$ der Nachricht $m$ berechnet, bis eine "Ubereinstimmung
   des R"uckw"artshashwerts $H_{r"uck}(m'_2,H(m))$ mit einem der
   Vorw"artshashwerte $H(m'_1)$ auftritt. Sie kann effizient
   mit Bin"arsuche festgestellt werden.  Aufgrund des 
   Geburtstagsparadoxon sind verh"altnism"a\3ig wenige Ver"anderungen 
   an $m_2$ notwending. Eine Kollisionsnachricht 
   $m'$ ergibt sich durch Hintereinanderschreiben von $m_1'$ und $m_2'$.
\end{itemize}

{\bf Vergleich des Aufwands f"ur k=64}

\begin{itemize}
   \item Brute-force-Angriff \\
     \begin{tabular}{rcl}
       Zeitaufwand     & $\approx$ & $2^{64}$ Schritte ($2^{64}$ $m'$ hashen) \\
       Speicheraufwand & $\approx$ & nur eine Nachricht
     \end{tabular} 
   \item Meet-in-middle-Angriff \\
     \begin{tabular}{rcl}
       vorw"arts hashen      & $\approx$ & $2^{32}$ Schritte ($2^{32}$ $m'_1$ hashen)          \\
       Sortieren der Tabelle & $\approx$ & $32 \times 2^{32}$ Schritte                         \\
       Speicheraufwand       & $\approx$ & $2^{32} \times 8$ byte $=$ 32 GByte f"ur $H(m'_1)$
    \end{tabular}
\end{itemize}
{\bf L"osung des Meet-in-the-middle Problems:} Will man $k$ bit
Sicherheit, so mu\3 der Hashwert $2k$ bit lang sein.


%% *************** Schl"usselaustausch
\section{Protokolle f"ur Schl"usselaustauch} 

Es werden Protokolle ben"otigt, die den sicheren Austausch eines
gemeinsamen Sitzungs\-schl"ussels $K$ zwischen zwei Teilnehmern
koordinieren, um ihnen sp"ater den Einsatz eines schnellen
symmetrischen Verfahrens zu erm"oglichen.

\subsection{Protokoll mit einem symmetrischen Verfahren und einer Schl"usselzentrale (SZ)}\label{SZmitSymmSchl}
  Die Schl"ussel von Alice und Bob sind der SZ bekannt. 
  \begin{enumerate}
     \item Alice teilt der SZ mit, da\3 sie mit Bob kommunizieren 
           will.
     \item Die SZ erzeugt einen Sitzungsschl"ussel $K$, berechnet 
           $E_A(K)$ und $E_B(K)$ und schickt dies an Alice.
     \item Alice berechnet $D_A(E_A(K))=K$.
     \item Alice schickt $E_B(K)$ an Bob.
     \item Bob berechnet $D_B(E_B(K))=K$.
     \item Alice und Bob besitzen jetzt einen gemeinsamen 
           Sitzungschl"ussel $K$, der nur ihnen und der SZ bekannt ist.
  \end{enumerate}
{\bf Problem:} Die SZ kennt die Schl"ussel aller Teilnehmer und alle
Sitzungsschl"ussel und kann somit alle Nachricht abh"oren.

\subsection{Protokoll mit Public-Key und einer Schl"usselzentrale}\label{SZmitPK}
Die "offentlichen Schl"ussel von Alice und Bob sind in einer Datenbank
der Schl"usselzentrale gespeichert.
\begin{enumerate}
   \item Alice holt Bobs "offentlichen Schl"ussel $B$ aus der Datenbank.
   \item Alice erzeugt einen zuf"alligen Sitzungsschl"ussel $K$,
         berechnet $E_B(K)$ und schickt dies mit ihrem Namen versehen
         an Bob.
   \item Bob entschl"usselt $D_{B'}(E_B(K))=K$ mit seinem geheimen 
         Schl"ussel $B'$.
   \item Alice und Bob haben jetzt einen gemeinsamen Sitzungsschl"ussel $K$.
\end{enumerate}
{\bf Vorteil:} Die SZ kennt nur die "offentlichen der Teilnehmer. Ihr sind die
Sitzungschl"ussel nicht bekannt.

Dieses Protokoll bietet noch keine Sicherheit, da sich ein aktiver Angreifer
in einem Man-in-the-middle Angriff zwischen die SZ und die Teilnehmer
schalten kann. Er kann dann den Teilnehmer falsche Schl"ussel schicken,
ohne da"s es von ihnen bemerkt wird.

\subsection{Man-in-the-middle Angriff}
Ein passiver Angreifer Eve kann nur das Public-Key-Verfahren oder das
symmetrische Verfahren angreifen. Eve sollte keine Chance haben, diese
Verfahren zu brechen.  Ein aktiver Angreifer Mallory kann aber die
Leitung unterbrechen und die Kommunikation, die "uber sie l"auft,
kontrollieren.
\begin{enumerate}
  \item Alice schickt ihren "offentlichen Schl"ussel $A$ an
        Bob. Mallory f"angt $A$ ab und schickt seinen "offentlichen
        Schl"ussel $M$ an Bob mit Absender Alice.
  \item Bob schickt seinen "offentlichen Schl"ussel $B$ an
        Alice. Mallory f"angt $B$ ab und schickt seinen "offentlichen
        Schl"ussel $M$ an Alice mit Absender Bob.
  \item Alice verschl"usselt die Nachricht $m$ und schickt das Chiffrat
        $c=(E_M(m))$ an Bob.
  \item Mallory entschl"usselt $D_{M'}(c)=m$, liest die Nachricht
        und leitet $E_B(m)$ weiter an Bob.
  \item Der R"uckweg verl"auft analog.
\end{enumerate}
Die von Alice und Bob ausgetauschten Nachrichten k"onnen von Mallory
nicht nur gelesen, sondern auch in beliebiger Weise manipuliert
werden. Alice und Bob haben keine M"oglichkeit festzustellen, ob
ihre Nachrichten abgeh"ort oder manipuliert wurden.

\subsection{Interlock-Protokoll}
Die besten Public-Key-Verfahren sind also wertlos, wenn die Teilnehmer
falsche Schl"ussel von einem Angreifer Mallory bei einem
Man-in-the-middle-Angriff bekommen. Dieser Schwachpunkt kann
teilweise durch das folgende Protokoll behoben werden.
\begin{enumerate}
  \item Alice schickt Bob ihren "offentlichen Schl"ussel $A$.
  \item Bob schickt Alice seinen "offentlichen Schl"ussel $B$.
  \item Alice berechnet $E_B(m_A)=e_1 e_2$ und schickt die erste
        H"alfte $e_1$ an Bob.
  \item Bob berechnet $E_A(m_B)=f_1 f_2$ und schickt die erste H"alfte
        $f_1$ an Alice.
  \item Alice schickt die zweite H"alfte $e_2$ an Bob.
  \item Bob berechnet $D_{B'}(e_1 e_2)=m_A$. Bob beantwortet {\em sofort}
        die Frage von Alice und schickt die zweite H"alfte $f_2$.
  \item Alice beantwortet Bobs Frage $D_{A'}(f_1 f_2)$.
\end{enumerate}
Ein aktiver Angreifer Mallory kann bis zum 3-ten Schritt alles wie
bisher machen. Dann mu\3 aber Mallory eine neue Nachricht erfinden,
sie verschl"usseln und die erste H"alfte an Bob schicken. Im 4-ten
Schritt mu\3 Mallory nochmals eine neue Nachricht erfinden, sie
verschl"usseln und die erste H"alfte an Alice schicken.  Im 5-ten
Schritt wei\3 Mallory Alices Frage, mu\3 aber die zweite H"alfte der
erfundenen Nachricht an Bob schicken. Im 6-ten Schritt erwartet Alice
eine Antwort von Bob, die jetzt aber Mallory beantworten
mu\3. Mallory wei\3 entweder die Antwort oder mu\3 raten. Der 7-ten
Schritt verl"auft analog. Durch das Interlock-Protokoll wird der
Man-in-the-midde Angriff Mallory erschwert, jedoch nicht unm"oglich
gemacht.

Das Protokoll versagt, wenn  Mallory die beiden Teilnehmer kennt oder wenn
Alice und Bob sich nicht direkt kennen und daher keine geeigneten Fragen
stellen k"onnen, d.h. kein gemeinsames Geheimnis besitzen. Das ist vor
allem dann der Fall, wenn Alice und Bob keine Menschen, sondern
Computer sind, die Fragen nur automatisch erzeugen k"onnen.

%% ******************* Schl"usselaustausch mit Signaturen
\subsection{Schl"usselaustausch mit Signaturen}
Das Interlockprotokoll erschwert einen Man-in-the-middle Angriff,
macht ihn aber nicht unm"oglich. Das Protokoll~\ref{SZmitPK} 
kann mit Hilfe der digitalen Signatur so erweitert werden, da"s
ein Man-in-the-middle Angriff f"ur einen Au"senstehenden Mallory
umm"oglich wird.

Die Schl"usselzentrale unterschreibt die "offentlichen Schl"ussel
von Alice und Bob, wobei sie pers"onlich anwesend sein m"ussen. Diese
von der Schl"usselzentrale signierten Schl"ussel sind in einer
Datenbank zug"anglich. Jeder Teilnehmer kennt den "offentlichen
Schl"ussel der SZ und kann damit ihre Unterschrift pr"ufen. Mallory
kann keinen Man-in-the-middle Angriff durchf"uhren, da Alice
und Bob die Authentizit"at der Schl"ussel "uberpr"ufen k"onnen. 
Mallory kann jetzt nur noch zuh"oren oder die Leitung
st"oren. Die SZ kann nicht wie bei der Vergabe von
Sitzungsschl"usseln im Protokoll~\ref{SZmitSymmSchl} alle Nachrichten
abh"oren, ist also nicht mehr der gro\3e Schwachpunkt im System.

Die Schl"usselzentrale kann jetzt nur noch als einzige Instanz einen
Man-in-the-middle Angriff durchf"uhren.

Mit Hilfe einer SZ lassen sich auch sehr einfach Nachrichten
verschicken. Der "offentliche Schl"ussel von Bob wird von der SZ
geliefert. Alice erzeugt Sitzungsschl"ussel $K$, verschl"usselt
Nachricht $m$ mit symmetrischen Verfahren und schickt
$(E_B(K),E_K(m))$ an Bob. Hier ist wieder wichtig, da\3 der "offentliche
Schl"ussel $B$ signiert ist.

\section{Zero Knowledge Protokolle}

Bei einem Zero-knowledge Protokoll weist ein Beweiser $B$ einem
Verifizierer $V$ nach, die L"osung eines bestimmten Problems $P$ zu
kennen. Bei diesem Nachweis wird keine Information "uber die Struktur
der L"osung preisgegeben. Genauer ist gefordert, da\3 bei polynomial
vielen Abl"aufen des Protokolls nicht mehr Information preisgegeben
wird, als die, die der Verifizierer selbst mit polynomialen Aufwand h"atte
bekommen k"onnen.

Die Zero Knowledge Protokolle k"onnen f"ur Authetifikationsverfahren
benutzt werden.

\begin{beispiel} 
Der Keller hat eine Geheimt"ur. $B$ will $V$ beweisen, da\3 er die
Geheimt"ur "offnen kann, ohne aber $V$ das Geheiminis, wie sie
ge"offnet wird, zu verraten.
  \begin{figure}[htb]
  \centerline{\includeeps{keller.eps}{1.5}}
  \caption{Keller mit Geheimt"ur}
  \end{figure}
  \begin{enumerate}
    \item $V$ steht bei $e$.
    \item $B$ geht entweder zu $g$ oder $h$ im Keller.
    \item $V$ geht zu $f$ und ruft $B$ zu, er solle links bzw. 
          rechts herauskommen.
    \item $B$ "offnet bei Bedarf die Geheimt"ur und kommt am 
          gew"unschten Ort heraus.
    \item Diese Prozedur wird n-mal wiederholt. Danach ist die
          Wahrscheinlichkeit, da\3 $B$ in Wirklichkeit die Geheimt"ur
          gar nicht "offnen kann und jedesmal nur Gl"uck hatte, nur
          $1/2^n$.
  \end{enumerate}
\end{beispiel}

\begin{definition}[Perfekt zero-knowledge]
Ein Protokoll hei"st {\em perfekt zero-knowledge}, wenn es einen
Simulator gibt, der mit polynomialen Aufwand das Protokoll so
simuliert, da\3 es identisch zur richtigen Protokollausf"uhrung
ist. Der Simulator kennt dabei das Geheiminis nicht.
\end{definition}

Die abstrakte Definition wird weiter unten eingehender erl"autert. Es
l"a\3t sich auch eine etwas schw"achere Definition erstellen, die die
absolute Sicherheit durch {\em nicht berechenbar mit heutigen Mitteln}
ersetzt:

\begin{definition}[Berechenbar zero-knowledge]
Ein Protokoll ist {\em berechenbar zero-knowledge}, wenn es einen
Simulator gibt, der mit polynomialen Aufwand das Protokoll so
simuliert, da\3 es mit heutigen Mitteln nicht von einer richtigen
Protokollausf"uhrug zu unterscheiden ist.
\end{definition}

Berechenbar zero-knowledge bietet zwar nicht den unbedingten Schutz von
perfekt zero-knowledge, da durch neue Algorithmen oder schnellere Computer die
Unterscheidung doch m"oglich werden kann. Es ist in vielen F"allen
das einzige, was "uber ein Protokoll bewiesen werden kann.

Der Simulator ist das zentrale Instrument, mit dem die Zero-knowledge
Eigenschaft definiert und bewiesen wird. Die Aufgabe dieses Simulators
ist es, in Interaktionen mit einem echten Verifizierer die Rolle
des Beweisers zu "ubernehemen. Dabei steht dem Simulator das Geheimnis
des Beweisers nicht zur Verf"ugung. Gefordert ist nun, da\3 er es in
einem nicht vernachl"a\3igbaren (polynomial gro\3en) Teil der
Protokollausf"uhrungen schafft, die erwarteten Antworten zu geben und
das Protokoll damit korrekt ablaufen zu lassen. Der Simulator stellt
einerseits keinen erfolgreichen Angriff auf das Protokoll dar, da er
bei konkreten Implementierungen nur in einem Teil der F"alle die
richtigen Antworten geben kann. Andererseits sind die Antworten, in
denen der Simulator korrekte Antworten geben konnte, nicht von den
Antworten des Beweisers unterscheidbar. Daher kann sich ein
Verifizierer mit Hilfe des Simulators beliebig (polynomial) viele
Ausf"uhrungen des Protokolls beschaffen, ohne mit dem Beweiser
interagieren zu m"ussen. Der Verifizierer erh"alt daher durch den
Ablauf des Protokolls mit dem Beweiser keine Informationen, die er
nicht auch mit dem Simulator h"atte erhalten k"onnen. Deswegen kann
der Beweiser sicher sein, da\3 er keine Information "uber sein
Geheimnis preisgibt, wenn er das Protokoll polynomial oft ausf"uhrt.

\begin{beispiel}
Der Keller mit Geheimt"ur ist perfekt zero-knowledge. Was w"urde
passieren, wenn der Eingang keinen Knick machen w"urde?
\end{beispiel}

Ein Beispiel f"ur einen Zero-knowledge Beweis wird am folgenden
Protokoll, mit dem ein Beweiser nachweist, die L"osung des
Rubik W"urfels aus einer bestimmen Stellung zu kennen, gezeigt.

\begin{beispiel}[Rubik W"urfel]
$B$ will $V$ beweisen, da\3 er den W"urfel aus einer bestimmten
Stellung $P$ l"osen kann. Es wird davon ausgegangen, da\3 es 
schwierig ist, eine allgemeine L"osung des W"urfel zu finden.
\begin{enumerate}
  \item B transformiert den W"urfel in eine andere Stellung $P'$ mit 
        Hilfe einer Transformation $\tau$.
  \item V entscheidet sich, ob er die L"osung aus der Stellung $P'$ 
        heraus oder die Transformation $\tau : P \leftrightarrow P'$ 
        sehen will.
  \item B f"uhrt das Verlangte vor.
\end{enumerate}
Ist der Beweiser echt, so kann er auf jede der beiden m"oglichen
Fragen des Verifiziers eine korrekte Antwort geben. Ist er jedoch ein
Betr"uger, so kann er auf die Aufgabe "`L"ose den W"urfel aus der
Stellung $P'$"' nicht l"osen, da er daf"ur eine allgemeine L"osung des
Rubik W"urfels kennen m"u\3te, was aber nach der Grundannahme des
Protokolls schwierig ist. Die Aufgabe "`Zeige die Transformation 
$\tau : P \leftrightarrow P'$"' kann er trivialerweise l"osen.

Um zu ermitteln, ob das Protokoll zero-Knowledge ist, mu\3 ein
Simulator angegeben werden. 

Der Simulator hat zwei Rubik W"urfel, die sich im Anfangszustand bzw. im 
Zustand $P$ befinden. Er r"at zuerst die Herausforderung des
Verifiziers. Nun kann er das Problem $P'$ so konstruieren, da\3 er die
Antwort kennt, falls sich der Verifizierer f"ur die Herausforderung
entscheidet, die er vorher geraten hat. Nachdem der Simulator nun das
Problem~$P'$ pr"asentiert hat, fordert der Verifizierer den Simulator
eine der L"osungen anzugegen. Stimmt die Herausforderung mit der
geratenen "uberein, so gibt der Simulator die korrekte
Antwort. Stimmen die Herausforderungen nicht "uberein, so gibt der
Simulator auf, und fordert zu einem neuen Protokolldurchlauf auf.

Dieses Protokoll ist zwar perfekt zero-knowledge, {\em aber} der 
Verifizier kann aus dem Weg die L"osung erraten.
\end{beispiel}

Aus diesem Beispiel l"a\3t sich die typische Strukur eines
kryptographisch nutzbaren Zero-knowledge Protokolls ableiten.

{\bf Typische Struktur eines zero-knowledge Protokolls:}

Der Beweiser $B$ will dem Verifizierer $V$ beweisen, die L"osung eines
Problems $P$ aus einer Problemklasse ${\cal P}$ zu kennen:
\begin{enumerate}
  \item B bestimmt ein zu $P$ "aquivalentes Problen $P'\in {\cal P}$
        mit Hilfe einer Transformation $\tau : P \mapsto P'$ und
        sendet $P'$ als Referenzwert an V.
  \item V kann sich f"ur eine Herausforderung (challenge) $r$
        entscheiden. Im Regelfall ist dies die Frage, ob er die
        L"osung des Problems $P'$ oder die Transformation $\tau : P
        \mapsto P'$ sehen will.
  \item B sendet die zum Referenzwert $P'$ und der Herausforderung $r$
        passende Antwort.
\end{enumerate}

Kennt der Beweiser die Herausforderung $r$ bevor er seinen
Referenzwert $P'$ offenlegen mu\3, so braucht er die L"osung gar nicht
zu kennen, da er sich zuerst die Antwort w"ahlen kann und dann einen zur
Herausforderung passende Referenzwert w"ahlen kann. Daraus ergibt sich

{\bf die Strukur eines Simulators:}
\begin{enumerate}
  \item W"ahle die Herausforderung $r'$ zuf"allig und gleichverteilt.
  \item Berechne einen zu $r'$ passenden Referenzwert und eine 
        dazugeh"orige Antwort.
  \item Der Verifizierer stellt die Herausforderung $r$.
  \item Falls $r \neq r'$ gehe zu Schritt 1.
  \item Gebe den Protokollablauf mit Referenzwert, Herausforderung und 
        Antwort aus.
  \item Weiter mit Schritt 1.
\end{enumerate}
Eine M"oglichkeit f"ur die Realisierung von Zero-Knowledge Protokollen
ist relativ offensichtlich. Diese M"oglichkeit basiert auf einem
Beweiser, der behauptet eine L"osung einer Klasse von Problemen zu
besitzen. Der Beweis kann dann einfach angetreten werden, indem der
Verifizierer ein Problem aus dieser Klasse benennt und der Beweiser
eine L"osung f"ur dieses Problem vorlegt.
\begin{beispiel}[Faktorisieren] 
Der Beweiser behauptet, effizient gro\3e Zahlen faktorisieren zu k"onnen.
  \begin{enumerate}
    \item V gibt eine gro\3e Zahl an.
    \item B faktorisiert sie und V "uberpr"uft das Ergebnis.
  \end{enumerate}
Der Simulator S w"ahlt sich einige Primzahlen und berechnet das Produkt
$r'$. Falls der Verifizierer V eine Zahl $r$ w"ahlt mit $r=r'$, so
kann S ihre Faktorisierung angeben. Das Problem ist es, da"s
die Wahrscheinlichkeit f"ur $r=r'$ unendlich klein ist, da V jede
der unendlich vielen Zahlen w"ahlen kann. Gefordert ist aber, da"s der
Simulator in einem nicht vernachl"a"sigbaren Teil der 
Protokollausf"uhrungen die erwarteten Antworten geben kann.
\end{beispiel}

\begin{beispiel}[Kenntnis von Quadratwurzeln]
Es sei $n$ ein \"offentlich bekannter, ge\-n\"u\-gend gro"ser Modulus. 
Die Faktorisierung
sei nicht "offentlich bekannt. Das Berechnen von Wurzeln modulo $n$
ist "aquivalent zu der Faktorisierung von $n$ und damit schwierig. Es
sei $y_B$ ein "offentlich bekannter Wert. Der Beweiser $B$ behauptet, ein
$x_B$ mit $x_B^2 \equiv y_B \bmod n$ zu kennen.  Das folgende
Protokoll kann von ihnen benutzt werden, um diese Behauptung zu
"uberpr"ufen.
\begin{enumerate}
   \item B w"ahlt x zuf"allig und legt $y \equiv x^2 \bmod n$ offen.
   \item V w"ahlt sich $r \in \{ 0,1 \}$ und teilt es $B$ mit.
   \item B berechnet  $z \equiv x_B^r x \bmod n$.
   \item V "uberpr"uft ob $z^2 \equiv y_B^r y \bmod n$.
\end{enumerate}
Durch die Multipikation mit $x$ wird das urspr"ungliche Problem auf
ein neues Problem, die Quadratwurzel aus $y_B y$ zu ziehen,
transformiert. Falls V diese Quadratwurzel sehen will, kann B sie
berechnen. Falls aber V die Transformation sehen will, so liefert V
einfach die Zahl x. Die Transformation ist hier auch durch
Wurzelziehen zu berechen und deshalb genauso schwierig wie das
urspr"ungliche Problem.

Um einen Simulator zu konstruieren, halten wir uns an die oben 
beschriebene typische Struktur eines Simulators. Der Simulator w"ahlt 
zuf"allig und gleichverteilt die Herausforderung $r'\in\left\{0,1\right\}$. 
Falls der Verifizierer die andere Herausforderung $r\neq r'$ w"ahlt, gibt der 
Simulator auf. Sonst kann er die richtigen Antworten geben:

Simulator S: $r'=0$
\begin{enumerate}
  \item S w"ahlt $x$ zuf"allig und legt $y\equiv x^2$ offen.
  \item V entscheidet sich f"ur $r=0$.
  \item S berechnet $z:=x_B^0 x \equiv x$.
  \item V "uberpr"uft, ob $z^2\equiv y_B^0 y$.
\end{enumerate}
Simulator S: $r'=1$
\begin{enumerate}
  \item S w"ahlt $x$ zuf"allig und legt $y\equiv y_B^{-1} x^2$ offen.
  \item V entscheidet sich f"ur $r=1$.
  \item S berechnet $z:=x$.
  \item V "uberpr"uft, ob $z^2\equiv y_B^1 y$.
\end{enumerate}

In der H"alfte der F"alle liefert der Simulator richtige Antworten, d.h.
das Protokoll ist ein Zero-knowledge-Protokoll. Es ist sogar perfekt
zero knowledge, da der Simulator nur quadratische Reste ver"offentlicht.
\end{beispiel}
\begin{beispiel}[Graphenisomorphie]
Der Beweiser kennt einen Isomorphismus $\phi$ zwischen den Graphen
$G_1$ und $G_2$ und will dies dem Verifizierer beweisen. Dazu eignet
sich folgendes Protokoll:
\begin{enumerate}
  \item B transformiert den Graphen $G_1$ durch eine zuf"allige 
        Permutation $\tau$ und erh"alt einen neuen Graphen $H=\tau(G_1)$.
  \item V w"ahlt, ob er $G_1 \mapsto H$ oder $G_2 \mapsto H$ sehen will.
  \item B gibt den geforderten Isomorphismus an, da er die beiden Isomorphismen
        \begin{eqnarray*}
          \tau                  & : & G_1 \longrightarrow H \\
          \tau \circ \phi^{-1}  & : & G_2 \longrightarrow H 
        \end{eqnarray*}
        kennt.
\end{enumerate}
Das Protokoll ist perfekt zero-knowledge.
\end{beispiel}




   
  


\chapter{Public Key Algorithmen}
  Das Konzept der Public-Key-Kryptographie basiert auf der 
  Tatsache, da\3 es m"oglich ist zwei Schl"ussel 
  (einen "offentlichen $E$ zum Verschl"usseln und einen geheimen
  $D$ zum Entschl"usseln) zu benutzen, wobei es aber praktisch 
  unm"oglich ist, den geheimen Schl"ussel aus dem "offentlichen 
  zu berechnen. Viele der vorgeschlagenen Public-Key-Algorithmen 
  stellten sich als unsicher heraus und viele der sicheren Algorithmen sind nicht praktisch 
  einsetzbar, da sie entweder zu gro\3e Schl"ussel ben"otigen oder das Chiffrat
  viel zu gro\3 wird im Vergleich zu der Nachricht. Von den sicheren und praktisch 
  einsetzbaren Algorithmen eignen sich nur einige f"ur eine Schl"usselvergabe. 
  Einige der Algorithmen sind nur zur Verschl"usselung und andere nur zur digitalen 
  Signatur geeignet. Es gibt nur zwei Klassen von Algorithmen, die sich sowohl f"ur 
  Verschl"usselung als auch f"ur digitale Signatur eignen und auch "uber l"angere
  Zeit h"aufig eingesetzt werden. 

  Die eine Klasse basiert auf dem Faktorisierungsproblem; 
  das RSA-Verfahren ist der bekannteste Vertreter. Die andere Klasse basiert 
  auf dem diskreten Logarithmusproblem in Gruppen; hier ist das Verfahren von 
  ElGamal das bekannteste Verfahren. Diese Public-Key-Algorithmen sind 
  langsam im Vergleich zu symmetrischen Verschl"usselungsalgorithmen. 
  Daher werden in heutigen Systemen sowohl Public-Key-Komponenten, als auch 
  symmetrische Komponenten eingesetzt; man spricht von Hybridkryptosystemen.
%% ************** Knapsack
  \section{Knapsack}
     Die Sicherheit der Knapsack-Algorithmen beruht auf der Tatsache, 
     da\3 das allgemeine Knapsackproblem ein NP-vollst"andiges Problem ist. 
     Obwohl sich der urspr"ungliche Algorihtmus sp"ater als unsicher 
     heraustellte, lohnt es sich denoch, ihn zu untersuchen, um zu sehen, 
     wie ein NP-vollst"andiges Problem f"ur die Entwicklung eines 
     Public-Key-Algorithmus verwendet werden kann.
     \begin{definition}[Das Knapsackproblem] Es seien beliebige Gewichte 
       $a_1, \ldots ,a_n\in \N$ und eine beliebige Summe $s \in \N$ gegeben. 
       Es wird eine Teilmenge von $\{ a_1, \ldots ,a_n \}$ gesucht, so da\3 
       sich die in ihr enthaltenen Elemente zu $s$ addieren. Anders formuliert 
       lautet das Problen: Finde ein 
       $(x_1,\ldots ,x_n) \in \{0,1\}^n$ mit $\sum_{i=1}^{n}x_i a_i = s$.
     \end{definition}
     Obwohl das allgemeine Knapsackproblem NP-vollst"andig ist, gibt es 
     spezielle Knaps"acke, die einfach zu l"osen sind. Ein superincreasing 
     Knapsack ist ein Beispiel daf"ur.
     \begin{definition}[Der superincreasing Knapsack]
        Es seien $b_1,\ldots, b_n\in \N$ und $s\in \N$ ein Knapsack, 
        wobei f"ur alle $i$ gilt $\sum_{j=1}^{i-1}b_j < b_i$.
     \end{definition}
     Wegen dieser speziellen Struktur l"a\3t sich der Knapsack 
     $(b_1,\ldots ,b_n)$ mit einem Greedy-Algorithmus von oben her l"osen.

     Das n"achste Beispiel zeigt einen sehr einfachen superincreasing Knapsack.
     \begin{beispiel}[Ein einfaches Knapsackproblem]
        Es seien $a_i = 2^{i-1}$ f"ur $i=1,\ldots , n$. Damit ist es sehr einfach,
        das Knapsackproblem zu l"osen. Der L"osungsvektor $(x_1,\ldots ,x_n)$ 
        entspricht der r"uckw"arts gelesenen Bin"ardarstellung von $s$.
     \end{beispiel}
     Das Knapsackproblem kann wie folgt in einem Public-Key-Algorithmus verwendet werden:
     \begin{itemize}
       \item "offentlicher Schl"ussel: $(a_1,\ldots,a_n) \in \N^n$;
       \item geheimer Schl"ussel: Ein Algorithmus, der den Knapsack 
             $(a_1,\ldots,a_n)$ in einen einfach l"osbaren Knapsack transformiert;
       \item Verschl"usselung: $(x_1, \ldots ,x_n ) \mapsto s=\sum_{i=1}^{n} x_i a_i$;
       \item Entschl"usselung: L"ose das Knapsackproblem $(a_1,\ldots,a_n)$ f"ur die Summe $s$.
     \end{itemize}
     Merkle und Hellman stellten 1976 ein Verfahren vor, das die oben 
     beschriebene Methode benutzt. Sie verwendeten einen superincreasing 
     Knapsack und verschleierten ihn mit einer oder mehreren modularen 
     Multiplikationen.
     \subsection*{Verfahren:}
     \begin{itemize}
        \item W"ahle einen superincreasing Knapsack $(b_1,\ldots ,b_n)$ und 
          verschleiere ihn durch eine lineare Trasformation in einen anderen 
          Knapsack. W"ahle dazu $W,M\in \N$ mit $M > \sum_{i=1}^n b_i$ und 
          $\ggt(M,W)=1$. Berechne 
          \begin{equation} a_i := W b_i \pmod M \label{verschleiern} \end{equation}
          und ver"offentliche das Tupel $(a_1,\ldots,a_n)$ als "offentlichen 
          Schl"ussel. Eine bin"are Nachricht $(x_1,\ldots ,x_n)$ wird verschl"usselt zu 
          $s:=\sum_{i=1}^n x_i a_i$.
        \item Entschl"usselt wird durch Zur"uckf"uhrung auf das urspr"ungliche 
          Problem. Berechne $s' := W^{-1} s \pmod M$ 
          ($W^{-1}$ existiert, da $\ggt(W,M)=1$) und l"ose das 
          Knapsackproblem $(b_1,\ldots,b_n)$ f"ur die Summe $s'$.
     \end{itemize}   
     \subsection*{Angriff auf das Knapsackverfahren:}
     Das Knapsackverfahren weist folgende Schwachstellen auf:
     \begin{itemize}
        \item Es ist linear, d.h. es gilt 
          $\sum_{i=1}^n (x_i + y_i)a_i = \sum_{i=1}^n x_i a_i + \sum_{i=1}^n y_i a_i$ und 
        \item Wenn einige $a_i$ ungerade sind, so ist ein Bit Information 
          "uber den Klartext direkt zu sehen (s gerade oder ungerade).
     \end{itemize}
     Ein echter Angriff basiert auf der Beobachtung, da"s aufgrund der 
     Formel~\ref{verschleiern} Zahlen $k_1, \ldots ,k_n \in \Z$ mit
     \begin{equation} 
       a_i U - k_i M = b_i \label{beobachtung} 
     \end{equation}
     existieren, wobei $U \equiv W^{-1} \pmod M$. Dividiert man beide Seiten der 
     Formel~\ref{beobachtung} durch $a_i M$, so erh"alt man
     \begin{equation} 
       \frac{U}{M} - \frac{k_i}{a_i} = \frac{b_i}{a_i M}. \label{formel} 
     \end{equation}
     Weiterhin gilt $b_i \le 2^{i-n} M$, da die Zweierpotenzen den 
     kleinsten superincreasing Knapsack bilden. Substrahiert man 
     Gleichung~\ref{formel} f"ur $i:=1$ von der
     allgemeinen Gleichung~\ref{formel}, so ergibt sich
     $$ -\frac{k_i}{a_i}+\frac{k_1}{a_1}  =  \frac{b_i}{a_i M} - \frac{b_1}{a_1 M} $$
     Multipliziert man diese Gleichung mit $a_1 a_i$, so erh"alt man
     \begin{equation} 
       -a_1 k_i + a_i k_1 = \frac{a_1 b_i}{M} - \frac{a_i b_1}{M}. 
     \end{equation}
     Es gilt 
     \begin{eqnarray}
       |a_i k_1 - a_1 k_i| & \le & \max \left\{ \frac{a_1 b_i}{M}, \frac{a_i b_1}{M}\right\} \\
                           & \le & 2^{i-n} M.
     \end{eqnarray}
      F"ur sehr kleine $i$ lassen sich damit Ungleichungen in $\Z$ angeben. 
      Dabei ist $a_i$ sehr gro\3, $2^{i-n}$ sehr klein. Es reichen nur wenige 
      Ungleichungen, um die $k_i$ eindeutig festzulegen. Mit dem Algorithmus 
      von Lenstra zum linearen Programmieren lassen sich solche Ungleichungen 
      in polynomialer Zeit l"osen. Im Algorithmus "`Faktorisieren von 
      ganzzahligen Polynomen"' von Lenstra ist ein Teilalgorithmus 
      "`Finden von fast orthogonalen Basen in ganzzahligen Gittern"' enthalten. 
      Dieser Algorithmus wurde 1983 implementiert und konnte auch zum Brechen 
      des Knapsackverfahrens ausgenutzt werden. Eine zweite g"unstige Eigenschaft, 
      die das Verfahren vollst"andig zusammen brechen l"a\3t,ist es, da\3 nicht 
      die urspr"unglichen $U$, $M$ gefunden werden m"ussen. Wenn ein 
      $\frac{u}{m} \approx \frac{U}{M}$ verwendet wird, so entsteht ein 
      "aquivalenter Knapsack, der dieselbe L"osung $(x_1,\ldots, x_n)$ besitzt.
%% ************** RSA (Rivest, Shamir, Adleman)
  \section{RSA-Verfahren}
     Das RSA-Verfahren (benannt nach den Entdeckern Rivest, Shamir, Adleman) ist
     das bekannteste Public-Key-Verfahren. Die Sicherheit von RSA beruht auf der 
     enormen Schwierigkeit, gro"se Zahlen zu faktorisieren. 
     \subsection*{Verfahren:}
     \begin{enumerate}
       \item W"ahle zwei gro"se Primzahlen $p$ und $q$ (70-200 Dezimalstellen oder mehr).
       \item Berechne das Produkt $n = p\cdot q$.
       \item W"ahle einen "offentlichen Schl"ussel $e$ mit $\ggt(e,(p-1)(q-1))=1$.
       \item Berechne den geheimen Schl"ussel $d$ mit 
         $e\cdot d \equiv 1 \pmod{(p-1)(q-1)}$ und 
         ver"offentliche das Tupel $(n,e)$ als 
         "offentlichen Schl"ussel.
     \end{enumerate}
     Mit Hilfe des Chinesische Restsatzes sieht man, da"s 
     $m^{ed} \equiv m \mod{n}$ f"ur alle $m\in\Z_n$ gilt.
     \begin{itemize}
       \item Verschl"usseln: \\
         Es sind der Modulus $n$ und der Schl"ussel $e$ "offentlich bekannt. 
         Zur Verschl"usselung wird eine Nachricht $m$ in $m_1,\ldots ,m_l$ 
         zerlegt mit $m_i < n$, und es werden die 
         Chiffrate $c_i \equiv m_i^e \pmod{n}$ berechnet.
       \item Entschl"usseln: \\
         Berechne $m_i \equiv c_i^d \equiv (m_i^e)^d \equiv m_i^{ed} \equiv m_i^1$ 
         f"ur $i=1,\ldots,l$ und setze daraus die Nachricht $m$ zusammen.
     \end{itemize}
     Das RSA-Verfahren kann auch zum Signieren verwendet werden, 
     indem man mit dem geheimen Schl"ussel $d$ verschl"usselt.
     Die Signatur kann durch Entschl"usseln mit dem "offentlichen
     Schl"ussel "uberpr"uft werden.

     Das RSA-Verfahren ist in Hardware etwa um einen Faktor 1000 und in 
     Software etwa 100 mal langsamer als symmetrische Blockchiffren. 
     Zur Beschleunigung w"ahlt man 
     geeignete einfache Exponenten (es wurden zum Beispiel $e=3, 17$ und 
     $2^{16}+1$ in verschiedenen Standards vorgeschlagen). 
     Es sind durch diese Einschr"ankung bisher keine gr"o"seren Probleme 
     mit der Sicherheit bekannt.
     \subsection*{Primzahlerzeugung:}
     Die Sicherheit des RSA-Verfahrens h"angt in ganz entscheidender Weise 
     davon ab, da\3 die verwendeten Primzahlen geheim bleiben. 
     Deshalb m"ussen beim Erzeugen der Primzahlen nicht vorhersehbare 
     Zahlen generiert werden; sie d"urfen beispielsweise nicht mit einem 
     Zufallsgenerator mit einem anderen Personen bekanntem Anfangszustand
     erzeugt werden.

     Alle Verfahren, gro\3e Primzahlen zu erzeugen, w"ahlen sich zun"achst eine gro\3e 
     Zufallszahl und testen diese auf Primalit"at. Bei den Primzahltests unterscheidet 
     man zwischen zwei Klassen von Algorithmen:
     \begin{itemize}
       \item Probabilistische Primzahltests (relativ schnell und einfach),
       \item Deterministische Primzahltests (recht langsam und aufwendig).
     \end{itemize}
     Die wichtigsten Vertreter der probabilistischen Primzahltests werden
     im n"achten Kapitel vorgestellt.

     \subsubsection*{Angriff auf Protokollebene:}
     Eve belauscht die von Bob an Alice geschickte Nachricht und kennt 
     dadurch das Chiffrat $c \equiv m^e \pmod n$, wobei $e$ der 
     "offentliche und $d$ der geheime Schl"ussel von Alice sind. 
     Eve w"ahlt sich ein $r$ zuf"allig mit $\ggt(r,n)=1$ und 
     berechnet $x \equiv r^e \mod n$ und $y \equiv x c \mod n$. 
     Eve schafft es weiterhin irgendwie, da\3 Alice die Nachricht 
     $y$ signiert (beispielsweise als Zeitstempelsignatur). Damit 
     kennt Eve $y^d$. Sie berechnet 
     $$r^{-1}y^d \equiv r^{-1}x^d c^d \equiv r^{-1}(r^e)^d c^d \equiv r^{-1} r m \equiv m \pmod{n}$$
     und erh"alt damit die Nachricht $m$. Alice f"uhrt eine blinde Signatur aus.

     \subsubsection*{Angriff aufgrund eines gemeinsamen Modulus:}
       Jeder Teilnehmer erh"alt von der Zentrale sein Schl"usselpaar 
       $(e_i, d_i)$, wobei f"ur alle Teilnehmer der Modulus $n$ gleich ist.
       Dadurch entsteht das folgende Problem: eine Nachricht $m$ soll an 
       zwei Personen geschickt werden. Es werden die beiden Chiffrate $m^{e_1}$ 
       und $m^{e_2}$ berechnent und and die entsprechenden Personen 
       geschickt. Es gelte dabei $\ggt(e_1,e_2)=1$ f"ur die geheimen 
       Schl"ussel der Personen. Wenn Eve die beiden Chiffrat $m^{e_1}$ 
       und $m^{e_2}$ belauscht, kann sie 
       $$ r e_1 + s e_2 = 1 $$
       und damit auch 
       $$ (m^{e_1})^r (m^{e_2})^s = m^{r e_1 + s e_2} = m^1$$ 
       berechnen. Deswegen mu\3 f"ur jeden Teilnehmer ein eigenes 
       Paar von Primzahlen $p$ und $q$ gew"ahlt werden.
%% Diffie Hellman
\section{Diffie-Hellman-Verfahren}
  Das Verfahren von Diffie und Hellman erlaubt es, Alice und Bob ein gemeinsames
  Geheimnis "uber eine unsichere Leitung auszutauschen. Das Protokoll kann mit einer
  beliebigen endlichen zyklischen Gruppe realisiert werden. Die Gruppe $G$ und ihr
  Erzeuger $\alpha$ sind "offentlich bekannt.
  \begin{enumerate}
    \item Alice w"ahlt eine Zufallszahl $a$, berechnet $\alpha^a$ in $G$ und sendet
      $\alpha^a$ an Bob.
    \item Bob w"ahlt eine Zufallszahl $b$, berechnet $\alpha^b$ in $G$ und sendet
      $\alpha^b$ an Alice.
    \item Alice empf"angt $\alpha^b$ und berechnet $(\alpha^b)^a$.
    \item Bob empf"angt $\alpha^a$ und berechnet $(\alpha^b)^a$.
  \end{enumerate}
  Alice und Bob haben jetzt ein gemeinsames Gruppenelement $\alpha^{ab}$, das nur ihnen
  bekannt ist. Eve kennt durch Abh"oren nur die Gruppenelemente $\alpha^a$ und 
  $\alpha^b$. Es wird allgemein angenommen, da"s die Gruppenelemente $\alpha^a$ und
  $\alpha^b$ nicht ausreichen, um $\alpha^{ab}$ auszurechnen, wenn es praktisch
  unm"oglich ist, den diskreten Logarithmus von Gruppenelementen zu bestimmen.
  
  Das {\em Diffie-Hellman-Problem} ist es, einen effizienten Algorthmus zu finden,
  der $\alpha^{ab}$ aus $\alpha^a$ und $\alpha^b$ berechnet. Es ist klar, da"s die
  L"osung des diskreten Logarithmusproblem erlaubt, das Diffie-Hellman-Problem zu
  l"osen. Die Umkehrung ist nicht bekannt.

%% ************ ElGamal
  \section{ElGamal-Verfahren}
     Die Sicherheit des Verfahrens von ElGamal
     beruht auf der Schwierigkeit, den diskreten Logarithmus 
     in endlichen zyklischen Gruppen zu bestimmen.
     Es kann sowohl zum Chiffrieren als auch zum Signieren 
     von Nachrichten verwendet werden. 

     Sei $\alpha$ ein Erzeuger der Gruppe $G$. Die Gruppe $G$ und der
     Erzeuger seien "offentlich bekannt. Die Nachrichten werden durch
     Gruppenelemente dargestellt. Alice will Bob eine Nachricht $m$ senden.
     Bob w"ahlt eine Zufallszahl $x$ (sein geheimer Schl"ussel), berechnet
     $y=\alpha^x$ und ver"offentlicht das Gruppenelelement $y$
     (sein "offentlicher Schl"ussel). F"ur den Nachrichtenaustausch
     wird folgendes Protokol verwendet:
     \begin{enumerate}
       \item Alice w"ahlt eine Zufallszahl $k$ und berechnet $a=\alpha^k$.
       \item Alice holt sich Bobs "offentlichen Schl"ussel $y$ und
         berechnet $b=m y^k$.
       \item Alice sendet das Tupel $(a, b)$ an Bob.
     \end{enumerate}
     Bob kann die Nachricht $m$ berechnen, da er seinen geheimen Schl"ussel
     $x$ kennt und $a$ von Alice erh"alt.

     Typischerweise werden die Gruppen $\F_{2^m}^*$, $\F_p^*$ ($p$ Primzahl) oder
     die Gruppe $E(\F)$ der Punkte auf einer elliptischen Kurve "uber einem endlichen
     K"orper benutzt, da sie sich besonders einfach implementieren lassen. Die Gruppe 
     $E(\F)$ ist entweder eine zyklische Gruppe oder
     das direkte Produkt zweier zyklischer Gruppen. Normalerweise ist die eine Komponente des
     direkten Produkts eine kleine zyklische Gruppe und die andere eine gro"se,
     die dann zum Verschl"usseln benutzt wird.
     
     Ein Signaturverfahren l"a"st sich nur mit Primk"orper realisieren. Betrachten
     wir den Fall $G=\F_p^*$ genauer. Bob w"ahlt eine Primzahl $p$, ein primitives 
     Element $g$ und eine Zufallszahl $x$. Er berechnet $y \equiv g^x \pmod{p}$. 
     Sein "offentlicher Schl"ussel ist das Tupel $(y,g,p)$. Sein geheimer Schl"ussel 
     ist der diskrete Logarithmus von $y$ zur Basis $g$, die Zufallszahl $x$.
    
     \subsubsection*{Chiffrierverfahren: }
       Alice verschl"usselt, eine Nachricht $m$, indem sie $k$ zuf"allig mit $\ggt(p-1,k)=1$ 
       w"ahlt und $a \equiv g^k \pmod{p}$ und $b \equiv y^k m$ berechnet. Das Chiffrat
       ist das Tupel (a,b). Bob kann das Chiffrat $(a,b)$ entschl"usselt, indem er
       $$ m \equiv b (y^k)^{-1} \equiv b (g^{xk})^{-1} \equiv b ((g^k)^x)^{-1} \equiv b (a^x)^{-1} \pmod{p}$$
       mit seinem geheimen Schl"ussel $x$ ausrechnet.
     
     \subsubsection*{\bf Signaturverfahren:}
       Bob will eine Nachricht $m$ signieren. Die Idee bei der Signatur 
       ist es, die Nachricht $m$ als
       $$ m \equiv a x \pmod{p-1}$$
       auszudr"ucken. Die Signatur $a$ kann dann durch
       $$ g^{m} \equiv g^{a x} \equiv y^{a} \pmod{p}$$
       "uberpr"uft werden. Der geheime Schl"ussel $x$ kann aber aus der Signatur $a$
       $$ x \equiv m a^{-1} \pmod{p-1} $$
       berechnet werden. Deswegen wird eine lineare Verschiebung benutzt. Dazu 
       w"ahlt Bob $k$ zuf"allig mit $\ggt(k,p-1)=1$ und berechnet 
       $a \equiv g^k \pmod{p}$. Bob berechnet au"serdem noch ein $b$, so da"s
       $$ m \equiv x a + k b \pmod{(p-1)} $$
       gilt. Die Signatur ist dann das Tupel $(a,b)$. Andere Teilnehmer k"onnen 
       Bobs Signatur "uberpr"ufen, indem sie
       $$ g^m \equiv g^{x a + k b} \equiv g^{x a} g^{k b} \equiv y^a a^b \pmod{p} $$
       nachrechnen. Die Signatur funktioniert nur bei Primk"orper,
       weil sonst Probleme mit dem Datentyp auftreten w"urden. Man m"u"st mit einem 
       Gruppenelement potenzieren. Bei Primk"orper werden die Gruppenelmente 
       als nat"urliche Zahlen aufgefa"st. In den anderen F"alle ist dies nicht
       m"oglich.

       Es ist wichtig, da"s jedesmal $k$ verschieden gew"ahlt wird. Falls 
       $a=g^{k} \pmod{p}$ zweimal gleich ist, kann das Linearegleichungssystem
       \begin{eqnarray*}
         m_1 & = & x a + k b_1 \\
         m_2 & = & x a + k b_2 
       \end{eqnarray*}
       aufgestellt und der geheime Schl"ussel $x$ bestimmt werden.
     
%% ******** McElice-Kryptosystem    
   \section{McElice-Kryptosysstem}
     Das McElice-Kryptosystem benutzt fehlerkorrigierende Goppa Codes, die einfach zu 
     decodieren sind, und transformiert sie auf einen "`allgemeinen"' linearen Code. 
     Die Decodierung von allgemeinen linearen Codes ist ein NP-vollst"andiges Problem. 
     Das McElice-System ist folgenderma\3en aufgebaut:
     \begin{itemize}
       \item Geheimer Schl"ussel
         \begin{itemize}
           \item $G'$ ist eine $k\times n$-Matrix, die Generator-Matrix eines Goppa-Codes, 
                 der maximal $t$ Fehler korrigieren kann;
           \item $P$ eine $n\times n$-Permutationsmatrix;
           \item $S$ eine invertierbare $k\times k$-Matrix.
         \end{itemize}
       \item "Offentlicher Schl"ussel
         \begin{itemize}
           \item $G=S G' P$ eine $k\times n$-Matrix;
           \item die Nachricht ist ein Vektor der L"ange $k$ "uber $\{0,1\}$;
           \item verschl"usselt wird durch $c = mG + z$, wobei $z$ ein zuf"alliger Vektor vom Gewicht ein
             wenig kleiner als $t$ ist.
         \end{itemize}
       \item Entschl"usseln
          \begin{itemize}
            \item $c' = c P^{-1}$;
            \item man decodiert $m'$, so da\3 $m'G'=c'$;
            \item es gilt dann $m = m'S^{-1}$.
          \end{itemize}
      \end{itemize}
      Der Vorschlag von McEliece war
      \begin{itemize}
         \item $n = 1024$,
         \item $t \approx 50$ Fehler,
         \item maximal m"ogliches $k$, f"ur das ein Goppa-Code existiert, ist  $k=524$.
      \end{itemize}
      \subsubsection*{Bewertung}
      \begin{itemize}
        \item Der Schl"ussel hat etwa $2^{19}$ bit L"ange.
        \item Die Nachrichtenl"ange wird verdoppelt; der Rechenaufwand ist aber wegen der 
              einfachen Verschl"usselung 
              und Entschl"usselung sehr gering (deutlich schneller als RSA).
        \item Praktisch wird das Verfahren wegen der Schl"usselgr"o\3e kaum eingesetzt.
      \end{itemize}
      \subsubsection*{Angriffe}
      Direktes Decodieren braucht ungef"ahr $2^{500}$ Schritte. Dieser Aufwand l"a\3t 
      sich erheblich reduzieren, wenn man $k$ Spalten der Matrix (mit der Hoffnung, da\3 
      dort kein Fehler dazuaddiert wurde) ausw"ahlt und das lineare Gleichungssystem $mM_k=c_k$ 
      l"ost. Falls kein Fehler in diesen Spalten aufgetreten ist, kann man $m$ 
      durch Matrixinversion zur"uckgewinnen. Es m"ussen allerdings viele Matrizen getestet 
      werden, bis bei 524 Spalten keiner der 50 Fehler enthalten ist. Der Aufwand 
      betr"agt etwa $2^{81}$ Schritte.

 
\chapter{Mathematische Grundlagen}

Dieses Kapitel ist stark an das Skript von
Schaefer-Lorinser~\cite{frank} angelehnt. Hier werden f"ur das
Verst"andnis der Public Key Kryptographie unbedingt notwendigen
Teile aus~\cite{frank} behandelt. F"ur ein besseres Verst"andnis
der Primtests und Faktorisierungsalgorithmen ist die Lekt"ure des
ausf"uhrlichen und weitergehenden Skripts w"armstens empfohlen.

\section{Primzahlen}

In diesem Kapitel werden wir uns mit ein paar grundlegenden
Eigenschaften der Primzahlen besch"aftigen, insbesondere werden wir
ein paar S"atze kennenlernen, die die Verteilung der Primzahlen in den
nat"urlichen Zahlen beschreiben. Beginnen wir jedoch mit der
Definition von Primzahlen:

\begin{definition}[Primzahl]
Jede nat"urliche Zahl $n > 1$, die innerhalb der nat"urlichen Zahlen
nur durch sich selbst und durch 1 teilbar ist, hei\3t {\em Primzahl}.
\end{definition}

Die Primzahlen bilden die fundamentalen Atome der Arithmetik auf den
nat"urlichen Zahlen, da sich jede nat"urliche Zahl in eindeutiger
Weise als Produkt der Primzahlen darstellen l"a\3t. Dies wird auch als
der {\em Hauptsatz der Arithmetik} bezeichnet.

Um auf eine effiziente Weise eine Liste der Primzahlen kleiner einer
vorgegebenen Zahl $n$ zu erstellen, ist das auf den griechischen
Mathematiker Eratosthenes (276-196 v.Chr.) zur"uckgehendes Verfahren
{\em "`Sieb des Eratosthenes"'} geeignet.

Schon aus der Antike ist der von Euklid ca. 325 v.Chr. in seinen
Elementen notierte Satz bekannt:

%%
%% Satz von Euklid
%%
\begin{satz}[von Euklid]
Es gibt unendlich viele Primzahlen.
\end{satz}

W"ahrend uns dieser Satz also zeigt, da\3 in der Folge der
nat"urlichen Zahlen immer wieder eine Primzahl vorkommt, zeigt der
n"achste Satz, da\3 auch beliebig gro\3e L"ucken auftauchen k"onnen,
das hei\3t Abschnitte der nat"urlichen Zahlen, die keine Primzahl
enthalten.

\begin{satz}
Zu jedem $k \in \N$ gibt es in den nat"urlichen Zahlen einen Abschnitt
der L"ange $k$, der keine Primzahl enth"alt.
\end{satz}

\beweis "Ubung \hfill $\Box$

Das hei\3t also, die Primzahlen treten in gewisser Weise
unregelm"a\3ig auf. Es gibt keine geschlossene Formel, die beschreibt,
welche nat"urlichen Zahlen Primzahlen sind. Dennoch kann eine recht
pr"azise Aussage "uber die asymptotische Verteilung der Primzahlen
gemacht werden, wie der folgende von Tschebyscheff bewiesene Satz
zeigt. Wir definieren zun"achst:

\begin{definition}
Es bezeichne f"ur reelle Zahlen $x\in\R$ die Funktion $\pi(x)$ die
Anzahl der Primzahlen, die nicht gr"o\3er als $x$ sind.
\end{definition}
%%
%% Satz von Tschebyscheff
%%
\begin{satz}[Tschebyscheff] 
Es gibt $a,b \in \R_{>0}$, so da\3 f"ur $x\ge 2$ gilt
\begin{equation}
  a\frac{x}{\ln x} < \pi(x) < b\frac{x}{\ln x}.
\end{equation}
\end{satz}

\beweis "Ubung \hfill $\Box$

Betrachten wir einen Restklassenring $\ring{n}$ modulo einer nat"urlichen
Zahl $n$, dann ist unmittelbar klar, da\3 dieser Ring nicht die
K"orperaxiome erf"ullt, falls $n$ keine Primzahl ist. Denn unter
dieser Voraussetzung ist jeder echte Teiler von $n$ nicht
invertierbar. Die teilerfremden Zahlen $a$ mit $1 \le a \le n$, f"ur
die $\ggt(a,n)=1$ gilt, sind in diesem Restklassenring $\ring{n}$
invertierbar. Sie bilden die Einheitengruppe $\einheit{n}$, und die
Anzahl der Elemente dieser Gruppe wird durch die sogenannte Eulersche
$\Phi$-Funktion angegeben.

%%
%% Eulersche Phi-Funktion
%%
\begin{definition}[Eulersche $\Phi$-Funktion] 
  Die Funktion $\Phi(n)$ bezeichnet die Anzahl der nat"urlichen Zahlen 
  kleiner $n$, die relativ prim zu $n$ sind
  und wird die Eulersche $\Phi$-Funktion genannt:
  $$ \#\einheit{n} = \euler{n} = \#\{ m \mid 1 \le m < n, \ggt
(m,n)=1\} $$
\end{definition}   

Wie kann diese Funktion explizit berechnet werden? Betrachten wir
zun"achst den einfachen Fall, da\3 $n$ eine Primzahlpotenz $p^e$ ist.

%%
%% Phi-Funktion f"ur Primpotenz
%%
\begin{bem}[Primzahlpotenz $p^e$]
  Es sei $p$ eine Primzahl und $e \ge 1$. Dann gilt
  $$ \euler{p^e} = p^e - p^{e-1}. $$
\end{bem} 
\beweis
F"ur alle Produkte $s p$ mit $s = 1, \ldots ,p^{e-1}$ gilt $\ggt(s p,
p^e) \ne 1$. Damit sind die $p^e - p^{e-1}$ restlichen Zahlen teilerfremd
zu $p^e$. \hfill $\Box$

Es gilt ferner, da\3 die Eulersche $\Phi$-Funktion f"ur teilerfremde
Zahlen multiplikativ ist:

%%
%% Bem. Multipikativit"at der Phi-Funktion
%%
\begin{bem}[Multiplikativit"at der Eulerschen $\Phi$-Funktion]
  Es seien $m, n \ge 1$ zwei nat"urliche Zahlen mit $\ggt(m,n)=1$. 
  Dann gilt
  $$ \euler{m \cdot n} = \euler{m} \cdot \euler{n} .$$
\end{bem}

\beweis
Da $m, n$ teilerfremd gilt $\ggt(m \cdot n ,s)=1$ genau dann, wenn
$\ggt (m,s)$ und $\ggt(n,s)=1$. Aufgrund der Isomorphie nach dem
Chinesischem Restesatz
$$ \ring{mn} \cong \ring{m} \times \ring{n} $$
ergibt sich die Multiplikativit"at
$$ \euler{m \cdot n} = \euler{m} \cdot \euler{n}$$
unmittelbar, da
$$ \einheit{mn} \cong \einheit{m} \times \einheit{n}. $$
\hfill $\Box$

Aus diesen Bemerkungen folgt unmittelbar das Korollar, das die
$\Phi$-Funktion f"ur alle nat"urlichen Zahlen bestimmt.

%%
%% Kor. Phi-Funktion f"ur beliebige Zahlen
%%
\begin{kor}[Eulersche $\Phi$-Funktion f"ur beliebige Zahlen]
  Jede nat"urliche Zahl $n$ l"a\3t sich eindeutig als $n =
  \prod_{i=1}^{r} p_i^{e_i}$ in seine Primfaktoren zerlegen. Damit
  gilt
  $$ \euler{p_1^{e_1} \cdots p_n^{e_r}} = \prod_{i=1}^r \left(
  p_i^{e_i} - p_i^{e_i-1} \right). $$
\end{kor}

Der folgende Satz geht auf Fermat (1601-1665) zur"uck und wird oft 
als {\em der kleine Satz von Fermat} bezeichnet.

%%
%% Satz kleiner Fermat
%%
\begin{satz}[der kleine Satz von Fermat] \label{fermat}
  Es sei $p$ prim, dann gilt f"ur alle $b\in\einheit{p}$
  $$ b^{p-1} \equiv 1 \pmod{p} .$$
\end{satz}

Das folgende weitergehende Resultat geht auf Euler (1707-1783)
zur"uck. Es stellt eine Verallgemeinerung des kleinen Fermats dar, da
es auch die F"alle, in denen $p$ keine Primzahl ist, erfa\3t.

%%
%% Satz von Euler
%%
\begin{satz} 
Es seien $n,b\in\N$ mit $\ggt(b,n)=1$, das hei"st $b$ ist eine Einheit
in dem Ring $\ring{n}$. Dann gilt
$$ b^{\euler{n}} \equiv 1 \pmod{n} .$$
\end{satz}

Da f"ur $n$ prim gilt: $\euler{n} = n-1$, ist Satz~\ref{fermat} ein
Spezialfall davon.
%%
%% *********** Quadratische Reste
%%

\section{Quadratische Reste}

Wir wollen im folgenden die Frage untersuchen, wann sich eine Einheit $a$
des Restklassenrings $\ring{n}$ als Quadrat einer anderen Einheit $x$
darstellen l"a"st.

\begin{definition}[Quadratische Reste]
  Es seien $a,n \in \N$ mit $\ggt(a,n)=1$. Alle Elemente $a$, f"ur die
  ein $x$ existiert mit
  $$ x^2 \equiv a \pmod{n} $$ hei"sen {\em quadratische Reste} modulo
  $n$, die "ubrigen Einheiten hei"sen {\em quadratische Nicht-Reste}.
\end{definition}

F"ur den Fall, da\3 es sich bei dem Modulus $n$ um eine Primzahl
handelt, definieren wir das {\em Legendre Symbol}:

%%
%% Def. Legendre Symbol
%%
\begin{definition}[Legendre Symbol]
Es sei $p$ eine Primzahl und $a\in\N$. Dann hei\3t das durch 
\begin{equation}
  \leg{a}{p} := \left\{ 
                  \begin{array}{rl}
                        1, & \mbox{falls $a$ ein quadratischer Rest ist} \\
                       -1, & \mbox{falls $a$ kein quadratische Rest ist}
                  \end{array}
                \right.
\end{equation}
definierte Symbol $\leg{a}{p}$ das Legendre Symbol von $a$ nach $p$.
\end{definition}

Der folgende Satz gibt uns schon zwei wichtige Regeln f"ur
quadratische Reste bzw. quadratische Nicht-Reste.

%%
%% Satz
%%
\begin{satz} \label{wannQR}
Das Produkt von zwei quadratischen Resten ist wieder ein quadratischer
Rest, und das Produkt von einem quadratischen Rest und einem
quadratischen Nicht-Rest ist ein quadratischer Nicht-Rest.
\end{satz}

\beweis 
Der erste Teil des Beweises ist direkt einzusehen. Denn f"ur zwei
quadratische Reste $a$ und $b$ exisitieren $x$ und $y$ mit $x^2 \equiv
a \pmod{n}$ und $y^2 \equiv b \pmod{n}$. Daraus folgt aber $(x\cdot
y)^2 \equiv a\cdot b \pmod{n}$.

Den zweiten Fall beweisen wir durch einen Widerspruch. Angenommen, das
Produkt von einem Quadratischen Rest $x^2$ und einem Quadratischen Nicht-Rest 
$z$ w"are ein Quadratischer Rest
$y^2$. Daraus w"urde
$$ x^2 \cdot z \equiv y^2 \pmod{n}$$
folgen. Da $x$ als Einheit invertierbar ist, w"are dann $z$ wegen
$$ z \equiv (x^{-1} \cdot y)^2 \pmod{n}$$
ein Quadratischer Rest, was der Voraussetzung widerspricht. \hfill $\Box$

F"ur den Fall eines Produkts von zwei quadratischen Nicht-Resten gibt
es keine allgemeine Regel. F"ur den Fall aber, da\3 der Modulus $n$
eine Primzahl ist, ergibt dies immer einen quadratischen Rest.

%%
%% Satz Mulitplikativ
%%
\begin{satz} \label{legmultiplikativ}
  Im Falle von quadratischen Resten modulo einer Primzahl $p$ gilt:
  \begin{equation}
    \leg{a}{p} \leg{b}{p} = \leg{ab}{p}. \label{legendreMult}
  \end{equation}
  Ferner ist f"ur $p\ne 2$ die H"alfte der Elemente der Einheitengruppe 
  $\einheit{p}$ Quadratische Reste und die andere H"alfte Quadratische 
  Nicht-Reste.
\end{satz}

\beweis 
F"ur $p=2$ ist die Behauptung offensichtlich. Sei nun $p\ne 2$.  Die
Einheitengruppe $\einheit{p}$ ist eine zyklische Gruppe der Ordnung
$p-1$. Es gibt also einen Erzeuger $g$ der Ordung $p-1$ und
$$ g,g^2,\ldots,g^{p-1} $$ 
ist bereits die ganze $\einheit{p}$. Zu jedem Element
$a\in\einheit{p}$ existiert ein Index $i$ mit
$$ g^i \equiv a \pmod{p}. $$ Der Erzeuger $g$ ist kein Quadratischer
Rest, da sonst ein $x$ mit $x^2 \equiv g \pmod{p}$ existieren
w"urde. Nach dem Satz~\ref{fermat} w"urde f"ur dieses $x$
$$ x^{p-1} \equiv 1 \pmod{p} $$ gelten und somit $g^{(p-1)/2} \equiv 1
\pmod{p}$. Das ist aber ein Widerspruch dazu, da\3 $g$ die ganze
Einheitengruppe $\einheit{p}$ erzeugt.  Offensichtlich sind alle
geraden Potenzen von $g$ Quadratische Reste
$$ g^2,g^4,g^6,\ldots ,g^{p-1}$$
und da $g$ ein Quadratischer Nicht-Rest ist, sind nach Satz~\ref{wannQR}
alle ungeraden Potenzen von $g$ Quadratische Nicht-Reste
$$ g,g^3,g^5, \ldots ,g^{p-2}.$$ 
Damit ist die Formel~\ref{legendreMult} bewiesen, da das
Produkt von zwei ungeraden Potenzen eine gerade Potenz ergibt. 
\hfill $\Box$

Auf Euler geht das folgende Ergebnis zur"uck, das uns ein
Kriterium liefert, f"ur ein gegebenes Element $a\in\einheit{p}$ direkt
zu entscheiden, ob es ein quadratischer Rest oder nicht.
%%
%% ********* Euler Kriterium
%%
\begin{satz}[Euler Kriterium] 
Es sei $a\in\N$ mit $\ggt(a,p)=1$, wobei $p$ eine ungerade Primzahl
ist. Dann gilt
$$ \leg{a}{p} \equiv a^{\frac{p-1}{2}} \pmod{p}. \label{eulerkrit} $$
\end{satz}

\beweis 
Es sei $g$ ein Erzeuger der Einheitengruppe $\einheit{p}$. Wir
bestimmen zun"achst ein $s\in\left\{0,\ldots,p-2\right\}$ mit 
$$ g^s \equiv -1 \pmod{p} $$. 
F"ur ein solches $s$ gilt 
$$ g^{2s} \equiv (g^s)^2 \equiv (-1)^2 \equiv 1 \pmod{p}. $$ 
Dies kann nur gelten, wenn $2s$ ein Vielfaches
der Gruppenordnung $\#\einheit{p} = p-1$ ist, da $p-1$ der kleinste
Exponent $e$ ist, f"ur den $g^e \equiv 1 \pmod{p}$ gilt. Also kann $s$
nur einen der Werte
$$ 0,{p-1 \over 2}, p-1,\ldots $$
annehmen. Es gilt wegen Satz~\ref{fermat}
$$ g^0 \equiv g^{p-1} \equiv 1 \pmod{p}. $$
Also mu\3 $s={p-1 \over 2}$ sein und somit 
$$ g^{(p-1)/2} \equiv -1 \pmod{p}. $$ 
Da $g$ ein Erzeuger der Einheitengruppe $\einheit{p}$ ist, kann jedes Element
$a\in\einheit{p}$ als $a \equiv g^i \pmod{p}$ dargestellt
werden. Damit gilt
$$ a^{(p-1)/2} \equiv (g^i)^{(p-1)/2} \equiv
\left(g^{(p-1)/2}\right)^i \equiv (-1)^i \pmod{p}. $$ Das hei\3t aber,
da\3 $a^{(p-1)/2} \equiv 1 \pmod{p}$ genau dann gilt, wenn $i$ ein
gerader Exponent ist, und $a^{(p-1)/2} \equiv -1 \pmod{p}$ genau dann,
wenn $i$ ein ungerader Exponent ist. Dies ist aber das Kriterium f"ur Quadratische
Reste aus Satz~\ref{legmultiplikativ}. \hfill $\Box$

\begin{kor} \label{modulo}
Es sei $p$ eine ungerade Primzahl und es gelte
$\ggt(a,p)=\ggt(b,p)=1.$ Aus $a \equiv b \pmod{p}$ folgt, da\3
$\leg{a}{p} = \leg{b}{p}$.
\end{kor}

\beweis 
Es ist $a \equiv b \pmod{p}$ gleichbedeutend damit, da\3 ein $l\in \Z$
existiert mit $b=a+l\cdot p$. Da modulo $p$ gerechnet wird, gilt:
\begin{eqnarray*}
  \leg{b}{p} & \equiv & \leg{a+l\cdot p}{p} \\
             & \equiv & (a+l\cdot p)^{(p-1)/2} \\
             & \equiv & \sum_{i=0}^{p-1} {p-1 \choose 2i} a^{(p-1)/2-i}(lp)^i \\
             & \equiv & a^{(p-1)/2} \pmod{p}.
\end{eqnarray*}
\rule{0pt}{0pt}\hfill $\Box$
                     
Ein sehr interessantes Ergebnis der Zahlentheorie, das wir ohne Beweis
angeben, ist {\em das quadratische Reziprozit"atsgesetz}, das von
Gau\3 bewiesen wurde. Es stellt eine Relation zwischen dem
Legendresymbol von zwei Primzahlen $\leg{p}{q}$ und der "`Umkehrung"'
$\leg{q}{p}$ dar.

%%
%% Quadratisches Reziprozit"atsgesetz
%%
\begin{satz}[Quadratisches Reziprozit"atsgesetz] 
Seien $p$ und $q$ zwei ungerade Primzahlen. Dann gilt:
  \begin{equation}
    \leg{p}{q} = \leg{q}{p} (-1)^{{1 \over 2}(p-1){1 \over 2}(q-1)} .
  \end{equation}
\end{satz}

Mit Hilfe des quadratischen Reziprozit"atssatzes und des
Korollars~\ref{modulo} l"a\3t sich effizient entscheiden, ob eine
gegebene Zahl $a$ modulo einer Primzahl $p$ ein Quadratischer Rest ist.

%%
%% Beispiel
%%
\begin{beispiel} Ist 19 ein Quadratischer Rest modulo 61? Es gilt
$$\leg{19}{61}=\leg{61}{19}(-1)^{30 \cdot 9}=\leg{61}{19}=\leg{4}{19} = 1, $$
da $4 \equiv 2^2 \pmod{19}$.
\end{beispiel}

%%
%% ***************** Jacobi Symbol
%%
Beim Legendre-Symbol sind Quadratische Reste modulo einer Primzahl $p$
zugelassen. Im folgenden wird eine Verallgemeinerung des
Legendre-Symbols, das Jacobi-Symbol, das auch f"ur zusammengesetzte
Zahlen definiert ist, eingef"uhrt.

%%
%% Def.
%%
\begin{definition}[Jacobi-Symbol]
  Es seien $a,n\in\N$ mit $n$ ungerade und $\ggt(a,n)=1$. Die
  Primfaktorzerlegung von $n$ sei gegeben durch $n=\prod_{i=1}^{r}
  p_i^{e_i}$. Dann ist das Jacobi-Symbol definiert durch
  $$\leg{a}{n} := \prod_{i=1}^{r} \leg{a}{p_i}^{e_i} .$$
\end{definition}

Das Jacobi-Symbol ist so definiert, da\3 es den gleichen Rechenregeln
wie das Legendre-Symbol gen"ugt. Es ist aber zu beachten, da\3 das
Jacobi-Symbol $\leg{a}{n}$ {\em nicht} angibt, ob $a$ ein
quadratischer Rest modulo $n$ ist.

Es ist leicht einzusehen, da\3 f"ur ein $a$, das ein Quadratischer
Rest modulo einer zusammengesetzten Zahl $n=\prod_{i=1}^{r}p_i^{e_i}$
ist, das Jacobi-Sybmol $\leg{a}{n}$ den Wert 1 annimmt. Da $a$ ein
Quadratischer Rest modulo $n$ ist, existiert ein $x\in\N$ mit $x^2
\equiv a \pmod{n}$. Betrachten wir diese Gleichung modulo $p_i$ f"ur
$i=1,\ldots,r$, so erhalten wir $ x^2 \equiv a \pmod{p_i}$. Es ist
also $a$ auch Quadratischer Rest modulo aller Primfaktoren $p_i$ 
f"ur $i=1,\ldots,r$ und damit ist das Jacobi-Symbol $\leg{a}{n}$ 
gleich 1.

Die Umkehrung dieser Aussage gilt nicht, wie das folgende Beispiel zeigt.

\begin{beispiel}
  Das Jacobi-Symbol kann auch f"ur Quadratische Nicht-Reste den Wert 1 annehmen:
  $$ \leg{2}{15} = \leg{2}{3}\leg{2}{5} = (-1)(-1)=1, $$
  obwohl 2 kein Quadratischer Rest modulo 15 ist.
\end{beispiel}    

%%
%% ************** Primtest
%%
\section{Primtests}

Die Primzahltests weisen eine ganze Reihe von Methoden und Parallelen
zu den Faktorisierungsalgorithmen auf

\subsection{Pseudoprimzahlen}

Betrachten wir zun"achst einen einfachen Satz, der ein Kriterium
liefert, mit dem sich entscheiden l"a\3t, ob eine Zahl $n$ prim oder
zusammengesetzt ist.

%%
%% Satz Prim-Kriterium
%%
\begin{satz}
   Eine Zahl $n$ ist prim genau dann, wenn
   \begin{equation} 
     a^{n-1} \equiv 1 \pmod{n}, \label{primkrit} 
   \end{equation}
   f"ur alle $a\in\N$ mit $2 \le a < n$.
\end{satz}

\beweis 
Die Hinrichtung gilt aufgrund des Satzes~\ref{fermat}.

F"ur die R"uckrichtung nehmen wir an, die
Fermat-Bedingung~\ref{primkrit} sei erf"ullt und $n$ sei
zusammengesetzt. Dann hat $n$ einen nichttrivialen Teiler $r$ mit $1 <
r < n$ und $n=r\cdot s$, und es gilt
$$ r^{n-1} \equiv 1 \pmod{n}. $$
Es existiert also ein $l\in\Z$ mit
$$ 1=r^{n-1}+l\cdot n = r^{n-1}+l\cdot s\cdot r = r^{n-1} + (l\cdot s) \cdot r. $$
Betrachten wir diese Gleichung modulo $r$, so ergibt sich 
$$ 0 \equiv 1 \pmod{r}. $$ 
Das ist ein Widerspruch. \hfill $\Box$

Mit Hilfe dieses Satzes k"onnen wir bereits einen einfachen Primtest realisieren.
Haben wir eine Zahl $n$ gegeben, von der wir entscheiden m"ochte, ob sie prim
oder zusammengesetzt ist, so w"ahlen wir eine Basis $a$ mit
$2\le a < n$ und testen die Fermat-Bedingung~\ref{primkrit}.
Erf"ullt die Zahl $n$ die Bedingung nicht, so ist $n$ sicher
eine zusammengesetzte Zahl. Andernfalls w"ahlen wir eine neue Basis
und f"uhren den Test erneut durch.

Dies w"urde, falls $n$ eine Primzahl ist, einen extrem aufwendigen
Primtest ergeben. Wir k"onnen aber nach einer gr"o\3eren Anzahl von
Durchl"aufen des Primtests schon mit einer gewissen Wahrscheinlichkeit
schlie\3en, da\3 die Zahl $n$ prim ist.

\begin{definition}
  Eine zusammengesetzte Zahl $n$ hei\3t {\em Pseudoprimzahl zur Basis $a$}, falls
  $$ a^{n-1} \equiv 1 \pmod{n} $$
  erf"ullt ist.
\end{definition}

\begin{beispiel}  
  Die Zahl $341=11\cdot 31$ ist eine Pseudoprimzahl zur Basis 2, denn
  $$ 2^{341-1} \equiv 1 \pmod{341}. $$
\end{beispiel}

\begin{satz}
  Es sei $n\in\N$ eine Pseudoprimzahl zur Basis $a_1$ und zur Basis
  $a_2$ mit $a_1, a_2\in\einheit{n}$. Dann ist $n$ auch eine
  Pseudoprimzahl zur Basis $a_1 \cdot a_2$ und zur Basis $a_1 \cdot
  a_2^{-1}$.

  Falls ein $a\in \einheit{n}$ existiert, das die Bedingung
  $$ a^{n-1} \equiv 1 \pmod{n} $$ nicht erf"ullt, so wird sie auch von
  mindestens der H"alfte der Elemente aus $\einheit{n}$ nicht
  erf"ullt.
\end{satz}

\beweis 
Der erste Teil ist durch direktes Nachrechnen zu beweisen. Es gilt
$$(a_1\cdot a_2)^{n-1} \equiv a_1^{n-1} \cdot a_2^{n-1} \equiv 1
\pmod{n},$$ und damit ist $n$ eine Pseudoprimzahl zur Basis $a_1
a_2$. Weiterhin gilt
$$(a_1 \cdot a_2^{-1}) \equiv a_1^{n-1} \cdot (a_2^{n-1})^{-1} \equiv
1 \cdot 1^{-1} \equiv 1 \pmod{n},$$ und damit ist $n$ eine
Pseudoprimzahl zur Basis $a_1 a_2^{-1}$.

Es sei $B:=\{a_1,\ldots, a_l\}$ die Menge aller $a\in\einheit{n}$ mit
$a^{n-1} \equiv 1 \pmod{n}$.  Die zusammengesetzte Zahl $n$ ist also
eine Pseudoprimzahl zur Basis $a_i\in B$ f"ur $i=1,\ldots, l$. Wir
w"ahlen $b\in\einheit{n}\setminus B$ beliebig aus. Wenn $n$ eine
Pseudoprimzahl zur Basis $b \cdot a_i$ w"are, dann w"are nach dem
ersten Teil $n$ auch eine Pseudoprimzahl zur Basis $(b \cdot a_i)
a_i^{-1} \equiv b \pmod{n}$. Dies kann aufgrund der Wahl von $b$ aber
nicht sein, und damit sind alle Elemente aus $\{ b \cdot a_1, \ldots,
b \cdot a_n\}$ Basen, zu denen $n$ keine Pseudoprimzahl ist. Damit
sieht man, da\3 mindestens die H"alfte der Elemente aus $\einheit{n}$
Basen sind, zu denen $n$ keine Pseudoprimzahl ist. \hfill $\Box$

Falls also zu einer zusammengesetzten Zahl $n$ "uberhaupt eine Zahl
$a\in\einheit{n}$ existiert, so da\3 $n$ keine Pseudoprimzahl zur
Basis $a$ ist, so finden wir mit einer Wahrscheinlichkeit gr"o\3er als
$1/2$ eine solche Basis. Eine solche Basis $a$ wird auch {\em Zeuge}
f"ur $n$ genannt, weil die Angabe von $a$ gen"ugt, um zu beweisen,
da\3 $n$ eine zusammengesetzte Zahl ist.

W"ahlen wir nun $k$ verschiedene Kandidaten f"ur die Basis $a$ auf
zuf"allige Weise, so ist $a$ f"ur eine zusammengesetzte Zahl $n$ nur mit
Wahrscheinlichkeit kleiner als $1/2^{k}$ unter diesen Elementen kein
Zeuge. Diese Aussage gilt wieder nur unter der Voraussetzung, da\3 es
f"ur $n$ "uberhaupt einen Zeugen gibt.

Es gibt Zahlen, f"ur die das nicht erf"ullt ist. Sie werden {\em
Carmichael-Zahlen} genannt. Wenn $n$ eine Carmichael-Zahl ist, dann
gilt
$$ a^{n-1} \equiv 1 \pmod{n}$$ f"ur alle $a\in\einheit{n}$.

Die Zahl $561=3\cdot 11\cdot 17$ ist beispielsweise eine Carmichael-Zahl.

\begin{satz}
  Es seien $m,n\in\Z$ ungerade. F"ur das Jacobi-Symbol gelten folgende
  Rechenregeln: \\
  (i)   $$ \leg{-1}{m} = (-1)^{(m-1)/2}, $$
  (ii)  $$ \leg{2}{m} = (-1)^{(m^2-1)/8}, $$
  (iii) F"ur $m,n$ teilerfremd gilt:
  $$ \leg{m}{n} = \leg{n}{m} (-1)^{\frac{m-1}{2}\frac{n-1}{2}}. $$
\end{satz}


%% ******** Primzahltests
\subsection{Probabilistischer Primtest nach Solovay, Strassen}

Im vorigen Abschnitt haben wir aus dem kleinen Fermat eine Art
heuristischen Primtest abgeleitet, indem wir f"ur die auf Primalit"at
zu testende Zahl $n$ "uberpr"uften, ob sie die Eigenschaft~\ref{primkrit}, die der
kleine Fermat f"ur alle Primzahlen liefert, erf"ullt. Dieser einfache
Primtest versagt mit sehr hoher Wahrscheinlichkeit bei den
Carmichael-Zahlen. Wir w"ahlen in dem Primtest eine Zufallszahl $a$ mit
$1 < a < n$. Ist $\ggt(a,n)\neq 1$, so hat man einen Nullteiler
gefunden, und wei\3 sofort, da\3 die Zahl $n$ zusammengesetzt ist. Wie
man aber leicht aus der Berechnungsvorschrift f"ur die Eulersche
$\Phi$-Funktion sieht, ist die Menge der Nullteiler im Vergleich zu
der Menge der Einheiten sehr klein. Ist $n$ eine Carmichael-Zahl, so
ist $n$ eine Pseudoprimzahl zu allen Einheiten aus $einheit{n}$.

Die Grundidee des folgenden Primtests ist es, einen analogen Test
bez"uglich des Euler-Kriteriums vorzunehmen.

%%
%% Def.
%%
\begin{definition}
  Analog zur Definition von Pseudoprimzahlen nennen wir
  zusammengesetzte Zahlen $n$, f"ur die
  $$ \leg{a}{n} \equiv a^{(n-1)/2} \pmod{n} $$ gilt, {\em
  Euler-Pseudoprimzahlen} zur Basis $a$.
\end{definition}

Das folgende einfache Lemma zeigt schon, da\3 dieser Primtest den auf
dem kleinen Fermat basierenden Primtest umfa\3t.

%%
%% Lemma
%%
\begin{lemma}
  Wenn $n$ eine Euler-Pseudoprimzahl zur Basis $a$ ist, dann ist $n$
  auch eine Pseudoprimzahl zur Basis $a$.
\end{lemma}

\beweis 
Quadrieren der Gleichung im Satz~\ref{eulerkrit} liefert die
Behauptung. \hfill $\Box$

Der auf dieser "Uberlegung basierende Primtest wurde 1974 von Solovay
und Strassen vorgeschlagen. Man nennt Algorithmen dieser Art
probabilistisch, da es zusammengesetzte Zahlen gibt, f"ur die der Test
behauptet, die Zahl $n$ sei mit einer gewissen Wahrscheinlichkeit
prim. Durch Wiederholen des Primtests kann die
Irrtumswahrscheinlichkeit exponentiell verringert werden. In der
Praxis treten Zahlen, f"ur die sich der Test auch nach mehrmaligem
Anwenden irrt, "au\3erst selten auf.

Der Primzahltest nach Solovay und Strassen war das erste
probabilistische Verfahren "uberhaupt. Heute werden in vielen
Bereichen der Algorithmentechnik probabilistische Methoden eingesetzt.

\subsection*{Algorithmus:}
  \begin{enumerate}
    \item Bestimme eine zuf"allig gew"ahlte Zahl $a$ mit $1 < a < n$.
    \item Teste, ob $\leg{a}{n} \equiv a^{(n-1)/2} \pmod{n}$ gilt.
    \item Falls ja, dann ist $n$ wahrscheinlich prim, sonst ist $n$
          mit Sicherheit zusammengesetzt.
  \end{enumerate}

Das Jacobi-Symbol kann hierbei offensichtlich nicht "uber die
Definiton berechnet werden, da dies die Primfaktorzerlegung von $n$
voraussetzt. Es l"a"st sich effizient mit den Rechenregeln ausrechnen.

Der gro\3e Vorteil dieses Primtests liegt nun darin, da\3 es f"ur
Euler-Pseudoprimzahlen nichts "Aquivalentes zu den Carmichael-Zahlen
gibt. Ferner l"a\3t sich zeigen, da\3 die Irrtumswahrscheinlichkeit
dieses Test kleiner als $1/2$ ist.

\begin{satz}[Solovay, Strassen 1974] 
Es sei $P:=\left\{a \mid \ a\in\einheit{n},\leg{a}{n}\equiv
a^{(n-1)/2}\right\}$.
  \begin{itemize}
    \item Falls $n$ prim ist, so gilt $\#P=n-1$.
    \item Falls $n$ nicht prim ist, so gilt $\#P \le {1 \over 2}(n-1)$.
  \end{itemize}
\end{satz}
\beweis 
Falls $n$ eine Primzahl ist, so liefert das Euler-Kriterium die
Richtigkeit der Behauptung. Es sei nun $n$ eine zusammengesetzte Zahl
mit der Primfaktorzerlegung $n=\prod_{i=1}^{r} p_i^{e_i}$. Wir
definieren zur Vereinfachung zwei Abbildungen

\begin{equation}
  \epsilon : \left\{ \begin{array}{ccc}
                      \einheit{n} & \rightarrow & \einheit{n} \\
                      a           & \mapsto         & \leg{a}{n}
                    \end{array} \right.
\end{equation}
und 
\begin{equation}
  \delta   : \left\{ \begin{array}{ccc}
                      \einheit{n} & \rightarrow & \einheit{n} \\
                      a           & \mapsto         & a^{(n-1)/2}
                    \end{array} \right. .
\end{equation}

Es ist also zu zeigen, da\3 $\#\{a\in\einheit{n} \mid \epsilon(a) =
\delta(a)\} \le {1 \over 2}(n-1)$ erf"ullt ist. Wir machen hierzu eine
Fallunterscheidung.

\begin{enumerate}
%% 1. Fall
\item Es gibt ein $j$ mit $e_j > 1$ ($n$ ist nicht quadratfrei).

      Angenommen es sei $\epsilon(a)=\delta(a)$ f"ur alle 
      $a\in\einheit{n}$. Dann gilt 
      \begin{equation}
        a^{n-1} \equiv \left(a^{(n-1)/2}\right)^2 \equiv \delta(a)^2
        \equiv \epsilon(a)^2 \equiv \leg{a}{n}^2 \equiv 1 \pmod{n}
        \label{gleich}
      \end{equation}
      f"ur alle $a\in\einheit{n}$. Es gilt ferner 
      \begin{eqnarray*}
        \#\einheit{n} & = & \euler{n} \\
                      & = & \euler{p_1^{e_1}} \cdots \euler{p_r^{e_r}} \\
                      & = & p_1^{(e_1-1)}(p_1-1) \cdots p_r^{(e_r-1)}(p_r-1).
      \end{eqnarray*}

      Daraus folgt $p_j \mid \#\einheit{n}$, da $e_j > 1$. Also gibt
      es ein Element $a\in\einheit{n}$ der Ordnung $p_j$, d.h. $p_j$
      ist der kleinste Exponent, f"ur den $a^{p_j} \equiv 1 \pmod{n}$
      gilt. Da $p_j$ kein Teiler von $n-1$ ist, gilt $a^{n-1}
      \not\equiv 1 \pmod{n}$.  Dies ist aber ein Widerspruch zu
      Gleichung~\ref{gleich}. Es ist also $\epsilon \neq \delta$ und
      damit ist die Menge
      $\{a\in\einheit{n}\mid\epsilon(a)=\delta(a)\}$ eine echte
      Untergruppe (sie ist multiplikativ und enth"alt das neutrale
      Element) von $\einheit{n}$.  Damit gilt f"ur die
      Untergruppenordung
      $$\#\{a\in\einheit{n}\mid\epsilon(a)=\delta(a)\} \le {1 \over 2}\euler{n} $$
      und somit die Behauptung, da $\euler{n} \le n-1$ immer erf"ullt ist.
%% 2. Fall
\item Es gilt $e_i=1$ f"ur alle $i=1,\ldots, r$ ($\ring{n}$ ist isomorph 
      zu dem direkten Produkt der Primk"orper $\ring{p_1},\ldots ,\ring{p_r}$).

      Angenommen, es sei $\epsilon(a)=\delta(a)$ f"ur alle
      $a\in\einheit{n}$. W"ahle $a_1$ mit $\leg{a_1}{p_1}=-1$ und
      $a_i$ mit $\leg{a_i}{p_i}=1$ f"ur
      $i=2,\ldots,r$. Nach dem Chinesischem Restesatz existiert ein
      $a\in\Z$ mit $a\equiv a_i \pmod{p_i}$ f"ur $i=1,\ldots,r$. Es
      gilt nach Definition des Jacobi-Symbols und
      Korollar~\ref{modulo}
      $$ \epsilon(a) = \leg{a}{n} = \prod_{i=1}^{r} \leg{a}{p_i} = 
         \prod_{i=1}^r \leg{a_i}{p_i} = -1 \cdot 1^{r-1} = -1 $$
      und nach der Annahme ist $\delta(a)=-1$. Betrachten wir den 
      durch den Chinesischen Restesatz gegebenen Isomorphismus
      $$ \phi :  \ring{n} \longrightarrow \ring{p_1} \times \ldots 
         \times \ring{p_r}.$$

      Da $\phi$ insbesondere ein Homomorphismus ist und aus $x \equiv
      -1 \pmod{n}$ f"ur alle Primteiler $p_i$ von die Gleichungen $x
      \equiv -1 \pmod{p_i}$ folgen, gilt speziell

      \begin{equation} 
         -1 = \delta(a) = a^{(n-1)/2} \mapsto 
         \left(a_1^{(n-1)/2},\ldots,a_r^{(n-1)/2}\right) = (-1,\ldots,-1). 
         \label{minuseins} 
      \end{equation}

      Nach dem Chinesischen Restesatz existiert ein $b\in\Z$ mit
      $b\equiv 1 \pmod{p_1}$ und $b\equiv a_i \pmod{p_i}$ f"ur
      $i=2,\ldots,r$. Damit gilt
      $$ \delta(b) = b^{(n-1)/2} \mapsto \left( 1^{(n-1)/2},
         a_2^{(n-1)/2},\ldots,a_r^{(n-1)/2)} \right) =
         (1,-1,\ldots, -1). $$

      Es kann also $b^{(n-1)/2} \equiv \pm 1$ auf keinen Fall gelten,
      da $b^{(n-1)/2}$ durch den Isomorphismus $\phi$ auf
      $(1,-1,\ldots,-1)$ abgebildet wird und da aufgrund der
      Eigenschaft des neutralen Elements und der
      Formel~\ref{minuseins} gilt

      \begin{eqnarray*}
         1 & \mapsto & \phi(1)  = (1,1,\ldots,1) \\
        -1 & \mapsto & \phi(-1) = (-1,-1,\ldots,-1).
      \end{eqnarray*}

      Das Jacobi-Symbol $\leg{b}{n}$ kann nur die Werte $1$ oder $-1$ annehmen. 
      Es gilt aber $\leg{b}{n}=\epsilon(b)\neq \pm 1$. Also ist auch 
      $\epsilon \neq \delta$ in diesem Fall gezeigt. Mit
      dem gleichen Untergruppenkriterium wie in ersten Fall folgt die
      Behauptung. \hfill $\Box$
\end{enumerate}


%% ************* Rabin und Miller
\subsection{Probabilistischer Primtest nach Rabin und Miller}

Der folgende Primtest wurde von Rabin aufbauend auf den
Ideen von Miller entwickelt. Er liefert eine Verbesserung des
Primtests von Solovay und Strassen, da die Menge der F"alle, in
denen er sich irrt, eine echte Untermenge der F"alle ist, in denen sich
der Primtest von Solovay und Strassen irrt.

Der Algorithmus geht von folgender grunds"atzlicher "Uberlegung aus:
Nehmen wir an, wir h"atten eine Pseudoprimzahl $n$ zur Basis $a$, also
mit
$$ a^{n-1} \equiv 1 \pmod{n}, $$ 
gegeben. Betrachten wir Quadratwurzeln aus $a^{n-1}$, so w"urde f"ur eine 
Primzahl $n$ gelten, da\3 die Quadratwurzel $\pm 1$ sind, da in einem K"orper 
zwei L"osungen f"ur $X^2 \equiv 1 \pmod{n}$ existieren. Das hei\3t, f"ur 
eine Primzahl $n$ w"urde die Folge
$$ a^{(n-1)/2}, a^{(n-1)/2^2},\ldots,a^{(n-1)/2^k} $$ 
f"ur $n-1 = 2^k m$ mit $m$ ungerade aus entweder beliebig vielen 1en, gefolgt von
einer $-1$ und dann beliebigen Zahlen aus $\einheit{n}$ oder lauter
1en bestehen. Diese Beobachtung wird in dem Miller-Rabin-Test
ausgenutzt. Zur Berechnung geht man hier nat"urlich von
$a^{(n-1)/2^k}$ aus und quadriert.

\subsection*{Algorithmus}

\begin{enumerate}
  \item Schreibe $n-1 = 2^k \cdot m$ mit $m$ ungerade.
  \item W"ahle eine Zufallszahl $a$ mit $1<a<n$ und $\ggt(a,n)=1$.
  \item Initialisiere $i:=0$ und $b:=a^m$.
  \item Falls $(i=0$ und $b=1)$, dann ist $n$ wahrscheinlich prim. Stop
  \item Falls $b=-1$, dann ist $n$ wahrscheinlich prim. Stop
  \item Falls $(i>0$ und $b=1)$, dann ist $n$ sicher zusammengesetzt. Stop \\
        Sonst $b:=b^2$ und $i:=i+1$
  \item Falls $i=k+1$, dann ist $n$ sicher zusammengesetzt. Stop \\ 
        Sonst gehe zu Schritt 5.
\end{enumerate}

Der folgende Satz liefert eine Absch"atzung der
Fehlerwahrscheinlichkeit dieses Algorithmus.

%%
%% Satz Rabin
%%
\begin{satz}[Rabin]
  Es sei $n\in\N$ ungerade und mit $n-1=2^k \cdot m$, wobei $m$
  ungerade ist. Zu $a\in\N$ mit $0<a<n$ definieren wir eine Folge
  $a_0,\ldots,a_k$ mit
  $a_i := a^{2^i \cdot m}$ f"ur $0\le i\le k$.
  Es sei 
    $$ P:=\left\{a\in\einheit{n} \mid (a_0 \equiv 1 \pmod{n}) \mbox{
       oder } (\exists i<k : a_i \equiv -1 \pmod{n}\right\}. $$
  \begin{itemize}
    \item Falls $n$ prim, so gilt $\#P = n-1$.
    \item Falls $n$ nicht prim, so gilt $\#P \le \frac{1}{4}(n-1)$.
  \end{itemize}
\end{satz}  

\beweis 
siehe \cite{frank} \hfill $\Box$

Der Algorithmus erkennt die Primzahlen mit einer Wahrscheinlichkeit
gr"o\3er als 3/4. Durch $t$-faches Wiederholen des Algorithmus kann
die Fehlerwahrscheinlichkeit auf $1/4^t$ gesenkt werden.

Der n"achste Satz zeigt, da\3 der Algorithmus von Rabin auch wirklich
eine Verbesserung des Algorithmus von Solovay und Strassen darstellt.
\begin{satz} 
  Es sei 
  $$ P':=\left\{a\in\einheit{n} \mid \leg{a}{n}\equiv a^{(n-1)/2}\right\}$$
  und
  $$ P:=\left\{a\in\einheit{n} \mid (a_0 \equiv 1 \pmod{n}) \mbox{ oder } 
     (\exists i<k : a_i \equiv -1 \pmod{n}\right\}, $$
  wobei $0<a<n$, $n$ ungerade und zusammengesetzt mit $n-1 = 2^k \cdot m$ und 
  $(a_i)_i$ eine 
  Folge mit $a_0,\ldots, a_k$. Dann gilt f"ur zusammengesetzte
  Zahlen $n$: $P \subset P'$. 
\end{satz}

\beweis 
siehe \cite{frank} \hfill $\Box$

\section{Faktorisierungsalgorithmen mit Gruppen}

Den beiden in diesem Abschnitt vorgestellten Algorithmen liegt eine
gemeinsame Idee zugrunde. Der chinesische Restesatz besagt, da\3 f"ur
zusammengesetztes $n$ der Restklassenring $\ring{n}$ isomorph ist zu einem 
direkten Produkt von kleineren Restklassenringen.  Wir definieren eine 
algebraische Struktur, die von dem Restklassenring $\ring{n}$ abh"angt. Wenn wir
in dieser Struktur rechnen, so rechnen wir implizit auch in der algebraischen
Struktur, die gewisserma\3en von den kleineren Komponenten des
direkten Produkts abh"angt, und erh"alten aufgrund des kleinen Fermats
Information "uber eine der kleinere Strukturen und k"onnen dadurch einen
Faktor von $n$ abspalten. Im Falle der $p-1$-Methode ist die
algebraische Struktur $\ring{n}$ selbst und bei der Faktorisierung mit
elliptischen Kurven eine elliptische Kurve $E(\ring{n})$ "uber dem
Restklassenring $\ring{n}$.

\subsection{Pollard $p-1$-Methode}

Der Aufwand dieses Algorithmus h"angt im wesentlichen von der Gr"o\3e
des zu findenden Faktors ab. Nehmen wir also an, die Zahl $n$, die
faktorisiert werden soll, enth"alt einen Primfaktor $p$, der die
Eigenschaft besitzt, da\3 $p-1$ in kleine Faktoren zerf"allt. Konkret
sollen alle Faktoren von $p-1$ kleiner als eine Schranke $k$ sein, und
es soll dar"uberhinaus gelten, da\3 $k!$ durch $(p-1)$ teilbar ist.

Da die Exponentiation modulo $n$ effizient berechnet werden kann, kann
die Berechnung von
$$ m:\equiv 2^{k!} \pmod{n} $$ f"ur nicht zu gro\3e $k$ und $n$
durchgef"uhrt werden. Der kleine Fermat liefert wegen $p-1 \mid k!$,
da\3
$$ m \equiv 2^{k!} \equiv 2^{k! \pmod{p-1}} \equiv 2^0 \equiv 1
\pmod{p} $$ gilt. Damit teilt $p$ also $m-1$. Falls $k$ nicht zu gro\3
angesetzt ist, teilt $n$ mit einer hohen Wahrscheinlichkeit $m-1$
nicht, so da\3 der gr"o\3te gemeinsame Teiler
$$ g:=\ggt(m-1,n) $$
den nicht trivialen Teiler $p$ von $n$ liefert.

Bei der praktischen Anwendung dieser Vorgehensweise, wei\3 man
nat"urlich nicht, in welcher Gr"o\3enordung der richtige Parameter $k$
gew"ahlt werden sollte. Damit nicht alle Primfaktoren von $n$ erfa\3t
werden, ist es wichtig von Zeit zu Zeit $\ggt(m-1,n)$ zu
berechnen. Falls der ggT gleich $1$ ist, so erh"oht man $k$ und
rechnet weiter.

Sollte der ggT schon $n$ sein, so sind alle Primfaktoren von $n$ schon
erreicht worden, so da\3 man entweder mit kleinerem $k$ oder mit einer
anderen Basis als 2 einen neuen Versuch macht.

Der Algorithmus f"uhrt nicht zum Erfolg, falls die Primfaktoren $p_i$ 
von $n$ eine "ahnliche Zerlegung f"ur $p_i-1$ haben. Dann gilt
$(p_i-1)\mid k!$. In diesem Fall wird ganz $n$ abgespalten.

\subsubsection{Verfeinerung des Algorithmus}

Anstatt $k!$ als Exponent zu verweden, ist es g"unstiger eine etwas
andere Wahl zu treffen, denn f"ur sehr kleine Primzahlen $q$ taucht
$q$ in $k!$ mit einem sehr hohem Exponenten auf. Es ist aber
unwahrscheinlich, da\3 dies f"ur die Bedingung $p-1\mid k!$ ben"otigt
wird. In der Praxis kann man beispielsweise so vorgehen, da\3 man bis
zu einer gewissen Schranke $B_1$, alle Primzahlen und Primpotenzen
kleiner $B_1$ betrachtet. An Stelle von $k!$ w"ahlt man dann das Produkt
dieser Primzahlen und Primpotenzen. Das Produkt sei $l$.

Man w"ahlt eine weitere Schranke $B_2 > B_1$ und testet
zus"atzlich zu 
$$ \ggt(2^l-1,n) $$
auch noch
$$ \ggt(2^{l \cdot q_i}-1,n) $$
f"ur alle Primzahlen $q_i$ mit $B_1 < q_i < B_2$, um
die Primfaktoren zu trennen.

\subsubsection{Laufzeitbetrachtungen}

Die Laufzeit dieses Algorithmus h"angt nicht nur von der Gr"o\3e der
zu faktorisierenden Zahl $n$ ab, sondern wird auch stark von der
Zerlegung von $p-1$ f"ur alle Primfaktoren $p$ beeinflu\3t. Die
Primfaktorzerlegung von $n$ ist im voraus nicht bekannt. Daher kann
man die Laufzeit des Algorithmus nur mit einer
Wahrscheinlichkeitsbetrachung absch"atzen. Es sei $p$ der kleinste
Primfaktor von $n$. Die Laufzeit h"angt dann nur von der gr"o\3ten
Primpotenz von $p-1$ ab, sofern $p$, wie im allgemeinen zu erwarten
ist, als erster Faktor gefunden wird.

Betrachten wir $p-1$ als zuf"allig gew"ahlte nat"urliche Zahl, dann
hat der gr"o\3te Primfaktor im Mittel eine Gr"o\3enordung von
$(p-1)^{0.63}$. Falls also die gr"o\3te Primpotenz kleiner der
vorgegeben Schranke $B$ ist, so f"uhrt der Algorithmus zum Erfolg,
wenn die Primfaktoren weit genug auseinander liegen.

%%
%% ************** Faktorisierung mit elliptischen Kurven
%%
\subsection{Faktorisierung mit elliptischen Kurven}

Der Ansatz, mit elliptischen Kurven zu faktorisieren, geht auf
H.W. Lenstra zur"uck. Dieser Ansatz stellt in gewisser Hinsicht eine
Verallgemeinerung der $p-1$-Methode von Pollard dar. Der Aufwand beim
Faktorisieren mit elliptischen Kurven ist im wesentlichen auch von der
Gr"o\3enordnung des gesuchten Primfaktors abh"angig.


\subsubsection{Elliptische Kurven "uber K"orpern}

Es sei $\K$ ein K"orper. Im Hinblick auf die betrachteten Anwendungen
gen"ugt es K"orper $\K$ der Charakteristik $\mbox{char}(\K)\neq 2,3$
zu betrachten. Dann lassen sich alle elliptischen Kurven beschreiben
durch die Gleichungen der Form
\begin{equation} 
  y^2 = x^3 + Ax + B, 
\end{equation}
wobei $A,B\in\K$ mit $4A^3 + 27B^2 \neq 0$. Diese Nebenbedingung
sichert, da\3 das Polynom $x^3 + Ax + B$ keine doppelten Nullstellen
enth"alt.

F"ur $\K=\R$ lassen sich die elliptischen Kurven geometrisch
veranschaulichen. Interessant werden die
elliptischen Kurven $E(\R)$ dadurch, da\3 man auf ihnen in einem
abstrakten Sinne addieren kann. Zu je zwei Punkten $P_1$ und $P_2$ der
Kurve findet man durch eine geometrische Konstruktion stets einen
dritten, der ebenfalls auf der Kurve liegt und $P_1 + P_2$ genannt
wird. Die so definierte Addition folgt den "ublichen Rechenregeln: Sie
ist assozitativ, kommutativ, es existiert ein neutrales Element, und zu
jedem Punkt existiert ein Inverses. Eine elliptische Kurve bildet also
eine kommutative Gruppe bez"uglich dieser Addition.

Das Inverse $-P$ eines Kurvenpunktes $P$ ist dessen Spiegelbild an der
$x$-Achse. Die Summe zweier Punkte $P_1$ und $P_2$ findet man im
allgemeinen, indem man die Gerade durch die Punkte $P_1$ und $P_2$
zieht. Sie schneidet die Kurve in genau einem dritten Punkt
$Q$. Dessen Spiegelbild $-Q$ bez"uglich der $x$-Achse ist die gesuchte
Summe $P_1 + P_2$.

In dem Sonderfall $P_1=P_2$ tritt an die Stelle der Geraden durch
$P_1$ und $P_2$ die Tangente an die Kurve im Punkt $P_1=P_2$. Wenn
schlie\3lich $P_1=-P_2$ ist, trifft die Gerade durch $P_1$ und $P_2$
keinen weiteren Punkt der Kurve. Man stellt sich ersatzweise vor,
da\3 die Gerade die Kurve im Unendlichen schneidet.  Darum f"ugt man
noch einen weiteren Punkt hinzu, den man sich im Unendlichen denkt.
Man nennt ihn ${\cal O}$ und setzt $P_1-P_1={\cal O}$. Der Punkt ${\cal
O}$ ist das neutrale Element der additiven Gruppe $E(\R)$.

Die geometrische Defintion der Punktaddition kann in Formeln umgesetzt
werden. Wenn $P_1=(x_1,x_2)$ und $P_2=(x_2,y_2)$ Punkte der
elliptischen Kurve sind, dann gilt
\begin{equation} -P_1 = (x_1,-y_1) , 
\end{equation}
\begin{equation} P_1 + P_2 = 
                    \left\{ 
                       \begin{array}{cl}
                         {\cal O} & \mbox{ falls } P_1 = -P_2 \\
                         P_2      & \mbox{ falls } P_1 = {\cal O} \\
                         P_1      & \mbox{ falls } P_2 = {\cal O}
                       \end{array}
                    \right.
\end{equation}

In den anderen F"allen ist die Summe $P_3 = P_1 + P_2 = (x_3, y_3)$
definiert durch
\begin{equation}
\begin{array}{ccc}
  x_3 & = & \lambda^2 - x_1 - x_2                     \\
      &   &                                           \\
    y_3 & = & \lambda\left( x_1 - x_3\right) - y_1
\end{array}
\end{equation}
wobei 
$$ \lambda = \frac{y_1 - y_2}{x_1 - x_2} $$
falls $P_1\neq P_2$, und
$$ \lambda = \frac{3 x_1^2 + A}{2y_1}$$
falls $P_1 = P_2$.

Diese Formeln k"onnen nun eine Art Eigenleben gewinnen. Man darf sie
anwenden, ohne sich unter $x$ und $y$ Koordinaten von Punkten in der
Ebene vorstellen zu m"ussen. Es kann ein beliebiger
endlicher K"orper $\F$ verwendet werden. Geraden und Tangenten haben dann 
nicht mehr viel mit geometrischer Anschauung zu tun.

Betrachten wir nun n"aher die elliptischen Kurven "uber Primk"orpern
$\F_p$. Offensichtlich ist $E(\F_p)$ endlich. Man schreibt
"ublicherweise die Anzahl der Elemente in der Form
$$ \# E(\F_p) = p+1-t, $$ wobei $t$ als die Spur bezeichnet wird. F"ur
die Spur gilt die sogenannte Hasse-Schranke
$$ -2\sqrt{p} < t < 2\sqrt{p}. $$ 
Die Gruppe $E(\F_p)$ ist entweder eine zyklische Gruppe oder das
Produkt zweier zyklischer Gruppen.


%%
%% ************* "uber Ringen
%%

\subsubsection{Elliptische Kurven "uber Ringen}

Die Definition einer elliptischen Kurve "uber einem Restklassenring
ist ganz analog zur Definition der elliptischen Kurven "uber einem
K"orper.

\begin{definition}
  Eine elliptische Kurve "uber dem Restklassenring $\ring{n}$ ist definiert
  als die Menge der Paare $(x,y)\in\ring{n} \times \ring{n}$, die die
  Weierstrass-Gleichung
  $$ y^2 = x^3 + Ax + B $$
  mit $A$,$B\in\ring{n}$ und $4A^3 + 27B^2 \neq 0$ erf"ullen zusammen
  mit einem Punkt im Unendlichen ${\cal O}$. Sie wird mit
  $E(\ring{n})$ bezeichnet. Wir setzen dabei voraus, da\3 6 eine
  Einheit im Ring $\ring{n}$ ist.
\end{definition}

Auf dieser Punktmenge wird analog zur Addition auf elliptischen Kurven
"uber K"orpern eine Operation definiert. Eine direkte Verwendung der
oben definierten Additionsformeln ist nicht m"oglich, da die dort
auftretenden inversen Elemente nicht notwendigerweise existieren, da
$\ring{n}$ nicht nullteilerfrei ist, falls $n$ nicht prim ist.

Es sei zur Vereinfachung $n$ als quadratfrei vorausgesetzt, also $n=p_1
\cdot\ldots\cdot p_r$ f"ur paarweise verschiedene Primzahlen $p_i$
f"ur $i=1,\ldots,r$.  Dann existieren die Abbildungen
$$ h_i : E(\ring{n}) \longrightarrow E(\ring{p_i}), $$
die komponentenweise mittels der Projektionen
$$ \pi_i : \ring{n} \longrightarrow \ring{p_i} $$ 
f"ur $i = 1,\ldots,l$ definiert sind. Durch Zusammensetzen dieser
Abbildungen $h_i$ erh"alt man eine Abbilung
$$ h : E(\ring{n}) \longrightarrow E(\ring{p_1}) \times 
   \ldots \times E(\ring{p_l}), $$
die den gesamten Zielraum aussch"opft, bis auf alle Tupel, in denen das
neutrale Element einer der Kurven $E(\Z_{p_i})$ vorkommt und die nicht
das Nulltupel sind. Dabei
definieren wir die Abbildung $h$ so, da\3 das neutrale Element ${\cal
O}$ auf das Tupel bestehend aus den neutralen Elementen der Kurven
$E(\ring{p_1})$ bis $E(\ring{p_r})$ abgebildet wird. Der Chinesische
Restesatz liefert, da\3 allen Elementen im Wertebereich dieser
Funktion ein Urbild in eindeutiger Weise zugeordnet werden kann. F"ur
diese F"alle ist damit auch eine Operation auf $E(\ring{n})$ eindeutig
definiert, in dem man die zu verkn"upfenden Elemente zun"achst mittels
$h$ nach $E(\ring{p_1}) \times \ldots \times E(\ring{p_r})$ abbildet,
dort koordinatenweise addiert und gegebenfalls dem so erhaltenen Tupel
wieder sein Urbild unter $h$ zuordnet.

Die so definierte Operation kann auch durch Rechnen modulo $n$ direkt
in $\ring{n}$ ausgef"uhrt werden. Da diese Verkn"upfung, falls sie
definert ist, der Addition auf elliptischen Kurven "uber K"orpern
entspricht, bezeichnet man sie als {\em Pseudoadditon}. Die mit dieser
Operation versehene Kurve $E(\ring{n})$ bildet offensichtlich keine
Gruppe, da die Operation nicht total definiert ist. Es gilt f"ur jeden
Punkt $P\in E(\ring{n})$
$$ N_n \cdot P = {\cal O} \mbox{ mit } N_n := \kgv(\# E(\ring{p_1}),\ldots,\# E(\ring{p_r})). $$


\subsubsection{Algorithmus}

"Ahnlich wie bei den Primtests f"uhrt man beim Faktorisieren mit
elliptischen Kurven eine Berechnung aus, die gelingen w"urde, falls
$n$ eine Primzahl w"are.  Weil $n$ aber zusammengesetzt ist, kann die
Rechnung scheitern. Im Verlauf der Berechnung mu\3 man inverse Elemente
berechnen, indem man den gr"o\3ten gemeinsamen Teiler von $n$ mit
anderen Zahlen ausrechet. Die Rechnung kann nur dann fortgesetzt
werden, wenn der gr"o\3te gemeinsame Teiler gleich $1$ ist, also die
Zahlen teilerfremd sind. Falls der gr"o"ste gemeinsame Teiler dann auch
noch ungleich $n$ ist, hat man einen echten Teiler von $n$ gefunden.

Die Idee bei diesem Faktorisierungsalgorithmus ist es die Gruppe
$\einheit{n}$, die bei der $p-1$-Methode von Pollard verwendet wird,
durch die allgemeinere Gruppe der rationalen Punkte "uber den
elliptischen Kurven zu ersetzen. Auch hier ist der Algorithmus
erfolgreich, falls $\# E(\F_p)$ in kleine Primfaktoren zerf"allt. Der
wesentliche Vorteil liegt jedoch darin, da\3 zu einem vorgegeben $n$
verschiedene Gruppen mit verschiedener Ordungung aus dem Intervall
$$ \left[ p+1-2\sqrt{p},p+1+2\sqrt{p} \right] $$
zur Verf"ugung stehen. 

F"ur den Algorithmus kann man konkret wie folgt vorgehen. Zu einer
Schranke $b$ w"ahlt man "ahnlich wie bei der $p-1$-Methode von Pollard
einen Multiplikator $k:=b!$ oder einen "ahnlichen Multiplikator, der
nur aus Faktoren kleiner $b$ besteht. Dann w"ahlt man sich eine
elliptische Kurve $E(\ring{n})$ und einen Punkt $P$ auf der Kurve. In
der Praxis kann man dabei so vorgehen, da\3 man zun"achst nur den
Kurvenparameter $A$ vorgibt, dann einen Punkt $P$ zuf"allig w"ahlt und
dann den Kurvenparameter $B$ so bestimmt, da\3 der Punkt $P$ auf der
Kurve liegt. Dann berechnet man $k\cdot P$ auf der Kurve. Diese
Addition kann sehr effizient berechnet werden, da es nicht n"otig ist
$k$ mal $P$ zu addieren. Man kann "ahnlich wie bei "`square and
multiply"' vorgehen.  Die Addition f"uhrt man wie in der
definierten Pseudoaddition durch. Bei der Berechnung m"ussen Inverse
in $\ring{n}$ bestimmt werden, die in $\ring{n}$ nicht immer
existieren m"ussen. Ein Nullteiler hat einen nichttrivalen Teiler mit $n$. 
In diesem Fall liefert die Berechnung des gr"o\3ten gemeinsamen Teilers einen
Teiler von $n$.

Der Zufall spielt bei diesem Algorithmus wirklich mit, da es sehr viele
elliptische Kurven zur Auswahl gibt, an denen man die Berechnungen
durchf"uhren kann.  Lenstra konnte zeigen, da\3 es f"ur jedes
zusammengesetzte $n$ Kurven gibt, die einen Teiler liefern. Falls man
keinen nichttrivialen Teiler findet, kann man ein neues $k$ bestimmen
oder eine neue elliptische Kurve w"ahlen.

Der Aufwand dieses Faktorisierungsalgorithmus l"a\3t sich unter
gewissen bisher unbewiesenen heuristischen Annahmen wie folgt
absch"atzen. Es sei $p$ der kleinste Faktor der zu faktorisierenden
Zahl $n$. Es sei ferner $g$ eine beliebige nat"urliche Zahl. Dann
findet der Algorithmus mit Wahrscheinlichkeit $ \ge 1-e^{-g}$ den
Faktor $p$ und hat einen asymptotischen Aufwand von
$$ g \cdot \exp\left({\sqrt{(2+o(1))(\log p\log\log p)}} \left(\log n\right)^2\right) $$
f"ur $p\rightarrow \infty $.

Der Algorithmus eignet sich in hervorragender Weise zur
Parallelisierung. Man w"ahlt f"ur jeden zur Verf"ugung stehenden
Prozessor eine elliptische Kurve und rechnet auf allen Kurven solange
parallel, bis ein Faktor gefunden wurde. Dabei ist nur minimale
Vorbereitung und Nachbereitung auf jedem der Prozessoren notwendig.

%%
%% Faktorisierung mit quadratischen Kongruenzen
%%
\section{Faktorisierung mit quadratischen Kongruenzen}

Der Aufwand der Faktorisierungsalgorithmen mit Gruppen h"angt von der
Gr"o\3e der Primfaktoren der zu faktorisierenden Zahl $n$ ab.  F"ur
die Suche nach gr"o\3eren Teilern benutzt man andere Algorithmen, die
auf einer anderen Grundidee basieren. Die Laufzeit dieser Algorithmen
h"angt nur von $n$ ab.  Wir werden hier zwei Vertreter dieser
Algorithmenklasse kennenlernen.

Man will eine zusammengesetzte Zahl $n$ faktorisieren, indem man eine
quadratische Kongrunenz der Art
\begin{equation}
  x^2 \equiv y^2 \pmod{n} 
\end{equation}
aufstellt. Diese Kongruenz ist gleichbedeutend mit der Aussage, da\3
$n$ die Differenz $x^2 - y^2$ teilt. Nach der dritten binomischen
Formel ist
$$ x^2 - y^2 = (x-y)(x+y). $$ Wenn $n$ nicht selbst Teiler von $x-y$
oder $x+y$ ist, so mu\3 ein echter Teiler von $n$ in $x-y$ und ein
anderer echter Teiler in $x+y$ enthalten sein.  Durch die Berechnung
von $\ggt(x+y,n)$ kann dann ein nicht trivialer Faktor gefunden
werden.

Die Tatsache, da\3 sich "uber solche Kongruenzen ein Faktor finden
l"a\3t, ist schon l"anger bekannt. Sie l"a\3t sich praktisch aber nur
nutzen, wenn man geeignete Verfahren kennt, solche Kongruenzen zu
finden. Dies wurde f"ur nicht triviale F"alle erst durch die
Verwendung von leistungsf"ahigen Computern m"oglich.

Zun"achst werden wir den relativ einfachen Algorithmus von Dixon
betrachten, an dem sich eine ganze Reihe von grundlegenden Ideen
besonders gut erkl"aren lassen.  Er nimmt sofern eine Sonderstellung
ein, als da\3 f"ur diesen Algorithmus der subexponentielle Aufwand mit
mathematischer Strenge bewiesen werden kann, w"ahrend f"ur die
meisten anderen Varianten der Aufwand immer unter Verwendung
unbewiesener, lediglich auf heuristische Annahmen gest"utzte
Vermutungen abgesch"atzt wird.

%%
%% Grundlagen der Laufzeitanalyse
%%

\subsection{Grundlagen der Laufzeitanalyse}

Bevor wir uns den eigentlichen Algorithmen zuwenden, f"uhren wir
einige Definitionen ein und stellen einige S"atze vor, die uns die
Absch"atzung der Laufzeit der Faktorisierungsalgorithmen erleichtern
werden.

Wir benutzen im folgenden die Abk"urzung
\begin{equation} 
  L(n) := \exp\left(\sqrt{\log n\log\log n}\right). 
\end{equation}
Ferner benutzen wir die Abk"urzung $L^{\alpha}$ f"ur die Funktionenklasse
\begin{equation} L(n)^{\alpha + o(1)}. \end{equation}

\begin{lemma} F"ur diese Funktionenklasse gelten die folgenden Rechenregeln:
  \begin{itemize}
    \item $2 L^{\alpha} = L^{\alpha}$;
    \item $L^{\alpha} + L^{\beta} = L^{\max\{\alpha,\beta\}}$.
  \end{itemize}
\end{lemma}

\beweis 
"Ubung \hfill $\Box$

\begin{lemma} 
F"ur die Funktion $\pi(x)$, die die Anzahl der Primzahlen kleiner $x\in\R$ angibt, gilt
  \begin{equation}
    \pi(L^{\alpha}) = L^{\alpha} 
  \end{equation}
\end{lemma}

\beweis 
"Ubung \hfill $\Box$

Eine wichtige Formel, die bei der Laufzeitanalyse ben"otigt wird, ist
die Absch"atzung der Anzahl der nat"urlichen Zahlen, die nur
Primfaktoren kleiner einer vorgegebenen Schranke haben. Solche Zahlen
werden als {\em glatt} (engl. {\em smooth}) bzw. {\em $b$-glatt} bezeichnet,
wobei $b$ die Schranke f"ur die Gr"o\3e der Primfaktoren ist. Dieser
Begriff wurde von R.Rivest eingef"uhrt.

\begin{definition}
  Die Funktion $\glatt{x}{y}$ gibt die Anzahl der Elemente $n<x$ an,
  die glatt bez"uglich der Schranke $y$ sind:
  \begin{equation}
    \glatt{x}{y} := \#\{ n\in\N \mid n < x \mbox{ und } 
     (( p \mid n, p \mbox{ prim } ) \Rightarrow p < y )\}.
  \end{equation}
\end{definition} 

\begin{satz}
  F"ur eine beliebige Konstante $\epsilon$ sei $D:=\{(x,y) \mid x>10
  \mbox{ und } y \ge (\log x)^{1+\epsilon}\}$. Dann gilt f"ur alle
  Paare $(x,y)\in D$
  $$ \glatt{x}{y} = x \cdot u^{-u+o(u)} $$
  wobei $u:=\frac{\log x}{\log y}$.
\end{satz}

\begin{lemma} \label{spezialfall}
  F"ur den f"ur die Laufzeitanalyse wichtigen Spezialfall 
  $\glatt{n}{L^{\alpha}}$ gilt
  \begin{equation}
    \glatt{n}{L^{\alpha}} = n \cdot L^{-\frac{1}{2\alpha}}.
  \end{equation}
\end{lemma}

%%
%% Dixon Algorithmus
%%

\subsection{Dixons Algorithmus}

Der Algorithmus von Dixon ist f"ur praktische Zwecke nicht brauchbar,
aber verschiedene Grundideen k"onnen anhand dieses Algorithmus
besonders einfach erkl"art werden.  Au\3erdem handelt es sich um den
historisch ersten Algorithmus, f"ur den eine subexponentielle Laufzeit
exakt nachgewiesen werden konnte.

\subsubsection{Beschreibung des Algorithmus}

Es sei $n$ eine zusammengesetzte Zahl, die nicht durch Primzahlen
$p<L(n)$ teilbar ist. Zu $m\in\Z$ bezeichne $Q(m)$ den kleinsten
nicht-negativen Rest von $m^2 \pmod{n}$. Es sei $a$ eine Konstante $0
< a < 1$, die wir sp"ater optimal festlegen werden. Um eine
quadratische Kongruenz modulo $n$ aufzustellen, ben"otigen wir den
folgenden Teilalgorithmus:

%% ****** Teilalgorithmus
\begin{enumerate}  
  \item W"ahle $m\in\{1,\ldots,n\}$ zuf"allig aus und berechne $Q(m)$.

  \item Teste, ob $Q(m)$ nur Faktoren $p\le L^a$ besitzt. Falls dies
        der Fall ist, so faktorisiere $Q(m)$ vollst"andig und trage
        $m$, $Q(m)$ und das Tupel $(e_1,\ldots, e_r)$, wobei
        $Q(m)=\prod_{i=1}^{r} p_i^{e_i}$ und $p_1,\ldots,p_r$ die
        Primzahlen kleiner $L^a$ sind, in einer Tabelle ein.
\end{enumerate}

Dieser Teilalgorithmus wird so oft ausgef"uhrt, bis $\pi(L^a)+1$
Quadrate $Q(m)$, die $L^a$-glatt sind, gefunden wurden. Zu jedem
Eintrag haben wir einen Vektor
$$ v(m) \in\koerper{2}^{\pi(L^a)}, $$
dessen $i$-te Komponente eine 0 oder eine 1 ist, je nachdem, ob die
Primzahl $p_i$ in gerader oder in ungerader Potenz in $Q(m)$
auftritt.

Da wir $\pi(L^a)+1$ viele Vektoren der L"ange $\pi(L^a)$ haben,
existiert eine lineare Abh"angigkeit unter diesen Vektoren. Wir
rechnen "uber dem Skalark"orper $\koerper{2}$, das hei\3t uns
interessiert nur, ob der Exponent gerade oder ungerade ist. Nach
Umnummerierung k"onnen wir die lineare Abh"angigkeit als
$$ v(m_1)+v(m_2)+\ldots +v(m_k) = {\bf 0} .$$
ausdr"ucken. Es ist also das Produkt $Q(m_1) \cdot Q(m_2) \cdot \ldots
\cdot Q(m_k)$ ein Quadrat modulo $n$ und es
existiert folglich ein $x\in\Z$ mit
$$ Q(m_1) \cdot Q(m_2) \cdot \ldots \cdot Q(m_k) \equiv x^2 \pmod{n}. $$

Dieses $x$ kann sehr einfach mit Hilfe der im 2-ten Schritt
gespeicherten Tupel berechnet werden.  Ferner l"a\3t sich $y\equiv m_1
\cdot m_2 \cdot \ldots \cdot m_k \pmod{n}$ berechnen und es gilt
hierf"ur:
$$ y^2 \equiv m_1^2 \cdot m_2^2 \cdot \ldots \cdot m_k^2 \equiv Q(m_1) 
   \cdot Q(m_2) \cdot \ldots \cdot Q(m_k) \equiv x^2 \pmod{n} $$
Wir haben also eine quadratische Kongruenz modulo $n$
aufgestellt. Durch Berechnen von $\ggt(x+y,n)$ erhalten wir mit
Wahrscheinlichkeit gr"o\3er $1/2$ einen nichttrivialen Faktor von $n$.
Falls $x \neq \pm y \pmod{n}$ liefert die Berechnung von $\ggt(x+y,n)$
sicher einen nichttrivialen Faktor von n. Das folgende Lemma besagt
genauer:

\begin{lemma} 
  Es seien $x,y\in\ring{n}$ mit $x^2 \equiv y^2 \pmod{n}$
  gegeben. Dann folgt daraus mit Wahrscheinlichkeit $1-1/2^{t-1}$,
  da\3 $x \ne \pm y \pmod{n}$, wobei $t$ die Anzahl der Faktoren von
  $n$ ist.
\end{lemma}

\beweis 
"Ubung \hfill $\Box$

\begin{bem}
  Man sieht also, da\3 bei dem RSA-Verfahren die Verwendung von mehr
  als zwei Primzahlen keine Sicherheitsvorteile bringt.
\end{bem}

Um die Erfolgswahrscheinlichkeit bei diesem Algorithmus zu erh"ohen,
berechnet man anstatt $\pi(L^a)+1$ beispielsweise $\pi(L^a)+10$ solche
Eintr"age und erh"alt damit 10 verschiedene lineare Abh"angigkeiten,
von denen dann eine mit sehr hoher Wahrscheinlichkeit zu der gesuchten
Zerlegung von $n$ f"uhrt.

%%
%% Laufzeitanalyse
%%

\subsection{Laufzeitanalyse}

Wir f"uhren jetzt eine grobe Laufzeitanalyse durch, bei der auch der
Parameter $a$ optimal bestimmt wird.

In dem Teilalgorithmus m"ussen wir f"ur $\pi(L^a)=L^a$ viele
Primzahlen testen, ob sie $Q(m)$ teilen. Wir definieren $b:=b(n)$
indirekt dadurch, da\3 wir sagen, wir ben"otigen $L^b$ Durchl"aufe des
Teilalgorithmus, bis wir $\pi(L^a)+1$ viele geeignete Eintr"age
gefunden haben.  Dies ergibt einen Aufwand von $L^{a+b}$ Schritten,
bis wir gen"ugend viele gute $Q(m)$ gefunden haben.  Die
Gau\3elimination ben"otigt naiv implementiert $L^{3a}$ Schritte. Die
restlichen Berechnung der Werte $x$, $y$ und $\ggt(x+y,n)$ ben"otigt
nur $L^a$ viele Schritte.  Also kann die gesamte Laufzeit des
Algorithmus durch
$$ L^{\max\{ a+b,3a\}} $$
abgesch"atzt werden.

Wir m"ussen zur Bestimmung der Laufzeit die Funktion oder Konstante
$b$ ermitteln und $a$ optimal w"ahlen. Wir benutzen die heuristische
Annahme, da\3 bei einer zuf"alligen Wahl von $m$ auch das modulo $n$
reduzierte Quadrat $Q(m)$ zuf"allig verteilt ist. Damit betr"agt die
Wahrscheinlichkeit, da\3 $Q(m)$ in Faktoren $p < L^a$ zerf"allt,
$$ \frac{\glatt{n}{L^a}}{n}, $$
da $\glatt{n}{L^a}$ die Anzahl der $L^a$-glatten Zahlen kleiner $n$
angibt und $\frac{1}{n}$ unter der heuristischen Annahme die
Wahrscheinlichkeit jeder der m"oglichen Zahlen ist. Mit
Lemma~\ref{spezialfall} sehen wir, da\3
$$ \frac{\glatt{n}{L^a}}{n} = L^{-\frac{1}{2a}} $$
gilt. Um also $\pi(L^a) + 1 = L^a$ viele $L^a$-glatte $Q(m)$ 
zu finden, m"ussen wir
$$ L^a \cdot \left( L^{-\frac{1}{2a}} \right)^{-1} = L^{a+\frac{1}{2a}} $$
viele Werte f"ur $m$ untersuchen. Also gilt f"ur die Anzahl $b$ der
notwendigen Durchl"aufe des Teilalgorithmus
$$ b = a + \frac{1}{2a} , $$
und somit betr"agt der Aufwand
$$ L^{a+b} = L^{2a+\frac{1}{2a}}. $$
Die Funktion $2a + \frac{1}{2a}$ hat ihr Minimum bei $a=1/2$. Dies
ergibt eine Laufzeit f"ur den gesamten Algorithmus von
$$ L^{\max\{ a+b, 3a \}} = L^{\max\{ 2a + \frac{1}{2a}, 3a\}} = 
   L^{\max\{ 2, 1.5\}} = L^2, $$
das hei\3t es werden $L(n)^{2+o(1)}$ viele Schritte und ein
Speicheraufwand von $L(n)^{1+o(1)}$ f"ur die $L^a \times L^a$-Matrix
ben"otigt.

\begin{satz}
  Der Faktorisierungsalgorthmus von Dixon ben"otigt asymptotisch eine
  Laufzeit von $L(n)^{2+o(1)}$ vielen Schritten und einen
  Speicherplatzbedarf von $L(n)^{1+o(1)}$ bit.
\end{satz}

Unser Beweis dieses Satzes ist nicht mathmatisch exakt durchgef"uhrt
worden, da wir eine heuristische Annahme benutzt haben. Diese Annahme
betrifft die Wahrscheinlichkeit, mit der ein $Q(m)$ f"ur ein zuf"allig
gew"ahltes $m$ in Primfaktoren $p < L^a$ zerf"allt. Die $Q(m)$ sind nicht 
zuf"allig gleichverteilt, was wir angenommen haben.
Es l"a\3t sich jedoch zeigen, da\3 innerhalb einer Abweichung, 
die in $L(n)^{o(1)}$ abgefangen werden kann, dies auch f"ur die 
Quadrate $Q(m)$ zutrifft.

%%
%% Quadratischer Siebalgorithmus
%%

\subsection{Quadratischer Siebalgorithmus}

Ein weiterer wichtiger Ansatz, wie eine Kongruenz der Art
$$ x^2 \equiv y^2 \pmod{n} $$ 
aufgestellt werden kann, geht auf Pomerance zur"uck und verwendet ein
"ahnliches Verfahren, wie es beim Sieb des Eratosthenes eingesetzt
wird. Da mit Hilfe eines solchen Sieb-Verfahrens eine Quadratzahl
gefunden werden soll, spricht man von einem {\em quadratischen Sieb}.

In dem quadratischen Siebalgorithmus gehen wir von einem quadratischen
Polynom
\begin{equation} 
   f(x) = (x + \unten{\sqrt{n}})^2 - n \label{quadpolynom} 
\end{equation}
aus. F"ur dieses quadratische Polynom mit ganzzahligen Koeffizienten
gilt f"ur beliebige $x\in\Z$
$$ \left(x+\unten{\sqrt{n}}\right)^2 \equiv f(x) \pmod{n}. $$

In dem quadratischen Siebalgorithmus betrachten wir statt $Q(m)$ die
Werte $f(m)$, die in einer vorgegebenen Menge von $B$ Primzahlen
faktorisieren. Finden wir $m_1,\ldots,m_k \in\N$, so da\3
$$ f(m_1)\cdot \ldots \cdot f(m_k) $$
ein Quadrat ist, so erhalten wir mit
$$ x^2 := f(m_1)\cdot\ldots\cdot f(m_k) $$
und 
$$ y:=(m_1 + \unten{\sqrt{n}}) \cdot\ldots\cdot (m_k + \unten{\sqrt{n}}) $$
eine quadratische Kongruenz
$$ x^2 \equiv y^2 \pmod{n}, $$
mit der wir einen Faktor von $n$ bestimmen k"onnen. Die Zahl $x$ kann
wie beim Dixon-Algorithmus bestimmt werden, indem man die
gespeicherte Faktorisierung der $f(m_i)$ f"ur $i=1,\ldots,k$
benutzt.

Wir haben hierbei nur eine gewisse Wahrscheinlichkeit, da\3 diese so
gewonnene Gleichung wirklich zu einer L"osung f"uhrt, aber da wir mit
geringf"ugigem Aufwand mehrere solche Gleichungen aufstellen k"onnen,
l"a\3t sich die Wahrscheinlichkeit einen Faktor zu finden, leicht
erh"ohen. Damit haben wir also das Problem, $n$ zu faktorisieren, auf
das Problem, $k$ Zahlen $m_1$,\ldots,$m_k$ zu finden, f"ur die
$f(m_1)\cdot\ldots\cdot f(m_k)$ ein Quadrat modulo $n$ ist, reduziert.

Wie gro\3 sollte nun $B$ gew"ahlt werden? Eine geeignete Wahl f"ur $B$ 
ist sehr wichtig f"ur den Algorithmus. Wenn wir $B$ klein
w"ahlen, dann m"ussen wir weniger $m_i$ testen, f"ur die $f(m_i)$ in
kleine Primfaktoren zerf"allt, und das lineare Gleichungsystem kleiner.
Andererseits wird durch diese Wahl das Finden einer Linearkombination des 
Nullvektors schwieriger. Ein geeignetes $B$ stellt also einen Kompromis zwischen
diesen gegenl"aufigen Aspekten dar.

Heuristische "Uberlegungen legen nahe, da\3 $B$ in der Gr"o\3enordunung von
$$ \exp \left(\frac{1}{2}\sqrt{\log n \log\log n} \right) $$
gew"ahlt werden sollte.

Nehmen wir nun an, da\3 wir $B+10$ ganze Zahlen $m_1,\ldots,m_{B+10}$
finden k"onnen, f"ur die
$$ f(m_i) = \prod_{j=1}^{B} p_j^{\alpha_{ji}} $$
"uber der Faktorbasis $\{ p_1,\ldots p_B\}$ f"ur $1 \le i \le B+10$ 
faktorisieren. Die Exponenten $a_{ji}$ sind dabei nat"urliche Zahlen.

F"ur jedes $i=1,\ldots, B+10$ betrachten wir wie beim Dixon-Algorithmus
den Vektor
$$ v(m_i) = (\alpha_{1i},\ldots,\alpha_{B i}) \in\koerper{2}^{B} $$
Die Koeffizienten $\alpha_{ji}$ geben an, ob der Exponenten von $p_j$
in $f(m_i)$ gerade oder ungerade ist. Wir haben also $B+10$
Vektoren der L"ange $B$, unter denen wir eine Linearkombination
des Nullvektors suchen. Es existieren bis zu 10 verschiedene L"osungen, 
f"ur die gilt
$$ f(m_{i_1})\cdot\ldots\cdot f(m_{i_k}) = 
   \prod_{j=1}^B p_j^{\alpha_{j,i_1} + \ldots + \alpha_{j,i_k}} $$
und da
$$ \alpha_{j,i_1} + \ldots + \alpha_{j,i_k} \equiv 0 \pmod{2} $$
gilt, sind alle Exponenten gerade und somit ist
$$ f(m_{i_1}) \cdot\ldots\cdot f(m_{i_k}) $$
ein Quadrat. Damit liefert dieses Vorgehen also 10 potentielle
L"osungen, mit der eine quadratische Kongruenz aufgestellt werden
kann.

Die Methode in der jetztigen Form eignet sich noch nicht zur
Faktorisierung gro\3er Zahlen. Je kleiner die $f(m)$ sind, um so
gr"o\3er ist die Wahrscheinlichkeit, da\3 $f(m)$ "uber einer
vorgegebenen Faktorbasis $\{ p_1,\ldots,p_B\}$ zerf"allt. Daher
sollten wir m"oglichst kleine Werte f"ur $m$ nehmen, um geeignete
$f(m)$ zu finden.

Aber selbst f"ur $m$ aus dem Intervall $1\le m \le n^{\epsilon}$ mit
$\epsilon > 0$ und klein, liegen die Werte von $f(m)$ wegen
$$ f(m) = \left( m + \unten{\sqrt{n}}\right)^2 - n = 
   2m\unten{\sqrt{n}} + m^2 + \left( \unten{\sqrt{n}}^2 - n \right) 
   \approx 2m\unten{\sqrt{n}} + m^2 $$
in der Gr"o\3enordung $C\unten{\sqrt{n}}$ f"ur eine Konstante $C$, 
wenn $\epsilon < 1/4$ gilt.

Mit einem geeigneten Siebproze\3 k"onnen wir auf wesentlich
effektivere Weise geeignete Werte f"ur $m$ bestimmen. Wir
konzentrieren uns dabei nicht auf einzelne Werte f"ur $m$, sondern
betrachten gleich eine ganze Reihe von Werten.

Dazu ben"otigen wir jedoch noch einige Vor"uberlegungen. Zun"achst ist
klar, da\3 f"ur einen Primfaktor $p$ von $f(m_0)$ gilt, da\3 $p$ auch
Primfaktor von $ f(m_0 + k \cdot p) $ f"ur alle $k\in\Z$ ist. Um die
Menge der Werte f"ur $m$ zu bestimmen, f"ur die $p$ ein Faktor von
$f(m)$ ist, gen"ugt es folglich, die quadratischen Kongruenzen
$$ f(m) \equiv 0 \pmod{p} $$
f"ur $m\in\{0,\ldots, p-1\}$ zu l"osen.

Aus der speziellen Form des quadratischen Polynoms~\ref{quadpolynom}
sieht man, da\3 f"ur $p>2$, das nicht selbst Primfaktor von $n$ ist,
die Kongruenz
\begin{equation} 
   f(m) = (m + \unten{\sqrt{n}})^2 - n \equiv 0 \pmod{p} \label{siebkongruenz} 
\end{equation}
genau dann zwei verschiedene L"osungen besitzt, falls $n$ ein Quadrat
modulo $p$ ist, d.h. das Legendresymbol $\leg{n}{p}=1$. Falls $n$ kein
Quadrat modulo $p$ ist, d.h. das Legendresymbol $\leg{n}{p}=-1$,
existiert keine L"osung.  In diesem Fall ist $p$ kein Teiler von
$f(x)$ f"ur alle $x\in\N$. Wir bezeichnen daher mit
$p_1,\ldots,p_B$ nicht wie bei dem Dixon-Algorithmus die ersten $B$
Primzahlen, sondern die ersten $B$ Primzahlen $p$ mit $\leg{n}{p}=1$
inklusive 2. Diese Menge ist unsere {\em Faktorbasis}.

Die quadratische Siebprozedur funktioniert nun wie folgt. Zu jeder
Primzahl aus der Faktorbasis werden die beiden L"osungen 
\begin{eqnarray*}
  a_p & \equiv & n^{(p-1)/2} - \unten{\sqrt{n}} \pmod{p} \mbox{ und} \\
  b_p & \equiv & -n^{(p-1)/2} - \unten{\sqrt{n}} \pmod{p} 
\end{eqnarray*} 
Kongruenz~\ref{siebkongruenz} berechnet.  Wir betrachten nun nicht nur
einzelne zuf"allig gew"ahlte Werte f"ur $m$, sondern wir untersuchen
ein Intervall von $T$ vielen aufeinander folgenden $m$-Werten.

Zur Vorbereitung berechnen wir zun"acht f"ur jedes $m$ den Wert
$\unten{\log_2 f(m)}$, also einfach die Anzahl bits, die f"ur die
Bin"ardarstellung von $f(m)$ ben"otigt werden. Dieser Wert ist f"ur
viele aufeinanderfolgende $m$-Werte gleich und daher effizient zu
berechnen. Die $m$-Werte werden zusammen mit $\unten{\log_2 f(m)}$ in
einer Tabelle gespeichert.

Danach wird beim Sieben f"ur jede Primzahl $p$ aus der Faktorbasis der
Wert $\unten{\log_2 p}$ von dem Wert $\unten{\log_2 f(m)}$ abgezogen,
falls
$$ m \equiv a_p \pmod{p} \mbox{ oder } m \equiv b_p \pmod{p} $$
gilt.

F"ur ein $f(m)$, das h"ohere Potenzen einer Primzahl $p$ aus der
Faktorbasis enth"alt, werden noch zus"atzliche Tests ben"otigt. Das
hei\3t f"ur Primzahlen $p$ mit $p^k<C$ mu\3 f"ur $i<k$
$$ f(m) \equiv 0 \pmod{p^i} $$
getestet und gegebenfalls mehrfach $\unten{\log_2 p}$ abgezogen werden. 

Am Ende des Siebens wurde also von jedem Eintrag $\unten{\log_2 f(m)}$
f"ur jeden Primteiler $p$ von $f(m)$ der Wert $\unten{\log_2 p}$
entspechend oft abgezogen.  Falls also $f(m)$ "uber der Faktorbasis
zerf"allt, ist dieser Eintrag innerhalb der Rechenungenaugigkeit f"ur
die Rundung etwa gleich Null.

Mit diesem Sieb-Verfahren kann also aus einem Intervall von $T$ vielen
aufeinanderfolgenden Werten f"ur $m$ eine Liste von Werten $f(m)$
erstellt werden, die mit hoher Wahrscheinlichkeit "uber der
Faktorbasis zerfallen. Diese Werte werden dann noch einmal im
Einzelnen betrachtet, da w"ahrend des Siebverfahren die Faktoren nicht
protokolliert wurden. Aber jeder dieser Werte ergibt mit hoher 
Wahrscheinlichkeit einen Vektor "uber $\F_2$, so da"s
auf die oben skizzierte Art und Weise "uber die lineare
Abh"angigkeit ein Faktor von $n$ gefunden werden kann.

Eine auf heuristische Annahmen gest"utzte Laufzeitanalyse ergibt f"ur
diesen Algorithmus einen Aufwand von
$$ L(n)^{1+o(1)} = \exp\left( (1+o(1))\sqrt{\log n\log\log n} \right) .$$


   
   
      

      
      
    
    

    

    


\chapter{Das diskrete Logarithmusproblem}

Alle bis heute bekannten Algorithmen liefern keine effizienten
Berechnungsverfahren zur Bestimmung des diskreten Logarithmus eines
Gruppenelements aus einer endlichen zyklischen Gruppe. Damit wird das 
diskrete Logarithmusproblem neben dem Faktorisieren als eines der beiden 
bekannten harten Probleme angesehen, die sich f"ur kryptographische 
Zwecke eignen. Das Verfahren von ElGamal st"utzt sich auf diese Tatsache.

\section{Einf"uhrung}

Es sei $G$ eine endliche zyklische Gruppe und $\alpha$ ein Erzeuger von $G$.
Dann gilt 
$$ G = \left\{\alpha^i\mid 0\le i < \#G\right\}, $$
wobei $\#G$ die Gruppenordnung von $G$ ist. Der diskrete Logarithmus
eines Elements $\beta$ zu der Basis $\alpha$ in $G$ ist eine ganze Zahl $x$ mit
$\alpha^x = \beta$. Wenn $x$ auf das Intervall $0\le x < \#G$ beschr"ankt ist,
dann ist der diskrete Logarithmus von $\beta$ zu der Basis $\alpha$
eindeutig. Wir schreiben $x = \dlog{\alpha}{\beta}$.

Das diskrete Logarithmusproblem besteht darin, einen effizienten
Algorithmus zu finden, der Logarithmen in einer gegebenen Gruppe 
berechnet.

Man kann sofort zwei Algorithmen zur Berechnung der diskreten Logarithmen
in einer endlichen zyklischen Gruppe $G$ angeben. Der erste berechnet
einmal eine Liste der Logarithmen aller Gruppenelemente. Der zweite
berechnet so lange Potenzen von $\alpha$, bis eine "Ubereinstimmung mit
$\beta$ auftritt. Diese trivialen Algorithmen sind wertlos, wenn $\#G$ 
gro"s ist.

\section{Baby-Step Giant-Step Algorithmus}
Der Baby-Step Giant-Algorithmus eignet sich, um Logarithmen in beliebigen
zyklischen Gruppen zu berechnen. Er stellt eine erhebliche Verbesserung
gegen"uber den in der Einf"uhrung beschrieben trivialen Algorithmen, ist
aber praktisch nicht einsetztbar, wenn die Gruppenordnung zu gro"s ist.

Wir wollen den diskreten Logarithmus $\dlog{\alpha}{\beta}$ berechnen.
Es sei $m=\oben{\sqrt{\#G}}$.
\begin{enumerate}
  \item Stelle eine Liste der Paare $(i,\alpha^i)$ f"ur $0\le i < m$
    (offensichtlich gilt $i=\dlog{\alpha}{\alpha^i}$) und sortiere
    diese Liste nach der zweiten Komponente.
  \item Berechne f"ur jedes $j$ mit $0 \le j < m$ $\beta \alpha^{-j m}$ und
    "uberpr"ufe (mit Bin"arsuche), ob dieses Element als zweite 
    Komponente eines Listeneintrags vorkommt.
  \item Falls $\beta\alpha^{-jm} = \alpha^i$ f"ur ein $i$ mit $0 \le i < m$,
    dann gilt $\beta = \alpha^{i+jm}$ und $\dlog{\alpha}{\beta} = i + j m$. 
\end{enumerate}
Der Algorithmus stellt eine Liste mit $\Oklasse{m}$ Eintr"agen und 
ben"otigt $\Oklasse{m \log m}$ Operationen, um die Liste zu sortieren
und dann in der Liste zu suchen. Mit einer Operation ist
eine Gruppenoperation oder ein Vergleich gemeint. Der Name beschreibt in
kennzeichnender Weise die Funktionsweise des Algorithmus.

"Ubung: Wie sieht die Giant-Step Baby-Step Version des Algorithmus?

\section{Index Calculus}

Wir beschreiben generisch die Funktionsweise der Index Calculus Methode.
Es sei $G$ eine endliche zyklische Gruppe der Ordnung $n$, die durch
$\alpha$ erzeugt wird. Wir wollen die Logarithmen der Gruppenelemente zu
dieser Basis berechnen. Es sei $S=\{p_1, p_2, \ldots, p_t\}$ eine
Untermenge von $G$ mit der Eigenschaft, da"s ein "`bedeutender"' Teile
der Gruppenelemente sich als Produkt der Elemente aus $S$ schreiben l"a"st.
Die Menge $S$ wird als die {\em Faktorbasis} bezeichnet.

Im ersten Schritt versuchen wir die Logarithmen aller Elemente aus $S$
zu bestimmen. Dazu w"ahlen wir eine Zufallszahl $a$ und versuchen $\alpha^a$
als Produkt der Elemente aus $S$ zu schreiben:
\begin{equation}
  \alpha^a = \prod_{i=1}^t p_i^{\lambda_i}.\label{produkt}
\end{equation}
Wenn wir so eine Darstellung finden k"onnen, erhalten wir aus \ref{produkt}
eine lineare Kongurenz
\begin{equation}
  a \equiv \sum_{i=1}^t \lambda_i \dlog{\alpha}{p_i} \pmod{n}.\label{relation}
\end{equation}
Nachdem wir eine gen"ugend gro"se Anzahl (gr"o"ser als $t$) an Relationen
der Form~\ref{relation} aufgestellt haben, k"onnen wir erwarten, da"s das 
zugeh"orige lineare
Gleichungssystem eine eindeutige L"osung f"ur die Unbekannten 
$\dlog{\alpha}{p_i}$ mit $1 \le i \le t$ besitzt.

Im zweiten Schritt berechnen wir individuelle Logarithmen in $G$.
Es sei ein $\beta\in G$ gegeben, von dem wir den Logarithmus 
$x=\dlog{\alpha}{\beta}$ bestimmen m"ochten. Wir w"ahlen solange Zufallszahlen
$s$, bis $\alpha^s\beta$ sich als Produkt von Elementen aus $S$ schreiben
l"a"st:
\begin{equation}
  \alpha^s\beta = \prod_{i=1}^t p_i^{b_i}. \label{inS}
\end{equation}
Es gilt dann
$$ \dlog{\alpha}{\beta} \equiv \sum_{i=1}^t b_i\dlog{\alpha}{p_i} - s
   \pmod{n}. $$
Um die Beschreibung der Index Calculus Methode abzuschlie"sen, m"ussen wir
angeben, wie ein geeignetes $S$ zu w"ahlen ist und wie effizient die
Relationen \ref{produkt} und \ref{inS} zu erzeugen sind. Ein geeignetes
$S$ ist klein (das Gleichungssystem im ersten Schritt ist dann nicht zu
gro"s) und gleichzeitig mu"s die Anzahl der Gruppenelemente, die "uber
$S$ faktorisieren, gro"s sein (die erwartete Zeit, um die Relationen
\ref{produkt} und \ref{inS} aufzustellen, ist nicht zu lang). Solche Spezifikationen
sind f"ur einige Endliche K"orper bekannt. Wir beschreiben kurz die
Index Calcus Methode f"ur endliche K"orper. Die Algorithmen machen Gebrauch 
von den unterschiedlichen Darstellungsm"oglichkeiten der K"orperelemente. 
Deswegen wird eine kurze Einf"uhrung in die Theorie der endlichen K"orper
gegeben.

Die erwarteten Laufzeiten der Algorithmen werden durch die Funktion
$$ L[x,\alpha,c] = \Oklasse{\exp((c+o(1))(\log x)^\alpha)
   (\log\log x)^{1-\alpha})}, $$
beschrieben, wobei $x$ die Gruppenordung ist, $0\le\alpha\le 1$ und $c$ Konstanten sind.
Wenn $\alpha = 0$, dann ist $L[x,\alpha,c]$ polynomial in
$\log x$. Bei $\alpha = 1$ ist $L[x,\alpha,c]$ dagegen exponentiell in
$\log x$. Wenn $0 < \alpha < c$, dann ist $L[x,\alpha,c]$ 
{\em subexponentiell} in $\log x$.

\subsection{Endliche K"orper}

  \begin{satz}[Primk"orper]
    F"ur $m\in\N$ sind 
    (i)  $m$ ist eine Primzahl und (ii) $\Z/m\Z$ ist ein K"orper 
    "aquivalent.
  \end{satz}
  Die endlichen K"orper $\koerper{p} = \Z/p\Z$, wobei $p$ eine Primzahl ist,
  hei\3en {\em Primk"orper}.

  \begin{definition}[K"orpererweiterung]
    Unter einer K"orpererweiterung versteht man ein Paar von K"orpern 
    $K \subset L$, wobei $K$ ein Teilk"orper von $L$ ist. Bezeichnung: $L/K$
    oder $K<L$. L hei\3t Erweiterungsk"orper von K.
  \end{definition}
  Es sei $L$ eine K"orpererweiterung von $K$. Man kann die Multipikation
  auf $L$ einschr"anken zu einer Multipikation $K \times L \rightarrow L$  
  und auf diese Weise $L$ als $K$-Vektorraum auffassen. Mit Hilfe dieser
  Beobachtung erh"alt man eine wichtige Strukturaussage "uber endliche
  K"orper. 

  Es sei $\koerper{q}$ ein endlicher K"orper mit q Elementen und 
  $\koerper{q}/\koerper{p}$. Dann ist $\koerper{q}$ ein endlicher 
  $\koerper{p}$-Vektorraum, da $\koerper{q}$ endlich ist. Es sei
  $ \mbox{dim}_{\koerper{p}}(\koerper{q}) = [\koerper{q} : \koerper{p}] = n$
  die Dimension und $\{\omega_1, \ldots, \omega_n\}$ eine Basis von
  $\koerper{q}$ "uber $\koerper{p}$. Jedes Element von $\koerper{q}$ l"a\3t
  sich in eindeutiger Weise als Linerakombination
  $ a_1 \omega_1 + \ldots + a_n \omega_n $
  mit $a_1, \ldots, a_n\in\koerper{p}$ schreiben. Die Anzahl der Elements
  des K"orpers $\koerper{q}$ ist also $p^n$.

  Durch die Vektorrauminterpretation erhalten wir die Strukturaussage, da\3
  die Anzahl der Elemente eines endlichen K"orpers eine Primzahlpotenz
  ist.
  \begin{definition}[Irreduzibel]
    Ein Polynom $f(x)\in K[x], f\neq 0$, wobei $K$ ein K"orper ist, hei\3t 
    {\em irreduzibel}, falls f"ur jede Zerlegung $f(x)=g(x) h(x)$ mit 
    $g(x),h(x)\in K[x]$ 
    gilt $g(x)\in K^*$ oder $h(x)\in K^*$. 
  \end{definition} 

  \begin{satz}[Konstruktion eines K"orpers mit $p^n$ Elementen]
    Es sei $f(x)\in\koerper{p}[x]$ ein irreduzibles Polynom mit $\grad f=n$.
    Dann ist der Restklassenring $\koerper{p}[x]/f(x)$ sogar ein K"orper.
  \end{satz}
  \beweis Die K"orpereigenschaft kann leicht dadurch verdeutlicht werden, 
  indem der erweiterte ggT-Algorithmus ELBA betrachtet wird. Der ELBA
  stellt den ggT zweier Elemente eines euklidischen Rings als eine
  Linearkombination von ihnen dar. 
  Wir fordern vor"ubergehend, da\3 $f(x)$ normiert ist, d.h. der 
  Leitkoeffizient ist gleich $1$. Da $f(x)$ ein normiertes irreduzibles 
  Polynom vom Grad $n$ ist, ist der $\ggt (f(x),g(x))=1$ f"ur alle 
  Polynome $g(x)\in\koerper{p}[x]$ mit $\grad(g(x)) < n$. 
  Der ELBA liefert die Darstellung 
  $$ \ggt (f(x),g(x)) = 1 = r(x)f(x)+s(x)g(x) $$
  mit $r(x), s(x)\in\koerper{p}[x]$. Wird die Berechnung modulo $f(x)$
  ausgef"uhrt, so ergibt sich direkt
  $$ 1\equiv s(x)g(x) \pmod{f(x)} $$
  und $s(x)$ ist somit das inverse Element von $g(x)$ modulo $f(x)$.
  Die Eigenschaft, da\3 $f(x)$ normiert ist, ist unwesentlich, da wir sonst
  mit dem inversen Element des Leitkoeffizients durchmultiplizieren k"onnen.
  Alle Element au\3er der Restklasse der Null haben bez"uglich der 
  Multiplikation ein Inverses. Der Restklassenring ist also ein K"orper 
  mit $p^n$ Elementen, da soviele Restklassen modulo $f(x)$ existieren. 
  \hfill $\Box$

  \begin{satz}[Verfahren von Kronecker]
    Es sei $K$ ein K"orper und $f(x)\in K[x]$ ein nichtkonstantes Polynom.
    Dann existiert ein Erweiterungsk"orper $K$ ($K<L$), so da\3 $f(x)$ eine
    Nullstelle in $L$ besitzt. Ist $f(x)$ irreduzibel, so kann man 
    $$ L:= K[x]/f(x) $$
    setzen.
  \end{satz}
  \beweis Es sei ohne Einschr"ankung $f(x)$ irreduzibel, ansonsten 
  faktorisieren wir $f(x)$ und ersetzen $f(x)$ durch einen 
  seiner irreduziblen Faktoren. Wir bilden die Komposition
  $$ \phi: K \stackrel{\iota}{\hookrightarrow} K[x] 
             \stackrel{\pi}{\rightarrow} K[x]/f(x) =:L. $$
  Der resultierende Homomorphismus $\phi = \pi \circ \iota$ ist injektiv.
  Wir k"onnen $L$ als Erweiterungsk"orper von $K$ auffassen, indem wir
  $K$ mit seinem Bild $\phi (K)$ unter $\phi: K \rightarrow L$
  identifizieren. Wir setzen $\bar{x} := \pi(x)$. Mit 
  $ f(X) = \sum_{i=0}^{n} c_i X^i$ gilt dann
  \begin{eqnarray*}
    f(\bar{x}) & = & \sum_{i=0}^{n} c_i \bar{x}^i \\
               & = & \sum_{i=0}^{n} \pi(x)^i \\
               & = & \pi\left( \sum_{i=o}^{n} c_i x^i\right) \\
               & = & \pi\left(f(x)\right) \\
               & = & \bar{0},
  \end{eqnarray*}
  d.h. $\bar{x}$ ist eine Nullstelle von $f(x)\in K[x]$ in $L$. 
  \hfill $\Box$

  Bei dem Verfahren von Kronecker ensteht $L$ aus $K$ durch
  {\em Adjunktiion einer Nullstelle $\bar{x}$ von $f(x)$}. Die Nullstelle
  $\bar{x}$ wird sozusagen erzwungen, indem wir ausgehend von dem 
  Polynomring $K[x]$ den Restklassenring $L:=K[x]/f(x)$ bilden. Wir
  k"onnen einen Linerfaktor von $f$ "uber $L$ abspalten und k"onnen das
  Verfahren wiederholen. Nach endlich vielen Schritten gelangt man zu einem
  Erweiterungsk"orper, "uber dem $f$ vollst"andig in Linearfaktoren
  zerf"allt.
  \begin{lemma}[Mehrfache Nullstellen]
    Ein normiertes Polynom $f(x)\in K[x]$ hat genau dann keine mehrfachen
    Nullstellen, wenn $\ggt (f(x),f'(x))=1$.
  \end{lemma}

  \begin{definition}[Zerf"allungsk"orper]
    Ein Zerf"allungsk"orper eines Polynoms $f\in K[x]$ ist ein
    Erweiterungsk"orper $L$ des K"orpers $K$, in dem $f$ vollst"andig
    in Linearfaktoren zerf"allt.
  \end{definition}  
  \begin{satz}[Isomorphie der minimalen Zerf"allungsk"orper]
    Es seien $L_1$ und $L_2$ zwei minimale Zerf"allungsk"orper eines
    Polynoms $f\in K[x]$. Dann sind $L_1$ und $L_2$ isomorph.
  \end{satz}
  \begin{lemma}[Gruppenordung]
    Es sei G eine Gruppe. Dann gilt f"ur jedes Gruppenelement 
    $ g^{\betrag{G}} = 1. $
  \end{lemma}
  \beweis mit dem Satz von Lagrange. \hfill $\Box$

  Die multiplikative Gruppe $\koerper{q}^{*}$ eines endlichen K"orpers 
  $\koerper{q}$ hat die Ordnung $q-1$, da alle Elemente au\3er des
  Nullelements invertierbar sind. Daher erf"ullt jedes
  $\alpha\in\koerper{q}$ die Gleichung $X^{q-1}=1$. Wenn man noch das
  Nullelement ber"ucksichtigt, erf"ullt jedes Element 
  $\alpha\in\koerper{q}$ die Gleichung
  $$ f(x) = x\left( x^{q-1} - 1\right) = x^q - x = 0. $$
  Das Polynom $f(x)$ hat $q$ verschiedne Nullstellen in $\koerper{q}$,
  n"amlich die Elemente von $\koerper{q}$. Das Polynom $f(x)$ zerf"allt
  in Linearfaktoren "uber $\koerper{q}$
  $$ f(x) = x^q-x = \prod_{\alpha\in\koerper{q}} \left( x - \alpha \right). $$
  Damit erhalten wir eine zweite Strukturaussage "uber endliche K"orper.
  Alle K"orper $K$ mit $q=p^n$ Elementen sind isomorph, da sie
  minimale Zerf"allungsk"orper des Polynoms $f(x) = x^q - x$ sind.
  
  Diese Aussagen lassen sich zu folgendem Theorem zusammenfassen.
  \begin{theorem}[Strukur endlicher K"orper]
    F"ur jede Primzahl $p$ und jedes $n\in\N$ existiert ein K"orper 
    $\koerper{p^n}$ mit $p^n$ Elementen. Er ist der minimale 
    Zerf"allungsk"orper des Polynoms $x^{p^n}-x$ und seine Elemente sind die 
    Nullstellen dieses Polynoms. 

    Jeder endliche K"orper ist isomorph zu genau einem K"orper
    $\koerper{p^n}$.
  \end{theorem}
  Der folgende Satz ist f"ur die Defintion des diskreten Logarithmus in 
  endlichen K"orpern wichtig.
  \begin{satz}[Struktur der multipikativen Gruppe $\koerper{q}^*$] 
    Die multiplikative Gruppe eines endlichen K"orpers ist zyktisch. Ein 
    Element $\omega\in\koerper{q}^*$ mit 
    $\left<\omega\right>=\koerper{q}^{*}$ hei\3t {\em primitives Element}.
  \end{satz}
%%
%% Index-Calculus
%%
\subsection{Index-Calculus f"ur $\koerper{p}^*$}

Wir stellen die Elemente von $\koerper{p}^*$ als die Menge der Zahlen
$\{1,2,\ldots, p-1\}$. Die Multiplikation wird modulo $p$ ausgef"uhrt.
Der Erzeuger von $\koerper{p}^*$ sei $\alpha$. Es sei $m$ eine Zahl, die
eine Funktion von $p$ ist. Die Faktorbasis $S$ ist die Menge aller
Primzahlen kleiner $m$. 

Der erste und zweite Schritt werden wie im generischen Fall ausgef"uhrt.

Wenn wir die Wahrscheinlichkeit, da"s eine Zufallszahl kleiner $p$ in
Primfaktoren kleiner $m$ zerf"allt, betrachten, k"onnen wir $m$ optimal
w"ahlen, um die Vorberechnungsphase (erster Schritt) und die Berechnung
eines beliebigen Logarithmus (zweiter Schritt) zu beschleunigen. Das f"uhrt
zu einer heuristischen aber subexponentiellen Implementierung des 
Index Calculus Algorithmus mit einer erwarteten Laufzeit von
$L[p,1/2,c]$. Die Laufzeit ist nicht rigoros, da angenommen wird, da"s
das Gleichungssystem im ersten Schritt vollen Rang hat.

\subsection{Index-Calculus f"ur $\koerper{p^n}$}

Die Elemente des endlichen K"orpers $\koerper{p^n}$ werden als Polynome vom 
Grad kleiner $n$ aus $\koerper{p}[x]$ aufgefa\3t. Die Addition ist
gew"ohnliche Polynomaddition und die Multiplikation wird modulo eines
festgelegten irreduziblen Polynoms $f(x)\in\koerper{p}[x]$ ausgef"uhrt.
Wir benutzen also die Isomorphie $\koerper{p^n} \cong \koerper{p}[x]/(f(x))$.

\begin{definition}
  Ein Polynom $w(x)\in\koerper{p}[x]$ hei\3t {\em glatt} bez"uglich
  der Schranke $b\in\N$, genau dann, wenn es in Polynome aus 
  $\koerper{p}[x]$ vom Grad kleiner $b$ faktorisiert.
\end{definition}

Mit der Polynomdarstellung des K"orpers $\koerper{p^n}$ ist es offensichtlich,
die Faktorbasis $S$ als die Menge aller irreduziblen $b$-glatten Polynome aus 
$\koerper{p}[x]$ zu w"ahlen. Im ersten Schritt werden die Logarithmen aller Elemente
der Faktorbasis bestimmt. Es kann gezeigt werden, da"s f"ur kleine $p$ und eine
geeignet Wahl f"ur $b$, der Index Calculus Algorithmus eine erwartete
Laufzeit von $L[p^n,1/2,c]$ hat.

{\bf Algorithmus: }

Es sei $\alpha(x)$ ein Erzeuger von $\koerper{p^n}^*$. Die Faktorbasis
$S$ wird als die Menge aller irreduziblen Polynome $p(x)\in\koerper{p}[x]$ 
mit $\deg(p(x)) \le b$ gew"ahlt. Die Schranke $b$ geben wir sp"ater an.
\begin{enumerate}
  \item Vorberechnungsphase: es werden die diskreten Logarithmen aller 
    Elemente aus der Faktorbasis $S$ berechnet.
    \begin{enumerate}
      \item W"ahle $a_i\in\{1,\ldots,p^n-1\}$ zuf"allig aus und
        berechne $y_i(x) := \alpha(x)^{a_i}$.  
        Wir k"onnen effizient mit dem ggT-Algoritmus "uberpr"ufen, ob
        die Faktoren $p_i(x)\in S$ $y_i(x)$ teilen. Es werden alle
        maximalen Potenzen von $p_i(x)\in S$ von $y_i(x)$ abgespalten.
        Falls wir danach ein Polynom ungleich $c_i\in\koerper{p}$
        erhalten, wissen wir, da"s $y_i(x)$ nicht $b$-glatt ist.
        Falls $y_i(x)$ $b$-glatt ist, so speichern wir die Faktorisierung
        $$ y_i(x) = c_i \cdot \prod_{p_j(x)\in S} p_j(x)^{e_{ij}} $$ mit
        $c_i\in\koerper{p}$. F"ur die diskreten Logarithmen $a_i$
        gilt dann
        \begin{equation} a_i = \dlog{\alpha(x)}{y_i(x)} =
            \dlog{\alpha(x)}{c_i} + \sum_{p_j(x)\in S} e_{ij} \cdot
            \dlog{\alpha(x)}{p_j(x)} \label{loggl} 
        \end{equation}
        Wir wiederholen dies so lange, bis wir mehr als $\#S$ Gleichungen
        der Form~\ref{loggl} haben. 
      \item 
        L"ose dieses lineare Gleichungssystem mit den unbekannten
        $\dlog{\alpha(x)}{p_j(x)}$ f"ur $p_j(x)\in S$.
    \end{enumerate}

  \item Berechnen eines konkreten Logarithmus von $y(x)=\alpha(x)^{s}$. 

        W"ahle $s$ so lange zuf"allig aus, bis $\alpha(x)^s y(x)$
        $b$-glatt ist. Das Polynom $\alpha(x)^s y(x)$ kann dann wie folgt
        $$ \alpha(x)^s y(x) = c \cdot \prod_{p_j(x)\in S} p_j(x)^{e_j} $$
        mit $c\in\koerper{p}$ faktorisiert werden. Daraus kann
        bereits der diskrete Logarithmus $\dlog{\alpha(x)}{y(x)}$ berechnet
        werden, denn es gilt
        $$ \dlog{\alpha(x)}{y(x)} = \dlog{\alpha(x)}{c} - s +
        \sum_{p_j(x)\in S} e_{j} \cdot \dlog{\alpha(x)}{p_j(x)}. $$ 
        Die diskreten Logarithmen $\dlog{\alpha(x)}{p_j(x)}$ f"ur $p_j(x)\in
        S$ wurden im Schritt 1.(a) bestimmt und $\dlog{\alpha(x)}{c}$
        ist im (deutlich kleineren Grundk"orper) einfach zu berechnen.
\end{enumerate}

Die Schranke $b$ sollte so gew"ahlt werden, da\3 f"ur die Menge $S$
$$ \# S \approx \exp\left( c \cdot \sqrt{\log q\log\log q} \right) $$
gilt, wobei $q=p^n$.
%%
%% Coppersmith-Algorithmus
%%

\section{Coppersmith-Algorithmus f"ur $\koerper{2^n}$}

Der Coppersmith-Algorithmus basiert auf der speziellen Darstellung des
K"orpers $\koerper{2^n} \cong \koerper{2}[x]/(x^n + f_1(x))$ mit
$\grad(f_1(x)) \le \log_2(n)$. Ein solches Polynom $f_1(x)$ existiert
f"ur gen"ugend gro\3e $n$.

%%
%% Algorithmus
%%

\subsubsection*{ 1. Erzeugung der Datenbasis}

Es werden die Logarithmen der Elemente aus der Menge $S$ berechnet,
wobei $S$ die irreduziblen Polynome vom Grad kleiner $s$
enth"alt. Dabei wird $s \approx (n\log^2(n))^{1/3}$ und das primitive
Element $\alpha(x)\in S$ gew"ahlt. Beim Index-Calculus-Algorithmus wird
ein beliebiges Polynom (K"orperelement) faktorisiert. Coppersmith
verwendet dagegen zwei Polynome vom Grad $n^{2/3}$. Da diese einen
kleineren Grad haben, steigt die Wahrscheinlichkeit, da\3 sie "uber
$S$ zerfallen.

Es sei $k\in\N$ so gew"ahlt, da\3 $2^k \le (n/\log_2(n))^{1/3}$ und
setze $h:=\oben{\frac{n}{2^k}}$.  W"ahle $u_1(x)$ und $u_2(x)$
teilerfremd vom Grad kleiner $s$ und berechne
\begin{equation} w_1(x) = u_1(x) \cdot x^h + u_2(x) 
\end{equation}
und
\begin{equation} w_2(x) \equiv w_1(x)^{2^k} \pmod{f(x)}. \label{w2} 
\end{equation}
Es gilt dann 
\begin{eqnarray*} 
  w_2(x) & \equiv & u_1(x)^{2^k} \cdot x^{h \cdot 2^k} 
                                      + u_2(x)^{2^k} \pmod{f(x)} \\
         & \equiv & u_1(x)^{2^k} \cdot x^{h \cdot 2^k - n} \cdot f_1(x) 
                                      + u_2(x)^{2^k} \pmod{f(x)},
\end{eqnarray*}
da $x^n \equiv f_1(x) \pmod{f(x)}$ und das Potenzieren mit $2^k$ in
K"orpern der Charakteristik $2$ ein Homomorphismus ist.  Zerfallen
beide Polynome $w_1(x)$ und $w_2(x)$ in irreduzible Polynome aus $S$,
so erhalten wir, wenn wir den diskreten Logarithmus der
Gleichung~\ref{w2} nehmen, die Beziehung
$$ 
\begin{array}{rcll}
  \dlog{\alpha(x)}{w_2(x)} & \equiv & 2^k \cdot \dlog{\alpha(x)}{w_1(x)} 
                  & \pmod{ 2^n - 1} \\
  \Leftrightarrow \dlog{\alpha(x)}{w_2(x)} - 2^k \cdot \dlog{\alpha(x)}{w_1(x)} 
                  & \equiv & 0 & \pmod{ 2^n - 1}.
\end{array}
$$
und damit eine lineare Gleichung f"ur die Logarithmen der Elemente aus $S$.
      
Der Aufwand f"ur die Aufstellung der Datenbasis ist $O(\exp( (c+o(1))
(n\log^2(n))^{1/3} ))$ mit $c \approx \mbox{1,4}$. Diese Laufzeitabsch"atzung
st"utzt sich auf heuristische Annahmen.


\subsubsection*{2. Berechnung des Logarithmus}

Es sei der Logarithmus des Elements $g(x)$ bez"uglich der Basis
$\alpha(x)$ gesucht. Falls $g(x)$ einen Grad $\deg(g(x))\le s$ hat,
dann kann der Logarithmus aus der Datenbasis entnommen werden.

Es sei nun $\deg(g(x))>s$. Wir gehen zun"achst wie beim
Index-Calculus-Algorithmus vor. Es wird $r\in\{1,\ldots,n-1\}$
zuf"allig ausgew"ahlt und
$$ h(x) = g(x) \cdot \alpha^r(x) $$ 
berechnet. Wir faktorisieren $h(x)$ und "uberpr"ufen, ob es glatt
bez"uglich der Schranke $\tilde{s} := (n^2\log_2(x))^{1/3}$ ist. Falls
ja, dann ergibt sich der Logarithmus von $g(x)$ zu
$$ \dlog{\alpha(x)}{g(x)} = \sum_{i=1}^{k}
   \log_{\alpha(x)}{p_i(x)}-r. 
$$ 
Die Polynome $p_i(x)$ sind vom Grad kleiner oder gleich $\tilde{s}$. Die
Logarithmen $\dlog{\alpha(x)}{p_i(x)}$ der Polynome $p_i(x)$ mit
$\deg{p_i(x)}\le s$ stehen in der Datenbasis. Die noch fehlenden
Logarithmen der Polynome $p_i(x)$ mit $s < \deg(p_i(x)) \le \tilde{s}$
werden noch berechnet. Die Berechnung wird vereinfacht, indem man
sie durch Polynome kleineren Grades ersetzen. Zur Bestimmung des
Logarithmus von $p_i(x)$ w"ahlen wir $k$, so da\3 $2^k \le \sqrt{n/b}$
mit $b\approx (n^2\log n)^{1/3} + \deg(p_i(x))/2$ und setzen
$h:=\oben{n/2^k}$.

Anschlie\3end w\"ahlt man zwei zuf"allige teilerfremde Polynome $u_1(x)$
und $u_2(x)$ vom Grad kleiner gleich $\tilde{s}$ und berechnen die
Polynome
$$ w_1(x) := u_1(x)\cdot x^h + u_2(x) $$
und
$$ w_2(x) := w_1(x)^{2^k}, $$
so da\3 $p_i(x) \mid w_1(x)$.
Zerfallen $\frac{w_1(x)}{p_i(x)}$ und $w_2(x)$ in irreduzible Polynome
kleineren Grades,
$$ w_1(x) = p_i(x) \prod_{j=1}^k s_i(x) $$
und
$$ w_2(x) = \prod_{j=1}^l t_j(x), $$
so erh"alt man
$$
\begin{array}{cccc}
   \ds \dlog{\alpha(x)}{w_2(x)}  
      & \equiv 
      & \ds 2^k\dlog{\alpha(x)}{w_1(x)} \equiv \sum_{j=1}^l \dlog{\alpha(x)}{t_j(x)}
      & \pmod{2^n - 1} \\
   \ds 
      & \equiv 
      & \ds 2^k(\dlog{\alpha(x)}{p_i(x)} + \sum_{j=1}^k \dlog{\alpha(x)}{s_j(x))} 
      & \pmod{2^n - 1} \\
   \ds \Rightarrow  \dlog{\alpha(x)}{p_i(x)} 
      & \equiv 
      & \ds 2^{-k} \sum_{j=1}^l \dlog{\alpha(x)}{t_i(x)} - \sum_{j=1}^k \dlog{\alpha(x)}{s_j(x)} 
      & \pmod{2^n - 1}
\end{array}
$$
Diese Gleichnung reduziert die Berechnung der Logarithmen
$\dlog{\alpha(x)}{p_i(x)}$ auf die Berechnung der
$\dlog{\alpha(x)}{s_j(x)}$ und $\dlog{\alpha(x)}{t_j(x)}$. Diese
k"onnen entweder aus der Datenbasis entnommen werden, oder werden
durch rekursive Anwendung des obigen Schrittes berechnet.

Der Aufwand betr"agt $O(\exp(c+o(1)) (n\log^2 n)^{1/3}))$ mit $c
\approx \mbox{1,1}$.
   


      
    
    
        
        
   
   

   

\chapter{Authentifikation}

Wenn sich Alice in einen Hostrechner einloggt, wie kann er
Rechner wissen, da\3 es wirklich Alice ist? Wie kann er
erkennen, da"s Eve versucht sich als Alice auszugeben?
"Ublicherweise wird dieses Problem mit Hilfe eines Passworts
gel"ost. Alice gibt ihr Passwort ein und der Rechner "uberpr"uft, ob es
korrekt ist. Alice und der Rechner haben ein gemeinsames Geheimnis, das
der er jedesmal von Alice verlangt, wenn sie sich einloggen will.

%%
%% Authentifikation mit Einwegfunktionen
%%

\section{Authentifikation mit Einwegfunktionen}

Roger Needham und Mike Guy bemerkten, da\3 der Host gar nicht die
Passw"orter zu kennen braucht. Der Host mu\3 nur die richtigen
Passw"orter von den falschen unterscheiden k"onnen. Dies kann einfach
mit Einwegfunktionen realisiert werden. Der Host speichert nur die
Werte $f(Passwort)$, wobei $f$ eine Einwegfunktion ist.
\begin{enumerate}
  \item Alice sendet ihr Passwort dem Host.
  \item Der Host wendet eine Einwegfunktion auf das Passwort an.
  \item Der Host vergleicht das Ergebnis mit dem zuvor gespeicherten Wert.
\end{enumerate}
Da jetzt der Host nicht die Passw"orter speichert, sondern nur die
Werte der mit einer Einwegfunktion verschl"usselten Passw"orter, ist
diese Liste f"ur einen Angreifer nicht direkt verwendbar, da er daraus die
Passw"orter nicht mehr so leicht rekonstruieren kann.

Diese Authentifikation ist nicht sonderlich sicher, da ein {\em
W"orterbuchangriff} unternommen werden kann.  Eine Datei, die mit
einer Einwegfunktion verschl"usselten Passw"orter enth"alt, stellt
immer noch einen Schwachpunkt dar. Mallory erstellt eine Liste der am
h"aufigsten benutzten Passw"orter, verschl"usselt jedes Passwort mit
der Einwegfunktion und speichert die Ergebnisse. Mallory schafft es die
Datei mit den verschl"usselten Passw"ortern, die sich auf dem
Hostrechner befindet, zu stehlen und sucht nach "Ubereinstimmungen in
den beiden Dateien. Dieser Angriff ist erstaunlich erfolgreich, was
nicht an dem Verfahren mit der Einwegfunktion liegt, sondern ist den
unachtsamen Benutzern anzulasten, da sie einfach erratbare Passw"orter
w"ahlen.

%%
%% Authentifikation mit Public Key
%%

\section{Authentifikation mit Hilfe von Public Key}

Die Protokolle zur Authentifikation mit Hilfe von Passw"ortern haben
eine weitere enorme Sicherheitsl"ucke. Wenn Alice das Passwort zu
ihrem Host sendet, kann Eve ihr Passwort lesen, wenn sie Zugang zu dem
Kommunikationsweg oder dem Prozessorspeicher des Hosts hat.

Mit Hilfe der Public-Key Kryptographie kann dieses Problem behoben
werden. Der Host hat eine Datei mit den "offentlichen Schl"usseln
aller Benutzer und alle Benutzer haben ihren "offentlichen
Schl"ussel. Ein einfaches Protokoll kann wie folgt aussehen:
\begin{enumerate}
  \item Der Host sendet Alice eine Zufallszahl $r$.
  \item Alice verschl"usselt $r$ mit ihrem geheimen Schl"ussel und
        schickt das Chiffrat an den Host.
  \item Der Host entschl"usselt das Chiffrat mit dem "offentlichen
        Sch"ussel von Alice und vergleicht es mit $r$.
\end{enumerate}
Ein Angreifer kann sich jetzt nicht dem Host gegen"uber als Alice
identifizieren, da er keinen Zugang zu Alices geheimen Schl"ussel
hat. Alice sendet niemals ihren geheimen Schl"ussel "uber den
Kommunikationsweg zu dem Host. Eve, die die Interaktion zwischen Alice
und dem Host abh"ort, kann keine Information gewinnen, die es ihr
erm"oglichen w"urde, Alices geheimen Schl"ussel zu berechnen oder sich
als Alice auszugeben, ohne da\3 es der Host bemerkt. Weder der Host
noch der Kommunikationsweg m"ussen sicher sein.

Es ist f"ur Alice nicht geschickt, eine beliebige Zufallszahl $r$, die
sie zugeschickt bekommt, zu verschl"usseln. Eve k"onnte dadurch
die an Alice geschickten Nachrichten entschl"usseln. Ein verbessertes 
Protokoll sieht wie folgt aus:
\begin{enumerate}
  \item Alice berechnet etwas, wobei sie eine eigene Zufallszahl und
        ihren geheimen Schl"ussel verwendet, und schickt das Ergebnis
        an den Host.
  \item Der Host sendet Alice eine andere Zufallszahl.
  \item Alice verwendet eine vorgegebene Funktion, um einen
        Funktionswert, der von den beiden Zufallszahlen und ihrem
        geheimen Schl"ussel abh"angt, zu berechnen und schickt das
        Ergebnis an den Host.
  \item Der Host "uberpr"uft dies.
\end{enumerate}
Falls Alice genauso wenig dem Host vertraut, wie der Host Alice traut,
fordet sie ihn auf, sich ihr gegen"uber zu authentifizieren.

%%
%% Authentifikation mit Schl"usselaustausch
%%

\section{Authentifikation mit Schl"usselaustausch}

Diese Protokolle verbinden die Authentifikation mit einem
Schl"usselaustausch, um ein generelles Problem zu l"osen: Alice und
Bob sind an zwei verschiedenen Enden eines Netzwerks und wollen sicher
Daten "ubertragen. Wie k"onnen Alice und Bob einen geheimen Schl"ussel
austauschen und dabei sicher sein, da\3 keiner von ihnen in
Wirklichkeit mit einem Angreifer Mallory spricht? Die meisten dieser
Protokolle basieren auf der Tatsache, da\3 eine Schl"usselzentrale mit
jedem Teilnehmer einen gemeinsamen geheimen Schl"ussel besitzt und
da\3 alle Schl"ussel am richtigen Ort sind, bevor das Protokoll
beginnt.

%%
%% Wide-mouth frog
%%

\subsubsection*{Wide-mouth frog}

Das wide-mouth frog Protokoll ist vielleicht das einfachste Protokoll,
das ein symmetrisches Verfahren und eine vertrauensw"urdige
Schl"usselzentrale benutzt.  Alice und Bob besitzen jeweils einen
geheimen Schl"ussel $A$ bzw. $B$ mit der Schl"usselzentrale. Diese
Schl"ussel werden nur f"ur den Schl"usselaustausch und nicht zur
Verschl"usselung von eigentlichen Nachrichten verwendet. Mit Hilfe von
nur zwei Nachrichten "ubermittelt Alice an Bob einen geheimen
Sitzungsschl"ussel $K$.  Mit $T_A$ und $T_B$ werden Zeitstempel
bezeichnet. Die Verschl"usselungsfunktionen $E_A$ und $E_B$ sind mit
den geheimen Schl"usseln parametrisiert, die Alice und Bob mit der
Schl"usselzentrale teilen.
\begin{enumerate}
  \item Alice erzeugt das Tupel $(Alice,E_A(T_A,Bob,K))$ und schickt es an
        die Schl"us\-sel\-zen\-trale.
  \item Die Schl"usselzentrale entschl"usselt die Nachricht von
        Alice. Sie erzeugt einen neuen Zeitstempel $T_B$ und sendet das
        verschl"usselte Tupel $E_B(T_B,Alice,K)$ an Bob.
\end{enumerate}
Die wichtigste Annahme bei diesem Protokoll ist, da\3 Alice kompetent
ist, gute (sichere) Sitzungschl"ussel zu erzeugen. Man mu\3 beachten,
da\3 es schwierig ist, Zufallszahlen zu erzeugen. Alice ist
m"oglicherweise nicht dazu in der Lage.

%%
%% Needham-Schroeder
%%

\subsubsection*{Needham-Schroeder}

Das Protokoll von Needham und Schroeder benutzt ein symmetrisches Verfahren
und eine Schl"usselzentrale.
\begin{enumerate}
  \item Alice erzeugt eine Zufallszahl $R_A$ und schickt das Tupel
        $(Alice,Bob,R_A)$ an die Schl"us\-sel\-zen\-trale.
  \item Die Schl"usselzentrale erzeugt einen zuf"alligen
        Sitzungsschl"ussel $K$ und schickt das Tupel
        $E_A(R_A,Bob,K,E_B(K,Alice))$ an Alice.
  \item Alice entschl"usselt und erh"alt den Sitzungsschl"ussel
        $K$. Sie "uberpr"uft die Zufallszahl $R_A$ und schickt das
        Tupel $E_B(K,Alice)$ an Bob.
  \item Bob entschl"usselt und erh"alt den Sitzungsschl"ussel $K$. Er
        erzeugt eine Zufallszahl $R_B$ und schickt das Tupel
        $E_K(R_B)$ an Alice.
  \item Alice entschl"usselt und erh"alt die Zufallszahl $R_B$. Sie
        berechnet $R_B-1$ und schickt das Tupel $E_K(R_B-1)$ an Bob.
  \item Bob pr"uft, ob $R_B-1$ geschickt wurde.
\end{enumerate}

%%
%% Kerberos
%%

\subsubsection*{Kerberos}

Das Kerberos-Protokoll ist eine Variante des
Needham-Schroeder-Protokolls. Die Zufallszahl $R_A$ wird durch einen
Zeitstempel ersetzt.
\begin{enumerate}
  \item Alice sendet das Tupel $(Alice,Bob)$ an die
        Schl"usselzentrale.
  \item Die Schl"usselzentrale schickt das Tupel
        $(E_A(T_{SZ},L,K,Bob),E_B(T_{SZ},L,K,Alice))$ an Alice. Die
        Nachricht wird mit einer G"ultigkeitsdauer $L$ und einem
        Zeitstempel $T_{SZ}$ versehen.
  \item Alice schick das Tupel
        $(E_K(Alice,T_A),E_B(T_{SZ},L,K,Alice))$ an Bob.
  \item Bob sendet das Tupel $E_K(T_A+1)$ an Alice.
\end{enumerate}

%%
%% Neuman-Stubblebine
%%

\subsubsection*{Neuman-Stubblebine}

\begin{enumerate}
  \item Alice erzeugt eine Zufallszahl $R_A$ und sendet das Tupel
        $(Alice,R_A)$ an Bob.
  \item Bob erzeugt eine Zufallszahl $R_B$ und sendet das Tupel
        $(Bob,R_B,E_B(Alice,R_A,T_B))$ mit einem Zeitstempel $T_B$
        versehen an die Schl"usselzentrale.
  \item Die Schl"usselzentrale erzeugt einen zuf"alligen
        Sitzungsschl"ussel $K$ und sendet das Tupel
        $E_A(Bob,R_A,K,T_B),E_B(Bob,K,T_B),R_B$ an Alice.
  \item Alice entschl"usselt und "uberpr"uft, ob $R_A$ mit dem Wert im
        1-ten Schritt "ubereinstimmt. Alice sendet
        $E_B(Alice,K,T_B),E_K(R_B)$ an Bob.
  \item Bob entschl"usselt und erh"alt den Sitzungsschl"ussel $K$. Er
        "uberpr"uft, ob $T_B$ und $R_B$ mit den Werten im 2-ten
        Schritt "ubereinstimmen.
\end{enumerate}

%%
%% Digital Cash
%%

\section{Digital Cash}

Ein ideales System f"ur "`Digital Cash"' soll folgende sechs
Eigenschaften erf"ullen:
\begin{itemize}
  \item {\em Unabh"angigkeit}: das digitale Bargeld kann "uber Netze
        verschickt und gespeichert werden.
  \item {\em Sicherheit}: das digitale Bargeld kann nicht kopiert und
        mehrfach benutzt werden.
  \item {\em Nichzur"uckverfolgbarkeit}: Die Privatsph"are des
        Benutzers ist gewahrt. Keiner kann feststellen, welche
        Transaktionen vom Benutzer vorgenommen wurden.
  \item {\em Off-line Bezahlung}: Bei der Bezahlung mit dem digitalen
        Bargeld ist keine Kommunikation mit der Bank notwendig. Die
        Bezahlung erfolgt off-line.
  \item {\em "Ubertragbarkeit}: Das digitale Bargeld kann an andere
        Benutzer "ubertragen werden.
  \item {\em Teilbarkeit}: Das digitale Bargeld kann in kleinere
        Betr"age aufgeteilt werden.
\end{itemize}
Das erste vorgeschlagene Protokoll, das alle diese Punkte erf"ullte,
ben"otigt 200 MByte Datentransfer f"ur eine Buchung. Mitterweile
existieren Protokolle, die nur 20 KByte Datentransfer brauchen und in
wenigen Sekunden arbeiten.

Im folgenden pr"asentieren wir ein Protokoll, das die ersten vier
Eigenschaft erf"ullt. Wir werden das Protokoll schrittweise ausbauen,
um Sicherheitsl"ucken zu beseitigen.

\subsection{Erste Protokolle}

\subsubsection*{1. Protokoll}
\begin{enumerate} 
  \item Alice stellt 100 anonyme Bankschecks "uber jeweils 1000 DM aus.
  \item Alice steckt jeden Bankscheck zusammen mit Durchschlagspapier
        in einen Umschlag und tr"agt dann alle Umschl"age zur Bank.
  \item Die Bank "offnet 99 Umsch"age und pr"uft, ob der richtige
        Betrag (1000 DM) auf jedem Scheck von Alice eingetragen wurde.
  \item Die Bank unterzeichnet den letzten unge"offneten Umschlag. Die
        Unterschrift der Bank wird durch das Durchschlagspapier auf
        den Scheck "ubertragen.  Die Bank gibt den unge"offneten
        Umschag an Alice zur"uck und belastet ihr Konto mit 1000 DM.
  \item Alice "offnet den Umschag und verwendet den Scheck bei einem
        H"andler.
  \item Der H"andler "uberpr"uft die Unterschrift der Bank und
        akzeptiert den Scheck.
  \item Der H"andler bringt den Scheck zur Bank.
  \item Die Bank "uberpr"uft ihre Unterschrift und zahlt 1000 DM dem
        H"andler aus.
\end{enumerate}
Bei diesem Protokoll kann die Bank nicht feststellen, von wem der
Scheck kommt. Die Wahrscheinlichkeit, da\3 ein Betrug von der Bank
unentdeckt bleibt, betr"agt durch die 1-aus-100-Wahl lediglich
$1/100$. Durch die Einf"uhrung einer Geldstrafe f"ur Betrug lohnt es
sich nicht f"ur Alice zu betr"ugen.

Allerdings gibt es noch eine Sicherheitsl"ucke, da Alice oder der
H"andler den Scheck kopieren und mehrfach verwenden k"onnten. Mit dem
folgenden Protokoll soll diese Sicherheitsl"ucke beseitigt werden.


\subsubsection*{2. Protokoll} 

\begin{enumerate}
  \item Alice stellt 100 anonyme Bankschecks "uber jeweils 1000 DM
        aus. Auf jeden Scheck tr\"agt sie eine Zufallszahl ein, die lang
        genug ist, um die Wahrscheinlichkeit, da\3 eine andere Person
        dieselbe Zufallszahl benutzt, verschwindened klein ist.
  \item Alice steckt jeden Bankscheck zusammen mit Durchschlagspapier
        in einen Umschlag und tr"agt alle Umschl"age zur Bank.
  \item Die Bank "offnet 99 Umsch"age und pr"uft, ob der richtige
        Betrag (1000 DM) und eine jeweils andere Zufallszahl auf jedem
        Scheck von Alice eingetragen wurde.
  \item Die Bank unterzeichnet den letzten unge"offneten Umschlag. Die
        Unterschrift der Bank wird durch das Durchschlagspapier auf
        den Scheck "ubertragen.  Die Bank gibt den unge"offneten
        Umschag an Alice zur"uck und belastet ihr Konto mit 1000 DM.
  \item Alice "offnet den Umschag und verwendet den Scheck bei einem
        H"andler.
  \item Der H"andler "uberpr"uft die Unterschrift der Bank und
        akzeptiert den Scheck.
  \item Der H"andler bringt den Scheck zur Bank.
  \item Die Bank "uberpr"uft ihre Unterschrift und schaut in einer
        Datenbank nach, ob ein Scheck mit derselben Zufallszahl
        eingel"ost wurde. Falls nicht, dann zahlt die Bank dem
        H"andler 1000 DM aus. Die Bank speichert die Zufallszahl in
        ihrer Datenbank ab.
  \item Falls aber ein Scheck mit derselben Zufallszahl eingel"ost
        wurde, dann akzeptiert die Bank den Scheck nicht.
\end{enumerate}
Die Bank kann jetzt zwar einen Betrug feststellen, kann aber den 
Betr"uger nicht identifizieren. Der H"andler ist im Nachteil, da er nicht 
feststellen kann, ob Alice ihm eine Kopie des Schecks aush"andigt. 

\subsubsection*{3. Protokoll}

\begin{enumerate}
  \item Alice stellt 100 anonyme Bankschecks "uber jeweils 1000 DM
        aus. Auf jeden Scheck tr\"agt sie eine Zufallszahl ein, die lang
        genug ist, um die Wahrscheinlichkeit, da\3 eine andere Person
        dieselbe Zufallszahl benutzt, verschwindened klein ist.
  \item Alice steckt jeden Bankscheck zusammen mit Durchschlagspapier
        in einen Umschlag und tr"agt alle Umschl"age zur Bank.
  \item Die Bank "offnet 99 Umsch"age und pr"uft, ob der richtige
        Betrag (1000 DM) und eine jeweils andere Zufallszahl auf jedem
        Scheck von Alice eingetragen wurde.
  \item Die Bank unterzeichnet den letzten unge"offneten Umschlag. Die
        Unterschrift der Bank wird durch das Durchschlagspapier auf
        den Scheck "ubertragen.  Die Bank gibt den unge"offneten
        Umschag an Alice zur"uck und belastet ihr Konto mit 1000 DM.
  \item Alice "offnet den Umschag und verwendet den Scheck bei einem
        H"andler.
  \item Der H"andler "uberpr"uft die Unterschrift der Bank.
  \item Der H"andler fordert Alice auf einen Zufallsstring auf den
        Scheck zu schreiben.
  \item Alice tut dies.
  \item Der H"andler bringt den Scheck zur Bank.
  \item Die Bank "uberpr"uft ihre Unterschrift und schaut in einer
        Datenbank nach, ob ein Scheck mit derselben Zufallszahl
        eingel"ost wurde. Falls nicht, dann zahlt die Bank dem
        H"andler 1000 DM aus. Die Bank speichert die Zufallszahl und den
        Zufallsstring in ihrer Datenbank ab.
  \item Falls aber ein Scheck mit derselben Zufallszahl eingel"ost
        wurde, dann akzeptiert die Bank den Scheck nicht. Sie
        "uberpr"uft den Zufallsstring auf dem Scheck mit dem in der
        Datenbank. Falls sie "ubereinstimmen, wei\3 die Bank, da\3 der
        H"andler den Scheck kopiert hat. Falls sie verschieden sind,
        dann wei\3 die Bank, da\3 die Person, die den Scheck gekauft
        hat, ihn kopiert hat.
\end{enumerate}

Damit ist ein Betrug des H\"andlers, der sich der Bank zu erkennen
gibt, nicht mehr m\"oglich, die anonyme Alice ist aber in der Lage zu
betr\"ugen.
Um Alices Anonymit\"at zu sichern, solange sie sich korrekt verh\"alt,
ihre Identit\"at aber offenzulegen, sobald sie versucht, zu betr\"ugen
bedarf es einiger Hilfsmittel, die jetzt bereitgestellt werden. Weiterhin
ist noch eine elektronische Version der "`Unterschrift mit
Durchschlagspapier"' zu finden.

\subsection{Hilfsprotokolle}

\subsubsection*{Blinde Unterschrift}

Die Bank soll eine Nachricht $m$ von Alice ungesehen
unterschreiben. F"ur die Signatur wird das RSA-Signaturverfahren
benutzt. Die Bank verwendet ihren geheimen Schl"usse~$d$. Der
"offentliche Schl"ussel $(e,n)$ ist allgeimein zug"anglich.

Alice berechnet 
$$ t := m \cdot k^e \pmod{n},$$
wobei $k$ eine Zufallszahl ist.
Die Bank signiert $t$, indem sie $t^d$ berechnet und das Ergebnis
Alice liefert. Alice berechnet
$$ s := t^d \cdot k^{-1} \equiv m^d \cdot (k^e)^d \cdot k^{-1}
   \equiv m^{d} \cdot k^1 \cdot k^{-1} \equiv m^d \pmod{n} $$
und erh"alt die mit dem geheimen Schl"ussel der Bank signierter
Nachricht $m^d$. Die Bank kennt die Nachricht $m$ nicht.

Die blinde Unterschrift entspricht dem Unterschreiben des Umschlags
durch die Bank, wobei die Unterschrift durch das Durchschlagspapier
auf den Scheck "ubertragen wird.


\subsubsection*{Verschleiern einer Nachricht}

Alice w"ahlt eine Zufallszahl $k$ und ver"offentlicht $(f(I,k),k)$,
wobei $f$ eine bekannte Einwegfunktion ist. Jeder kann nach der
Bekanntgabe von $I$ verifizieren, da\3 auch wirklich $f(I,k)$
ver"offentlicht wurde.

Mit dieser Technik l"a"st sich das Problem, die Anonymit"at einer
ehrlichen Alice zu sichern und bei einem Betrug doch die Identit"at
des Betr"ugers zu erhalten:

Alice spaltet eine Bitfolge $I$, die ihre Identit"at vollst"andig
beschreibt (Name, Adresse und alle anderen Informationen, die die Bank
verlangt), in zwei Bitfolgen derart auf, da\3 $I = I_L \oplus I_R$ gilt.
Dann ver\"offentlicht Alice einige dieser "`halbierten Identit\"aten"' in der Form
$\left(f\left(I_{i_L},k_{i_1}\right),k_{i_1}\right)$ und
$\left(f\left(I_{i_R},k_{i_2}\right),k_{i_2}\right)$.
Kann man es erreichen, da\3 bei korrektem Verhalten von jedem Paar nur
eine H\"alfte offengelegt wird, bei einem Betrug jedoch von mindestens
einem Paar beide Teile, so ist diese scheinbar widerspr\"uchliche
Anforderung von Anonymit\"at und Bekanntgabe der Identit\"at
erf\"ullt.


\subsection{Ein vollst\"andiges Protokoll}

\begin{enumerate}
  \item Alice stellt $n$ anonyme Bankschecks "uber einen festen,
        jeweils identischen Betrag aus. Auf jeden Scheck tr\"agt sie
        eine gen\"ugend lange Zufallszahl ein.
        Auf jedem Scheck wird $m$ mal der Identit\"atsstring von
        Alice eingetragen, jeweils auf verschiedene Weisen halbiert 
        und verschleiert (s.o.).
  \item Alle $n$ Schecks werden von Alice f\"ur eine blinde
        Unterschrift der Bank vorbereitet.
  \item Alice "offnet nach Wahl der Bank $n-1$ dieser Schecks;
        die Bank pr\"uft den Betrag, die Zufallszahl und alle (von Alice
        auch offengelegten) Identit\"atsstrings.
  \item Die Bank signiert den letzten Scheck, gibt ihn Alice und belastet
        ihr Konto.
  \item Alice entfernt die Verschleierung und verwendet den Scheck bei einem
        H"andler.
  \item Der H"andler "uberpr"uft die Unterschrift der Bank.
  \item Der H"andler fordert Alice auf von den $m$ Identit\"atsstrings
        jeweils, nach seiner Wahl, die linke oder rechte H\"alfte bekanntzugeben.
  \item Alice tut dies.
  \item Der H"andler bringt den Scheck zur Bank.
  \item Die Bank "uberpr"uft ihre Unterschrift und schaut in einer
        Datenbank nach, ob ein Scheck mit derselben Zufallszahl
        eingel"ost wurde. Falls nicht, dann zahlt die Bank dem
        H"andler den Betrag aus. Die Bank speichert die Zufallszahl in
        ihrer Datenbank ab.
  \item Falls aber ein Scheck mit derselben Zufallszahl eingel"ost
        wurde, dann akzeptiert die Bank den Scheck nicht. Sie
        vergleicht die offengelegten H"alften der Identit"atsstrings 
        mit denen des Schecks in der Datenbank. Falls sie alle "ubereinstimmen, 
        wei\3 die Bank, da\3 der
        H"andler den Scheck kopiert hat. Falls sie verschieden sind,
        dann wei\3 die Bank, da\3 die Person, die den Scheck gekauft
        hat, ihn kopiert hat.
        Da der zweite H\"andler h\"ochst wahrscheinlich eine andere 
        Kombination der halbierten Identit\"aten hat offenlegen lassen,
        kennt die Bank jetzt von einem der Identit\"atsstrings beide 
        H\"alften und damit die Identit\"at des Betr\"ugers.
\end{enumerate}


  

\begin{thebibliography}{99}
\bibitem{copper}D. Coppersmith, A.\,M. Odlyzko, R. Schroeppel;
{\em Discrete Logarithms in $GF(p)$};
Algorithmica, 1986; S. 1--15.

\bibitem{denning}D. Denning;
{\em Cryptography and Data Security};
Addison  Wesley, London, 1983.

\bibitem{odlyzko}A.\,M. Odlyzko;
{\em Discrete logarithms in finite fields and their cryptographic 
     significance};
S. 224--314.

\bibitem{frank}
F. Schaefer-Lorinser;
{\em Algorithmische Zahlentheorie};
Vorlesungsskript, Universit"at Karlsruhe.

\bibitem{schneier}
B. Schneier;
{\em Applied Cryptography};
2. Ausgabe, Wiley, 1996.

\bibitem{seberry}
J. Seberry, J. Pieprzyk;
{\em Cryptography};
Prentice Hall, New York, 1989.

\bibitem{simmons}
G.\,J. Simmons (Ed.);
{\em Contemporary Cryptology};
IEEE Press, 1992.

\bibitem{stinson}
D.\,R. STinson;
{\em Cryptography, Theory and Practice};
CRC Press, Boca Taton, Florida, 1995.

\end{thebibliography}

\end{document}




